\documentclass[11pt,a4paper]{article}
\usepackage[T1]{fontenc}
\usepackage[utf8]{inputenc}
\usepackage{lmodern,amsmath,amssymb,amsthm,mathtools,microtype,mathrsfs}
\usepackage[margin=25mm,top=23mm,bottom=25mm]{geometry}
\usepackage{booktabs,array,longtable,enumitem,xcolor,fancyhdr}
\usepackage[colorlinks=true,linkcolor=blue!45!black,urlcolor=blue!45!black,citecolor=blue!45!black]{hyperref}
\hypersetup{pdftitle={Reading the Alpoge--Claude RH and Dirichlet-GRH proof: working version 0.47},pdfauthor={ChatGPT},pdfsubject={Partial human-readable reconstruction with explicit arithmetic inputs},pdfkeywords={Riemann hypothesis, Lean, theta lifts, prime sums, Arakelov slopes}}
\setlength{\parindent}{0pt}
\setlength{\parskip}{5pt plus 1pt}
\setlist{topsep=3pt,itemsep=3pt,leftmargin=*}
\setlength{\headheight}{14pt}
\pagestyle{fancy}\fancyhf{}
\fancyhead[L]{\small RH and Dirichlet GRH: a reconstruction}
\fancyhead[R]{\small Working version 0.47}
\fancyfoot[C]{\thepage}
\renewcommand{\headrulewidth}{0.3pt}
\numberwithin{equation}{section}
\newtheorem{theorem}{Theorem}[section]
\newtheorem{proposition}[theorem]{Proposition}
\newtheorem{lemma}[theorem]{Lemma}
\newtheorem{corollary}[theorem]{Corollary}
\theoremstyle{remark}\newtheorem{remark}[theorem]{Remark}
\newcommand{\Rea}{\operatorname{Re}}
\newcommand{\HS}{\mathrm{HS}}
\newcommand{\tr}{\operatorname{tr}}
\newcommand{\codim}{\operatorname{codim}}
\newcommand{\covol}{\operatorname{covol}}
\newcommand{\OK}{\mathcal O_K}
\newcommand{\src}[2]{\href{https://github.com/CocoSgt/riemann-hypothesis-lean/blob/7cbf7a0c32fca1badf05931879a37c3f8a2b03b4/src/#1}{\texttt{#2}}}
\newcommand{\status}[1]{\textbf{Status:} #1}
\begin{document}
\begin{center}
{\LARGE\bfseries RH and Dirichlet GRH: new reconstruction chapters}\\[7pt]
{\large Version 0.47 --- 23 November 2026}\\[5pt]
{\large Rigidity, canonical connections, and original setting inputs}
\end{center}
\textbf{Scope.} This edition contains the current audit and the new
chapters. It remains a partial human reconstruction. Earlier named
claims are in the cumulative manuscript and preserved component files.
The underlying formal proof is accepted as correct; the exact remaining
human reconstruction boundaries are stated below.
\setcounter{tocdepth}{1}
\tableofcontents
\clearpage
\begingroup\emergencystretch=3em
\section{Version 0.47: rigidity, canonical connections, and original setting inputs}
\label{sec:dependency-audit-v47}

\subsection{Current result and the two completed selected contracts}

The pinned formal source is accepted as fully checked and correct.
The task is to reconstruct its mathematical argument in readable prose,
including any simpler argument that proves the same original statement.
Version 0.47 adds 83 numbered claims to the 2768 preserved predecessor
claims, for 2851 in all. The dependency graph records the human
reconstruction, rather than a mechanically extracted formal dependency
graph.

The full \texttt{RIG\_\allowbreak II} contract is now reconstructed. Its original
local point produces an explicit rank-two extension. The literal
extension-modularity antecedent, the matched Kato theorem, quadratic
base change and the polarized adjoint theorem give primitives in every
summand of the full nine-dimensional adjoint. No replacement local
point or independent automorphy premise is silently substituted.

The distinct record-construction task
\texttt{TASK\_\allowbreak INP\_\allowbreak Z\_\allowbreak ASSEMBLY} is also complete. The precise
\(\mathrm{CohInp}\) fields, the uniform \(\mathrm{LemmaK}\)
witness, and all six additional \(\mathrm{InpZ}\) fields are supplied
by named human proofs on the original objects. In particular,
\(\mathrm{RigII}\) and the previously completed rank-three
congruence integer occupy their actual original fields.
The one-character trace/length theorem and the common
unramified-coefficient automorphy package are separate producers;
neither is inferred merely from constructing this input record.

The selected sufficient route retains 121 nodes and 161 edges.
Its whole-obligation register decreases from 28 to 26:
19 class C and 7 class B. There are 19 shared RH/Dirichlet-GRH
obligations, 7 additional Dirichlet-GRH obligations, and no
RH-only obligations in this route. A claim count measures written
inventory. One whole obligation may still contain many substantial
producers, so these counts do not measure a percentage of completion.

\subsection{What the seven new chapters establish}

\begin{enumerate}
\item KTR proves the complete original \(\mathrm{KatoUse}\),
\(\mathrm{KatoTwo}\) and \(\mathrm{RigII}\). For weight
\(2e+1\), the actual Galois representation has Kato twist \(e\),
which lies in the required noncentral critical range. The local
ordered-product ideal gives an inertia flag on the actual corner;
inertia and the class-number hypothesis kill its middle and determinant
characters. The corner extension is irreducible on the full infinite
cyclotomic kernel. Its quadratic base change has the exact symmetric
polarization required for the adjoint theorem. Explicit projections
decompose all nine coordinates into the corner adjoint, two rank-two
pieces and a quadratic scalar line, and explicit sections recombine
their coboundary primitives.

\item KIZ assembles the complete original cohomological input record.
The local comparison, normalized Poitou--Tate, global finiteness and
Euler-characteristic fields use the same cohomology conventions.
The uniform \(\mathrm{LemmaK}\) constant is selected before every
varying later arithmetic datum. The exact \(\mathrm{LemmaM}\),
\(\mathrm{LemmaM'}\), \(\mathrm{LemmaP}\), \(\mathrm{RigII}\),
\(\mathrm{AdVanishK}\) and congruence-integer fields are then inserted.

\item RFD proves the entire literal \(\mathrm{RowsFinDim}\), with
arbitrary original locally constant supported test function and
arbitrary original unit vectors. A common compact-open level for the
test function reduces the actual adelic quotient to finitely many
arithmetic real components. Harish--Chandra finite multiplicity and
the exact Schur-projection compression identity give a finite-dimensional
range. No extra idempotence, smoothness, compact-finiteness, gap or
cuspidality condition is imposed on the given data.

\item ASG constructs eight dependent original adelic blocks:
the unitary instance, integral lattice and flags, finite character,
reference frame, transport matrices, automorphic bundle, finite
representation block, and subgroup/quotient measures.
The original \(\mathrm{Forms0}\) antecedent gives an extra
isotropic vector. An ideal-class construction supplies primitive
vectors of every required negative norm; a trace-one triangular
correction gives the integral hyperbolic frame without dividing by two
in the integer ring. Split local linear algebra supplies the complete
generic reference-frame and transport predicates.
The finite highest weights are exactly the original
\(\mathrm{lowK}(\mathrm{allPos},u)=u-\rho_C\).
The bundle retains the original gap hypothesis. The Bessel vector,
measured-field, Siegel--Weil ties, arithmetic count and common-family
producers remain explicit.

\item CNC constructs the full generic canonical gauge and the literal
original \(\mathrm{IsNab}\) connection over arbitrary original
coefficient algebras, including nonreduced and nonlocal rings.
Global formal jets, normalized line generators and proper-base-function
uniqueness give the exact biextension laws. A formal logarithm,
two symmetric-cocycle primitives, a separability-idempotent correction
and an antisymmetric half correction construct the canonical gauge.
The unique germ of every section gives the original derivative formula.
The separate connection input in ICL is thus supplied.

\item PCL reconstructs the actual Hermitian coordinates and the full
trace-dual lattice calculation. The exact different and fractional-ideal
identity make \(L'=\mathfrak A^{-1}\Lambda\) self-dual, without
principalizing any ideal. On each supplied original principal chart,
the actual half-different and quadratic ideal give the prescribed
\(N_A,N_C,k_0\), their lattice inclusions and both logarithmic bounds.
The full \(\mathrm{DataX}\) and all-row model endpoints follow with
\(\mathrm{IsUnf}\), \(\mathrm{HasPrinc}/\mathrm{GamAt}\) and
\(\mathrm{LawsCore}\) retained. The common uniformization is not
constructed by this conditional lattice argument.

\item LGS supplies the actual local-growth, canonical mean-deficit
and full auxiliary admission inputs retained by GLG. It proves the
small-order and band estimates, lattice projection and finite-walk
bounds, first growth and normalized lengths. Setting only the
auxiliary chain to zero supplies every field of the admitted record;
the original valuation lattices, comparisons, supports and family
stay fixed. All five \(\mathrm{LinkG}\) conclusions follow with
\(C_G=4C_{\mathrm{count}}+1\) and
\(b_o=b^{x^{-\epsilon/2}}\).
The actual common \(\mathrm{GroupCompatibilityFamily}\) and
\(\mathrm{Adm2F}\) remain inputs; \(\mathrm{TopsF}\) supplies
the latter. At class level the two original \(\mathrm{Closed}\)
predicates, including their earlier cutoffs and coefficient choices,
remain explicit. No whole global-gain closure is claimed.
\end{enumerate}

\subsection{Foundations, scope, preservation and review}

There are four new class-A applications: the exact noncentral critical
Kato theorem, quadratic \(\mathrm{GL}_2\) base change,
Harish--Chandra finite multiplicity, and the Hilbert class
field/Chebotarev ideal-selection argument. Each is matched to its
primary statement and the actual data in the corresponding chapter.
The other new claims are class-B human proofs at their printed scopes.
Previously matched cohomological, period, automorphic and geometric
foundations remain identified separately from the numbered dependency
graph.

Every new claim is independently reviewed by an agent other than its
author. The review manifests bind the exact whole proof file and each
reviewed block. Exact cited source bytes are compared with the pinned
git blobs and included in the archive. All predecessor claim records
and the entire historical proof body are preserved. Mutable predecessor
metadata is frozen in \texttt{history/\allowbreak version\_\allowbreak 0\_\allowbreak 46}.

Draft-mode TeX compilation checks the two complete readers for syntax,
cross-references and missing characters. It does not constitute PDF
visual review or formal replay. No Lean build, new formal-correctness
audit or completed human proof of RH/GRH is claimed.

The following register lists every remaining whole selected obligation.
Its version-0.47 qualifications are the current scope authority;
historical conditional statements retain their printed hypotheses even
when a new chapter now supplies one of those hypotheses.

\subsection{The complete current selected register}

\paragraph{1. GLOBAL\_\allowbreak DENOMINATOR — C — Actual LinkN local table and binding producer}

Closed(Des $\mathcal C$ 3) LinkN: the actual etaSum is bounded with exact dZ0 and same-host vvTerm, subject only to the literal ShaFinOdd and summability guards.

\emph{Shared objects and choice order.} one FamP; one class $\mathcal A$ and its target-compatible subclass; one coefficient field per family datum; one x and i at each application; original selected rows, words and levels; thresholds chosen before the common family; actual D.nd; actual etaJ and Aj; sha\_\allowbreak eq

\emph{Source.} \href{https://github.com/CocoSgt/riemann-hypothesis-lean/blob/7cbf7a0c32fca1badf05931879a37c3f8a2b03b4/src/Formalism/SharpRestriction/GlobalTheoremInputs.lean#L45-L48}{LinkN}.

\emph{Source.} \href{https://github.com/CocoSgt/riemann-hypothesis-lean/blob/7cbf7a0c32fca1badf05931879a37c3f8a2b03b4/src/Formalism/SharpRestriction/GlobalTheoremInputs.lean#L40-L44}{LinkN\_\allowbreak at}.

\emph{Source.} \href{https://github.com/CocoSgt/riemann-hypothesis-lean/blob/7cbf7a0c32fca1badf05931879a37c3f8a2b03b4/src/Development/AutomorphicArithmetic/CharacterBoxes/DivinNOfHypV.lean#L117-L134}{divinN\_\allowbreak of\_\allowbreak hypV}.



\paragraph{2. GLOBAL\_\allowbreak RETAINED\_\allowbreak DEFECT — C — Original retained-row defect bound LinkIRR}

Rows.DelpNonneg and delp$\le$exp(-x\textasciicircum{}(ez/7)) uniformly for every late actual row datum. CURRENT VERSION0.47 QUALIFICATION: Finish the same-D/common-family adelic setting after ASG: normalized nonzero Bessel W and decay, measured-field block, original Siegel-Weil rescaling and measure/character ties, Bessel uniqueness/local ties, then the original AutTie/RowDictTie/LabTie. RowsFinDim is now proved universally by RFD. Supply the actual CountIdempotentsClass theorem or its three sufficient arithmetic cylinder counts, with one simultaneous family threshold; reuse the completed ARO structure. Completed new original producers: ARO1-ARO30, RFD1-RFD9, ASG1-ASG18. Any earlier narrative naming these completed producers as outstanding is historical; this qualification is the current scope authority.

\emph{Shared objects and choice order.} one FamP; one class $\mathcal A$ and its target-compatible subclass; one coefficient field per family datum; one x and i at each application; original selected rows, words and levels; thresholds chosen before the common family; AllRows A; LabelCountGoodInjective; AllJoinA; CountAllRows

\emph{New completed inputs.} ARO1-ARO30, RFD1-RFD9, ASG1-ASG18

\emph{Current remaining subfrontier.} Finish the same-D/common-family adelic setting after ASG: normalized nonzero Bessel W and decay, measured-field block, original Siegel-Weil rescaling and measure/character ties, Bessel uniqueness/local ties, then the original AutTie/RowDictTie/LabTie. RowsFinDim is now proved universally by RFD. Supply the actual CountIdempotentsClass theorem or its three sufficient arithmetic cylinder counts, with one simultaneous family threshold; reuse the completed ARO structure.

\emph{Source.} \href{https://github.com/CocoSgt/riemann-hypothesis-lean/blob/7cbf7a0c32fca1badf05931879a37c3f8a2b03b4/src/Formalism/SharpRestriction/GlobalTheoremInputs.lean#L55-L58}{LinkIRR}.

\emph{Source.} \href{https://github.com/CocoSgt/riemann-hypothesis-lean/blob/7cbf7a0c32fca1badf05931879a37c3f8a2b03b4/src/Development/AbelianVarietyTheory/Proofs/GlobalInterpolation/LinkIRROfJoinEx.lean#L17-L38}{linkIRR\_\allowbreak of\_\allowbreak join\_\allowbreak ex}.

\emph{Source.} \href{https://github.com/CocoSgt/riemann-hypothesis-lean/blob/7cbf7a0c32fca1badf05931879a37c3f8a2b03b4/src/Development/AbelianVarietyTheory/Proofs/GlobalInterpolation/LinkIRROfJoinEx.lean#L143-L166}{link\_\allowbreak irr}.



\paragraph{3. GLOBAL\_\allowbreak PERIOD\_\allowbreak HEIGHT — C — Full actual period and height comparison LinkDIVt}

Closed conditional TuSize -> DGamSize -> (LinkEta -> DIVinN -> LinkDIV), with every original LinkDIV\_\allowbreak at clause.

\emph{Shared objects and choice order.} one FamP; one class $\mathcal A$ and its target-compatible subclass; one coefficient field per family datum; one x and i at each application; original selected rows, words and levels; thresholds chosen before the common family; same RR target; same Bu/besS; same Nu and Bk

\emph{Source.} \href{https://github.com/CocoSgt/riemann-hypothesis-lean/blob/7cbf7a0c32fca1badf05931879a37c3f8a2b03b4/src/Formalism/SharpRestriction/GlobalTheoremInputs.lean#L77-L90}{LinkDIV\_\allowbreak at}.

\emph{Source.} \href{https://github.com/CocoSgt/riemann-hypothesis-lean/blob/7cbf7a0c32fca1badf05931879a37c3f8a2b03b4/src/Formalism/SharpRestriction/GlobalTheoremInputs.lean#L94-L96}{LinkDIVt}.

\emph{Source.} \href{https://github.com/CocoSgt/riemann-hypothesis-lean/blob/7cbf7a0c32fca1badf05931879a37c3f8a2b03b4/src/Development/AbelianVarietyTheory/Proofs/GlobalInterpolation/LinkBrest.lean#L7-L25}{GlobalTheory.Top.link\_\allowbreak brest}.

\emph{Source.} \href{https://github.com/CocoSgt/riemann-hypothesis-lean/blob/7cbf7a0c32fca1badf05931879a37c3f8a2b03b4/src/Development/AbelianVarietyTheory/Proofs/GlobalInterpolation/WireSel.lean#L11-L41}{GlobalTheory.Top.wire\_\allowbreak sel}.



\paragraph{4. GLOBAL\_\allowbreak LOCAL\_\allowbreak GAIN — C — Original global LinkG installation}

One uniform LinkG with actual local lattice lengths, served-orbit coverage and MeanKappaBound; delta(x)->0 before varying actual indices. CURRENT VERSION0.47 QUALIFICATION: Construct the original common GroupCompatibilityFamily and selected tops/Adm2F for global local gain. LGS now supplies full auxiliary admission, actual first-growth and mean inputs; these should not be reopened. Continue the distinct denominator, period/height, normalized spectral and second-layer compatible-selection producers with the original common objects and thresholds. Completed new original producers: GLG1-GLG18, LGS1-LGS13. Any earlier narrative naming these completed producers as outstanding is historical; this qualification is the current scope authority.

\emph{Shared objects and choice order.} one FamP; one class $\mathcal A$ and its target-compatible subclass; one coefficient field per family datum; one x and i at each application; original selected rows, words and levels; thresholds chosen before the common family; actual G.Inst; actual GK; actual lenMN; same orbit representatives

\emph{New completed inputs.} GLG1-GLG18, LGS1-LGS13

\emph{Current remaining subfrontier.} Construct the original common GroupCompatibilityFamily and selected tops/Adm2F for global local gain. LGS now supplies full auxiliary admission, actual first-growth and mean inputs; these should not be reopened. Continue the distinct denominator, period/height, normalized spectral and second-layer compatible-selection producers with the original common objects and thresholds.

\emph{Source.} \href{https://github.com/CocoSgt/riemann-hypothesis-lean/blob/7cbf7a0c32fca1badf05931879a37c3f8a2b03b4/src/Formalism/SharpRestriction/GlobalTheoremInputs.lean#L122-L137}{LinkG\_\allowbreak at}.

\emph{Source.} \href{https://github.com/CocoSgt/riemann-hypothesis-lean/blob/7cbf7a0c32fca1badf05931879a37c3f8a2b03b4/src/Formalism/SharpRestriction/GlobalTheoremInputs.lean#L138-L141}{LinkG}.

\emph{Source.} \href{https://github.com/CocoSgt/riemann-hypothesis-lean/blob/7cbf7a0c32fca1badf05931879a37c3f8a2b03b4/src/Development/AbelianVarietyTheory/Proofs/GlobalInterpolation/LinkG.lean#L14-L27}{GlobalTheory.Top.link\_\allowbreak g}.



\paragraph{5. GLOBAL\_\allowbreak SPECTRAL — C — Original normalized kernel/Gram estimates LinkS}

Closed TuSize -> LinkS with exact Gamma normalization, off-diagonal Hilbert-Schmidt bound, diagonal error and Dall/Dstar estimate, errors exp(-x\textasciicircum{}(epsilon/8)).

\emph{Shared objects and choice order.} one FamP; one class $\mathcal A$ and its target-compatible subclass; one coefficient field per family datum; one x and i at each application; original selected rows, words and levels; thresholds chosen before the common family; actual KD; same S.Cu spaces; GamOf; Dstar; Dall

\emph{Source.} \href{https://github.com/CocoSgt/riemann-hypothesis-lean/blob/7cbf7a0c32fca1badf05931879a37c3f8a2b03b4/src/Formalism/SharpRestriction/GlobalTheoremInputs.lean#L145-L149}{LinkS\_\allowbreak at}.

\emph{Source.} \href{https://github.com/CocoSgt/riemann-hypothesis-lean/blob/7cbf7a0c32fca1badf05931879a37c3f8a2b03b4/src/Formalism/SharpRestriction/GlobalTheoremInputs.lean#L152-L157}{LinkS}.

\emph{Source.} \href{https://github.com/CocoSgt/riemann-hypothesis-lean/blob/7cbf7a0c32fca1badf05931879a37c3f8a2b03b4/src/Development/AbelianVarietyTheory/Proofs/GlobalInterpolation/LinkAlt.lean#L16-L50}{GlobalTheory.Top.link\_\allowbreak alt}.



\paragraph{6. TASK\_\allowbreak ONE\_\allowbreak CHAR\_\allowbreak PROOF — C — Reconstruct the original one-character argument}

From exact InpZ and AutomorphyUnramifiedCoefficients, prove the original OneChar trace/length estimate on the supplied coefficient datum and GD input. CURRENT VERSION0.47 QUALIFICATION: The actual original OneChar trace/length estimate and the common marked unramified-coefficient automorphy datum. The complete InpZ record and finite PropGK consumer are now available. Completed new original producers: KTR1-KTR13, KIZ input record, RNC16, PGK1-PGK6. Any earlier narrative naming these completed producers as outstanding is historical; this qualification is the current scope authority.

\emph{New completed inputs.} KTR1-KTR13, KIZ input record, RNC16, PGK1-PGK6

\emph{Current remaining subfrontier.} The actual original OneChar trace/length estimate and the common marked unramified-coefficient automorphy datum. The complete InpZ record and finite PropGK consumer are now available.

\emph{Source.} \href{https://github.com/CocoSgt/riemann-hypothesis-lean/blob/7cbf7a0c32fca1badf05931879a37c3f8a2b03b4/src/Development/AutomorphicArithmetic/AutomorphicVarieties/UnitaryThreeGroups/KatoUse.lean#L54-L64}{OneChar}.

\emph{Source.} \href{https://github.com/CocoSgt/riemann-hypothesis-lean/blob/7cbf7a0c32fca1badf05931879a37c3f8a2b03b4/src/Development/AutomorphicArithmetic/AutomorphicVarieties/ModularityArgument/KatoQuadraticCharacters/OneChar.lean#L158-L178}{UpperTheory.V.Z.Kato.oneChar}.



\paragraph{7. AUT\_\allowbreak UNRAM\_\allowbreak COEF — C — Actual unramified-coefficient automorphy}

For each actual Dat with coefficient ramification one and fixed iota, choose one marked g1 supplying FinA, AutB and LemW together.

\emph{Source.} \href{https://github.com/CocoSgt/riemann-hypothesis-lean/blob/7cbf7a0c32fca1badf05931879a37c3f8a2b03b4/src/Development/AutomorphicArithmetic/AutomorphicVarieties/AutomorphicCases/AutomorphyUnramifiedCoefficients.lean#L24-L27}{AutomorphyUnramifiedCoefficients}.

\emph{Source.} \href{https://github.com/CocoSgt/riemann-hypothesis-lean/blob/7cbf7a0c32fca1badf05931879a37c3f8a2b03b4/src/Development/AutomorphicArithmetic/AutomorphicVarieties/AutomorphicCases/AutomorphyUnramifiedCoefficients.lean#L12-L13}{FinA}.

\emph{Source.} \href{https://github.com/CocoSgt/riemann-hypothesis-lean/blob/7cbf7a0c32fca1badf05931879a37c3f8a2b03b4/src/Development/AutomorphicArithmetic/AutomorphicVarieties/AutomorphicCases/AutomorphyUnramifiedCoefficients.lean#L15-L17}{AutB}.

\emph{Source.} \href{https://github.com/CocoSgt/riemann-hypothesis-lean/blob/7cbf7a0c32fca1badf05931879a37c3f8a2b03b4/src/Development/AutomorphicArithmetic/AutomorphicVarieties/AutomorphicCases/AutomorphyUnramifiedCoefficients.lean#L19-L22}{LemW}.



\paragraph{8. TASK\_\allowbreak CASE3\_\allowbreak AUT\_\allowbreak JOIN — B — Complete the original Case-III automorphy/image join}

Install the actual Case-III ProModAlt, SeedIVAlt and shared PMExAlt with the inherited CohomologyHyp/Ferrero-Washington/HFit/SeedII facts into the original whole Case-III automorphy/image conclusion.

\emph{Source.} \href{https://github.com/CocoSgt/riemann-hypothesis-lean/blob/7cbf7a0c32fca1badf05931879a37c3f8a2b03b4/src/Development/AutomorphicArithmetic/AutomorphicVarieties/ModularityArgument/ThresholdArguments/DenseThresholdCases/PromodularityCaseThree.lean#L43-L44}{AutomorphyRankFourDense}.

\emph{Source.} \href{https://github.com/CocoSgt/riemann-hypothesis-lean/blob/7cbf7a0c32fca1badf05931879a37c3f8a2b03b4/src/Development/AutomorphicArithmetic/AutomorphicVarieties/ModularityArgument/ThresholdArguments/DenseThresholdCases/CuFinal.lean#L44-L44}{main4\_\allowbreak alt}.



\paragraph{9. TASK\_\allowbreak CASE3\_\allowbreak PROMOD\_\allowbreak JOIN — C — Reconstruct the actual Case-III chain-to-promodularity proof}

Use the original tensor-factor tower, ChainRankFour, ChainDesc4, seeded rank-four chain reduction, actual PMExAlt and SeedIVAlt to construct ProModAlt on every original Ctx3 star point at the same fixed iota, preserving all finite-coefficient, inertia and support maps.

\emph{Source.} \href{https://github.com/CocoSgt/riemann-hypothesis-lean/blob/7cbf7a0c32fca1badf05931879a37c3f8a2b03b4/src/Development/AutomorphicArithmetic/AutomorphicVarieties/ModularityArgument/ThresholdArguments/DenseThresholdCases/PromodularityCaseThreeOfLemmas.lean#L131-L153}{UpperTheory.V.Cu.promodularity\_\allowbreak case\_\allowbreak three\_\allowbreak of\_\allowbreak chain\_\allowbreak reduction}.

\emph{Source.} \href{https://github.com/CocoSgt/riemann-hypothesis-lean/blob/7cbf7a0c32fca1badf05931879a37c3f8a2b03b4/src/Development/AutomorphicArithmetic/AutomorphicVarieties/ModularityArgument/ThresholdArguments/DenseThresholdCases/PromodularityCaseThree.lean#L40-L42}{PromodularityCaseThreeDense}.



\paragraph{10. SEED\_\allowbreak IV\_\allowbreak ALT — C — Actual Case-III seed with prechosen residual datum}

Base Ctx3/Four tower, iota and Amb -> w,s before every later Dies tower -> compatible finite-coefficient seed with exact residual support/root entries.

\emph{Source.} \href{https://github.com/CocoSgt/riemann-hypothesis-lean/blob/7cbf7a0c32fca1badf05931879a37c3f8a2b03b4/src/Development/AutomorphicArithmetic/AutomorphicVarieties/ModularityArgument/ThresholdArguments/DenseThresholdCases/PromodularityCaseThree.lean#L19-L25}{SeedIVAlt}.



\paragraph{11. TENSOR\_\allowbreak FACTOR\_\allowbreak TOWER — B — Actual tensor-factor tower and characters}

Original finite-field, tame, factor and character constraints; source-level construction is sizeable and has not been reconstructed here.

\emph{Source.} \href{https://github.com/CocoSgt/riemann-hypothesis-lean/blob/7cbf7a0c32fca1badf05931879a37c3f8a2b03b4/src/Development/AutomorphicArithmetic/AutomorphicVarieties/ModularityArgument/ThresholdArguments/DenseThresholdCases/PromodularityCaseThreeOfLemmas.lean#L131-L153}{UpperTheory.V.Cu.promodularity\_\allowbreak case\_\allowbreak three\_\allowbreak of\_\allowbreak chain\_\allowbreak reduction}.



\paragraph{12. CHAIN\_\allowbreak RANK\_\allowbreak FOUR — C — Actual SomeIn-to-AllIn propagation}

Same source Space and PM ideal, residual point Good4, original bound 30(C+1)\textasciicircum{}2 <= d; geometry must preserve all identities.

\emph{Source.} \href{https://github.com/CocoSgt/riemann-hypothesis-lean/blob/7cbf7a0c32fca1badf05931879a37c3f8a2b03b4/src/Development/AutomorphicArithmetic/AutomorphicVarieties/ModularityArgument/ThresholdArguments/DenseThresholdCases/PromodularityCaseThree.lean#L26-L32}{ChainRankFour}.



\paragraph{13. CHAIN\_\allowbreak DESC4 — B — Original residual chain descent}

Finite residual descent retaining source Chain4 compatibility.

\emph{Source.} \href{https://github.com/CocoSgt/riemann-hypothesis-lean/blob/7cbf7a0c32fca1badf05931879a37c3f8a2b03b4/src/Development/AutomorphicArithmetic/AutomorphicVarieties/ModularityArgument/ThresholdArguments/DenseThresholdCases/PromodularityCaseThreeOfLemmas.lean#L131-L153}{UpperTheory.V.Cu.promodularity\_\allowbreak case\_\allowbreak three\_\allowbreak of\_\allowbreak chain\_\allowbreak reduction}.



\paragraph{14. RANK4\_\allowbreak CHAIN\_\allowbreak REDUCTION — C — Seeded reduction to the original residual chain}

Select T1 before every later Deep tower; retain one compatible seed and Chain4 for every original residual lift over Q.

\emph{Source.} \href{https://github.com/CocoSgt/riemann-hypothesis-lean/blob/7cbf7a0c32fca1badf05931879a37c3f8a2b03b4/src/Development/AutomorphicArithmetic/AutomorphicVarieties/ModularityArgument/ThresholdArguments/DenseThresholdCases/PromodularityCaseThree.lean#L33-L39}{RankFourChainReductionDense}.

\emph{Source.} \href{https://github.com/CocoSgt/riemann-hypothesis-lean/blob/7cbf7a0c32fca1badf05931879a37c3f8a2b03b4/src/Development/AutomorphicArithmetic/AutomorphicVarieties/ModularityArgument/ThresholdArguments/DenseThresholdCases/CuPChainIV.lean#L16-L82}{UpperTheory.V.Cu.chainIV\_\allowbreak alt}.



\paragraph{15. TASK\_\allowbreak COUNT\_\allowbreak V\_\allowbreak INSTALL — B — Complete the original CountV finite assembly}

From exactly the listed source host/cohomology inputs, complete countV\_\allowbreak of\_\allowbreak finiteness\_\allowbreak duality\_\allowbreak euler\_\allowbreak host, including original host/frame formation, normalized length, the CaseV guard and R1<=2\textasciicircum{}c2.

\emph{Source.} \href{https://github.com/CocoSgt/riemann-hypothesis-lean/blob/7cbf7a0c32fca1badf05931879a37c3f8a2b03b4/src/Development/AutomorphicArithmetic/AutomorphicVarieties/HeightAndVanishing/Counts.lean#L69-L69}{CountV}.

\emph{Source.} \href{https://github.com/CocoSgt/riemann-hypothesis-lean/blob/7cbf7a0c32fca1badf05931879a37c3f8a2b03b4/src/Development/AutomorphicArithmetic/AutomorphicVarieties/ModularityArgument/ThresholdArguments/DenseThresholdCases/CountVOf.lean#L33-L60}{UpperTheory.V.UnitaryThreeInertia.countV\_\allowbreak of\_\allowbreak finiteness\_\allowbreak duality\_\allowbreak euler\_\allowbreak host}.



\paragraph{16. THM\_\allowbreak S6\_\allowbreak S — C — Original unitary-three ordinary support and automorphy}

Original host/frame and fixed iota -> one marked g1 -> every universal R,u -> weight-zero finiteness and every Type-I BPrime\_\allowbreak s point.

\emph{Source.} \href{https://github.com/CocoSgt/riemann-hypothesis-lean/blob/7cbf7a0c32fca1badf05931879a37c3f8a2b03b4/src/Development/AutomorphicArithmetic/AutomorphicVarieties/ModularityArgument/ThresholdArguments/DenseThresholdCases/ThmS6CoreS.lean#L17-L26}{ThmS6\_\allowbreak s}.

\emph{Source.} \href{https://github.com/CocoSgt/riemann-hypothesis-lean/blob/7cbf7a0c32fca1badf05931879a37c3f8a2b03b4/src/Development/AutomorphicArithmetic/AutomorphicVarieties/ModularityArgument/ThresholdArguments/DenseThresholdCases/ThmS6CoreS.lean#L12-L16}{BPrime\_\allowbreak s}.



\paragraph{17. HOST\_\allowbreak FRAME\_\allowbreak LENGTH — B — Original Case-V host/frame and normalized-length installation}

Actual HostData, frame, local-condition equality, abstract count and normalized original Selmer length must all refer to the same objects.

\emph{Source.} \href{https://github.com/CocoSgt/riemann-hypothesis-lean/blob/7cbf7a0c32fca1badf05931879a37c3f8a2b03b4/src/Development/AutomorphicArithmetic/AutomorphicVarieties/ModularityArgument/ThresholdArguments/DenseThresholdCases/CountVOf.lean#L33-L60}{UpperTheory.V.UnitaryThreeInertia.countV\_\allowbreak of\_\allowbreak finiteness\_\allowbreak duality\_\allowbreak euler\_\allowbreak host}.



\paragraph{18. HOST\_\allowbreak NUMBER\_\allowbreak S — C — Uniform original-host rank-three catalogue}

One exponent before all closures, coefficient rings, levels, characters, frames, embeddings and marked points; one separating trace function controls all later finite pairwise nonisomorphic collections.

\emph{Source.} \href{https://github.com/CocoSgt/riemann-hypothesis-lean/blob/7cbf7a0c32fca1badf05931879a37c3f8a2b03b4/src/Development/AutomorphicArithmetic/AutomorphicVarieties/ModularityArgument/ThresholdArguments/DenseThresholdCases/HostNumberS.lean#L18-L28}{HostNumber\_\allowbreak s}.



\paragraph{19. TASK\_\allowbreak HOST\_\allowbreak VAN\_\allowbreak E\_\allowbreak INSTALL — B — Transport K/duality to literal HostVanE}

Use the already supplied LemmaK/local-global duality and finiteness on the actual CountV host/frame and exact original coefficient modules.

\emph{Source.} \href{https://github.com/CocoSgt/riemann-hypothesis-lean/blob/7cbf7a0c32fca1badf05931879a37c3f8a2b03b4/src/Development/AutomorphicArithmetic/AutomorphicVarieties/UnitaryThreeDeformations/HostNumber.lean#L9-L12}{HostVanE}.

\emph{Source.} \href{https://github.com/CocoSgt/riemann-hypothesis-lean/blob/7cbf7a0c32fca1badf05931879a37c3f8a2b03b4/src/Development/AutomorphicArithmetic/AutomorphicVarieties/LevelCohomology/SoHostVanE.lean#L11-L18}{UpperTheory.V.X.So.hostVanE}.



\paragraph{20. TASK\_\allowbreak HOST\_\allowbreak XPP\_\allowbreak INSTALL — B — Transport CorX to literal HostXpp}

Preserve Psi\textasciicircum{}-2 eta, the prime-to-ell quotient, actual local conditions and the original CountV host while installing BR7.

\emph{Source.} \href{https://github.com/CocoSgt/riemann-hypothesis-lean/blob/7cbf7a0c32fca1badf05931879a37c3f8a2b03b4/src/Development/AutomorphicArithmetic/AutomorphicVarieties/UnitaryThreeDeformations/HostNumber.lean#L14-L17}{HostXpp}.

\emph{Source.} \href{https://github.com/CocoSgt/riemann-hypothesis-lean/blob/7cbf7a0c32fca1badf05931879a37c3f8a2b03b4/src/Development/AutomorphicArithmetic/AutomorphicVarieties/LevelCohomology/Xpp.lean#L11-L84}{UpperTheory.V.X.xpp}.



\paragraph{21. DESIGN\_\allowbreak CHOICES\_\allowbreak COMMON — C — Construct the original common design family and bundle}

For the one target HostExW, construct the original compatible DesignChoices on Cdes=WithTw(WithSm Cprime), where Cprime=WithChN(A intersect target class), A is the one Built construction class, and cT(eps,y)=cTmin(y). Retain every original field, tuple, word, metric and common bundle duty. DSZ10 and RMP11 give size/rank properties of class members, not member existence. The original RecExP/RecAtG and common object/row route remains coupled. Construct ONE original common design family R and its recursive/geometric half with original Inputs/Pins/Holds/ExAt/InGeo and required CoreP. PKT now reconstructs full exists\_\allowbreak pkT, rowPk\_\allowbreak nonempty and exists\_\allowbreak moved with Copy/Coh on those original given inputs and a supplied fitting later stage: one catalogue, both ChT, complete LfT, actual integral lattices and all metric/measure/period constants are retained; one numerical record precedes every eta and V/S/N/G. RDR supplies one Side0 before all readers, ReadPk/pins and full ReadL with exact o.eta=q.eta. The initial actual geometry, Fits/division/central conditions, primitive rational tests and eta agreement must still be produced. Inherited CSR supplies finite columns, one ReadC before later twists, period/evaluation and scalar/error transports, RowSideDetermined and same-R Bundle/BTw/BWit/BunU consumers. RPI supplies exact RowPairInt for actual pairs and finite translates with original mu\_\allowbreak b/fd0; first literal DivDef3 at the same updated R,F,Y,Z,C,sEE follows from its retained original readings, Pins/Holds, stored nonzero IsItW chain, five real-group predicates and group tie. Remaining common fields include host/frame, MultEq nonzero-part support, that original nonzero chain/group tie, eigensystem/carrier/parts and all-embedding logLam, HD/d3b, NEb, actual rational tests and base rows. Retain every simultaneous family hypothesis and eventual threshold. Stronger second DivRows3 remains separate; universal quotient-height CuspRapidDecay and actual common member existence are not inferred from the comparison constructors. Version 0.41 supplies the finite phase-correct rational theta route: exact coefficient extension, split representatives, square-law/PhiIs, Pins from CoversAt, class-sum/point-identity transfer to OneMapOverL, actual-good-projector RatG and both-level LiftTestB. Its original higher covers, all-tuple class-sum uniqueness including phase rescaling, one PointIdentCore before O, base ClassSumsSpan and geometric LatFull remain explicit producers. MLat is available from the original LatFull input. This sharpens the rational-test work without claiming actual common geometry, Fits or eta agreement. Version0.46 proves finite generation of the original whole integral cusp lattice by actual integral-frame evaluation. It also proves the exact identity-cover, analytic unfolding, original-field coefficient descent, global clearing-scalar and both-level lattice-transport implications. Full spanning retains the supplied original IsNab connection and the five explicit principal-chart, level-system, raw-Gaussian, short-rule/dictionary and ramified-jet producers. The higher tuple covers, actual common class sums, PointIdentCore and ClassSumsSpan remain separate. CURRENT VERSION0.47 QUALIFICATION: Construct one original common geometric family and the higher CoversAt/all-tuple class sums and uniqueness with original phase rescalings, one PointIdentCore with common lambda,z, and base ClassSumsSpan. CNC now supplies the full original IsNab; PCL supplies the conditional trace-dual principal-lattice and quantitative normalization. Actual uniformization/HasPrinc and the five principal-chart, finite-level-system, Gaussian/reproduction, short-rule/simultaneous-dictionary and ramified-jet producer groups remain. Keep all original both-level, coefficient-field, Fits and eta agreements. Completed new original producers: ICL1-ICL21, CNC1-CNC16, PCL1-PCL12. Any earlier narrative naming these completed producers as outstanding is historical; this qualification is the current scope authority.

\emph{New completed inputs.} ICL1-ICL21, CNC1-CNC16, PCL1-PCL12

\emph{Current remaining subfrontier.} Construct one original common geometric family and the higher CoversAt/all-tuple class sums and uniqueness with original phase rescalings, one PointIdentCore with common lambda,z, and base ClassSumsSpan. CNC now supplies the full original IsNab; PCL supplies the conditional trace-dual principal-lattice and quantitative normalization. Actual uniformization/HasPrinc and the five principal-chart, finite-level-system, Gaussian/reproduction, short-rule/simultaneous-dictionary and ramified-jet producer groups remain. Keep all original both-level, coefficient-field, Fits and eta agreements.



\paragraph{22. LOCAL\_\allowbreak NBUILT\_\allowbreak COMMON — C — Actual local specialization on the common selected family}

Construct, with the required closed-class property on the same Des(C,3), one eventual family of actual NBuilt certificates for the original Fm.D(x,i), retaining the certificate before forgetting to TSel. HLP now reconstructs full HeckeLift0, all-index commutation, HeckeStd0 and HeckeIntTw from the original MtInput and complex Hecke facts, in the full ambient GenAt range k<=nV. Every maximal ideal is an actual supported residual character class; literal MtSub is the O-span of projections of the original MtLat generators, is finite free, and localizes onto the exact NCov.part. Literal Psi has its exact range and off-space projection error, both original span assertions hold, and ActEx/ActHK are constructed on the same lattice under original semisimplicity. Original MjLat relation quotients, universal coefficient-map relation fields, selected tame generators, coefficient embeddings and all identity transports give the complete nonempty\_\allowbreak spec constructor on its full printed input list. A single compositum combines the four original coefficient requirements with big fixed beforehand. Whole ndSpec\_\allowbreak exists still needs actual common geometric/choice MtInput, complex HeckeStdPlain and the no-Jordan producer, original F3 bridges with valuation/GoodCoef/iq, original F4 common-overfield reading of the entire coefficient field, and exact arithmetic character/reciprocity/admissibility suppliers. Original AutSpec, Galois specialization and the remaining NBuilt fields on that same installation also remain. Do not replace the actual projected lattice or maps by abstract objects with the same dimensions.



\paragraph{23. SECOND\_\allowbreak LAYER\_\allowbreak COMPATIBLE\_\allowbreak SELECTION — C — Original second-layer selector and class duties}

For the same A, supply Duty on A and a Closed common original TSel family whose selected records are the original true-column rec7 records (or have the entire dependent tuple transported), with their real witness and original coefficient-field root data. Existence/compatibility remains arithmetic; merely abstractly isomorphic fields or independently selected columns do not suffice. CSR21–22 now construct the same-R Bundle/BTw/BWit/BunU outputs from the exact selected realization/row-side packet and NEb, with the original EvP/ParEv quantifier order. Actual common selection, the five stated realization packets and NEb remain upstream; none is replaced by an independently chosen compatible record. Version 0.38 closes the selected analytic extension and reconstructs the first updated DivDef3 consumer; it leaves the original common selection, geometric readings, rational tests, base rows, host/parts/NEb and exact retained DI premises upstream, with unchanged EvP/ParEv thresholds. Version 0.39 supplies the original given-data reader constructors and exact RowPairInt/first DivDef3 on original premises. It leaves one actual common geometric/comparison/test member, eta agreement, compatible selection, rational tests, base rows, host/parts/NEb and original EvP/ParEv thresholds upstream. Version 0.40 also supplies original comparison/RowPk/fitting-stage Moved algorithms and exact common data transport. The actual initial geometric member, fitting tests, eta agreement and compatible second-layer selector with original eventual thresholds remain upstream. Version 0.41 supplies the finite phase-correct rational theta route: exact coefficient extension, split representatives, square-law/PhiIs, Pins from CoversAt, class-sum/point-identity transfer to OneMapOverL, actual-good-projector RatG and both-level LiftTestB. Its original higher covers, all-tuple class-sum uniqueness including phase rescaling, one PointIdentCore before O, base ClassSumsSpan and geometric LatFull remain explicit producers. MLat is available from the original LatFull input. This sharpens the rational-test work without claiming actual common geometry, Fits or eta agreement.



\paragraph{24. BASE\_\allowbreak THETA\_\allowbreak READINGS — C — Actual base reading on the common bundle}

For every family in the original class, eventually on the SAME selected Bundle and for each i<m, supply the actual R,F,Read,ReadL, weak models/Weil data and phase-one pins. Identify their LiftTie.theta, mtRead, mjRead and ThDef3 basis with that Bundle.P.A3, P.Zs, P.S0 and its stored maps. Exact field, embeddings, metrics, measures, row family and weak choices are retained. This remains one part of the coupled BTH12 producer. PKT now supplies the complete original comparison packet, RowPk and fitting-stage Moved with Copy/Coh on original given geometry, Inputs/Pins/Holds/ExAt/InGeo and Fits. Its ChT/LfT/Reps, actual lattice carriers, original Haar and norm constants, period maps, representatives and all-place RrT are constructed together; its one numerical record precedes every eta/V/S/N/G. RDR gives the original level-zero side, moved-Read and ReadPk/pin constructors on those objects, with one Side0 before every reader and all twelve Read ties. Full ReadL retains the exact o.eta(j,v)=q.eta(j,v) guard for every 1<=j<=m and place v; arbitrary Over does not imply it. Inherited CSR fixes one ReadC before all TwArgs, transports both metrics and all thirteen reading maps, and gives actual period/evaluation with kappa=(c*nP)\textasciicircum{}(-1/2). RPI supplies exact row-pair integrability, so first updated DivDef3 follows from the retained original Pins/Holds, nonzero chain, five real-group predicates and group tie. Initial theta extension remains complete. Remaining work is actual initial common geometry and its original predicates, Fits/primitive rational tests, eta agreement and outer family/eventual/same-Bundle duties, reFace compatibility, MultEq support, original chain/group data, actual base rows and any stronger DivRows3 theta square. No new compatible reading, family threshold or Haar normalization is substituted. Version 0.41 supplies the finite phase-correct rational theta route: exact coefficient extension, split representatives, square-law/PhiIs, Pins from CoversAt, class-sum/point-identity transfer to OneMapOverL, actual-good-projector RatG and both-level LiftTestB. Its original higher covers, all-tuple class-sum uniqueness including phase rescaling, one PointIdentCore before O, base ClassSumsSpan and geometric LatFull remain explicit producers. MLat is available from the original LatFull input. This sharpens the rational-test work without claiming actual common geometry, Fits or eta agreement.



\paragraph{25. BASE\_\allowbreak THETA\_\allowbreak RATIONAL\_\allowbreak TESTS — C — Simultaneous original base rational tests}

On the SAME original Bundle/reading/theta map, for every i<m supply one finite separating complex-linear test family whose values lie in the fixed embedded coefficient field on theta of EVERY original rational source vector and on EVERY stored target-coordinate basis vector. BTH8/BTH11 use this same LiftTestB for rational adjoint image and matrix coordinates. RTG now reconstructs prescribed-field coefficient extension, original split class representatives/RepsOK, the self-dual split-lattice square law and finite residue PhiIs; every Pins field follows from the actual CoversAt. With the base coefficient record literally fixed, exact class-sum existence/uniqueness including phase-rescaled O, PointIdentCore with common lambda,z before every O, and base ClassSumsSpan imply all-tuple finite lift span and OneMapOverL. The phase is exactly cj(twAt embLL A3.A.v0)(R.phs i). RatG uses the actual good projector eOp(y), and full LiftTestB follows at BOTH original levels. These finite implication proofs are complete. Remaining original producers are the prescribed higher covers, those actual class sums/uniqueness, PointIdentCore, base ClassSumsSpan and the base/common geometric datum; exact MLat is supplied only from original LatFull, whose geometric production remains separate. Finite dimension or an unrelated rational theta model cannot replace these same-object identities. Original outer eventual-family/Bundle quantifiers remain. Version0.46 proves finite generation of the original whole integral cusp lattice by actual integral-frame evaluation. It also proves the exact identity-cover, analytic unfolding, original-field coefficient descent, global clearing-scalar and both-level lattice-transport implications. Full spanning retains the supplied original IsNab connection and the five explicit principal-chart, level-system, raw-Gaussian, short-rule/dictionary and ramified-jet producers. The higher tuple covers, actual common class sums, PointIdentCore and ClassSumsSpan remain separate. CURRENT VERSION0.47 QUALIFICATION: Construct one original common geometric family and the higher CoversAt/all-tuple class sums and uniqueness with original phase rescalings, one PointIdentCore with common lambda,z, and base ClassSumsSpan. CNC now supplies the full original IsNab; PCL supplies the conditional trace-dual principal-lattice and quantitative normalization. Actual uniformization/HasPrinc and the five principal-chart, finite-level-system, Gaussian/reproduction, short-rule/simultaneous-dictionary and ramified-jet producer groups remain. Keep all original both-level, coefficient-field, Fits and eta agreements. Completed new original producers: ICL1-ICL21, CNC1-CNC16, PCL1-PCL12. Any earlier narrative naming these completed producers as outstanding is historical; this qualification is the current scope authority.

\emph{New completed inputs.} ICL1-ICL21, CNC1-CNC16, PCL1-PCL12

\emph{Current remaining subfrontier.} Construct one original common geometric family and the higher CoversAt/all-tuple class sums and uniqueness with original phase rescalings, one PointIdentCore with common lambda,z, and base ClassSumsSpan. CNC now supplies the full original IsNab; PCL supplies the conditional trace-dual principal-lattice and quantitative normalization. Actual uniformization/HasPrinc and the five principal-chart, finite-level-system, Gaussian/reproduction, short-rule/simultaneous-dictionary and ramified-jet producer groups remain. Keep all original both-level, coefficient-field, Fits and eta agreements.



\paragraph{26. SLR\_\allowbreak BASE\_\allowbreak ROWS\_\allowbreak PRODUCER — C — Actual base row identities on that same bundle}

For the SAME selected Bundle, common original fields and readings used by the other three BTH12 inputs, supply actual RowsHold0\_\allowbreak 3 uniformly eventually for every family. Its original primitive row realization, spectrum/type and basis identifications must be retained. BTH supplies a theta consumer only; it does not construct the base row identity. No separate witness may be chosen merely to satisfy this row requirement.



\clearpage
\section{Kato vanishing and the actual type-II rigidity producer}
\label{sec:type-ii-rigidity-v47}

We prove the original \texttt{Kato\allowbreak Two} and \texttt{Rig\allowbreak II} assertions.
The latter concerns the full nine-dimensional polarized adjoint of an
actual rank-three deformation point.  Its modularity hypothesis is
retained exactly: every rank-two extension specified by
\texttt{Dat.\allowbreak Is\allowbreak Ext} is modular for the supplied field isomorphism
\(\iota:\overline{\mathbf Q}_\ell\simeq\mathbf C\).
We construct one such extension and prove its required properties.
In particular, no adjoint-vanishing, newform-to-base-change, or Kato
predicate is left as an unnamed producer.

All source references are at commit
\texttt{7cbf7a0c32fca1badf05931879a37c3f8a2b03b4}.
The original tied period-ring model, finite Galois cohomology,
class-field theory, and the explicitly cited published automorphic
theorems remain classical foundations.  The proofs below do not
reconstruct those foundational theorems.  They reconstruct their
applications, coefficient comparisons, and every finite-dimensional
consumer needed for these two endpoints.

Write \(G=G_{\mathbf Q}\), \(H=G_K\), and
\(\epsilon:G\to\mathbf Z_\ell^\times\) for the cyclotomic
character.  Frobenius in an attachment formula is geometric, so
\(\epsilon(F_p)=p^{-1}\).  At a fixed point let
\(\eta:G\to\{\pm1\}\) be the quadratic character of the
original imaginary quadratic field \(K\), and let \(c\) be the
specified complex conjugation.  The original coefficient ring is
\(\mathcal O\); it is not replaced when constructing the point.
We use a finite extension of \(\mathbf Q_\ell\) only to apply
classical rational cohomology theorems, and return each primitive
to the original algebraic closure.

\subsection{The precise Selmer and Kato comparison}

% BEGIN KTR1
\paragraph{KTR1. Finite coefficients, functorial local conditions, and finite discrete Selmer groups.}
\label{claim:KTR1}
The source's \texttt{Bloch\allowbreak Kato\allowbreak Selmer\allowbreak Vanishes} is rational
Bloch--Kato vanishing with its full original cocycle scope.  The
following three operations preserve that scope: a continuous
equivariant linear map, finite coefficient extension, and restriction
to a finite number-field extension.  Moreover, if \(V\) is a
finite-coefficient representation and the discrete Bloch--Kato
Selmer group attached to a stable lattice \(T\subset V\) is
finite, then \(H^1_f(L,V)=0\).

\emph{Proof.}
ADK1 puts a continuous algebraic-closure-valued representation and
any one continuous cocycle over one finite coefficient field.
This uses compactness and Baire's theorem, not an assertion that
an arbitrary compact subset of the algebraic closure is contained
in a finite extension.  ADK5 identifies the source's chosen
inertias and decomposition-group graphs with the classical groups.
The graph identification is continuous because the inverse image
of the stabilizer of a finite set of global algebraic elements
fixes a finite local extension.  With BX1's tied period model, the
local condition at \(\ell\) is exactly the kernel of
\[
 H^1(L_v,V)\longrightarrow
 H^1(L_v,V\otimes_{\mathbf Q_\ell}B_{\rm cris}),
 \tag{KTR.1}
\]
and away from \(\ell\) it is the kernel of restriction to inertia.
Positive cohomology at real places vanishes by averaging.

If \(f:V\to V'\) intertwines the actions, apply \(f\) to a
cocycle and to each inertia primitive.  At \(\ell\), apply
\(f\otimes1\) to a crystalline primitive.  On a pure tensor,
\[
 (f\otimes1)(g v\otimes g b)=g f(v)\otimes g b,
\]
so this map commutes with the diagonal action, and the identity
extends by linearity.  Restriction of a primitive is again a
primitive, proving the restriction assertion.  Local extension of
scalars causes no hidden enlargement of the kernel: every given
algebraic-closure primitive uses finitely many vector coefficients,
contained in some finite \(E'/E\).  An \(E\)-linear retraction
\(E'\to E\) fixing \(1\), tensored with \(1\) on
\(B_{\rm cris}\), brings that primitive back to \(E\).
Different places may use different \(E'\).

For the last assertion choose a finite set \(S\) containing the
ramified places, the places above \(\ell\), and the infinite
places.  A rational Selmer class factors through \(G_{L,S}\):
outside \(S\) the action is trivial on inertia, so its
unramified condition says that its restriction to inertia is
zero, and the cocycle descends through their closed normal closure.
The finite-coefficient global cohomology foundation in DX1 gives
that \(H^1(G_{L,S},T)\) is a finitely generated
\(\mathbf Z_\ell\)-module.  Equivalently this follows from
finite cohomology of \(T/\ell T\), inverse limits, and
topological Nakayama.  Its image in \(H^1(G_{L,S},V)\) is a
bounded lattice.  The map to \(H^1(G_{L,S},V/T)\) sends the
rational Selmer group into the discrete Selmer group whose local
conditions are the images of the rational local conditions.
Its kernel is contained in this bounded lattice, by the long
exact sequence for \(0\to T\to V\to V/T\to0\).
The image of a divisible group in a finite group is zero.
Thus the whole rational vector space \(H^1_f(L,V)\) is contained
in a bounded lattice.  A nonzero vector and all its
\(\ell^{-m}\)-multiples cannot be so contained.  The vector
space is zero, as asserted.  This argument does not require the
order of the finite discrete Selmer group to be prime to \(\ell\).
% END KTR1

% BEGIN KTR2
\paragraph{KTR2. The exact noncentral critical Kato application.}
\label{claim:KTR2}
Let \(e\ge1\), let \(f\) be an original normalized newform
of weight \(k=2e+1\), and suppose the continuous representation
\(W\) satisfies
\[
             f.\mathrm{Att}\bigl(\iota,W\epsilon^{-e}\bigr).
 \tag{KTR.2}
\]
For every finite abelian extension \(L/\mathbf Q\),
\(H^1_f(L,W)=0\), with the exact coefficient and local meaning
of KTR1.

\emph{Published theorem and normalization check.}
Kato, \emph{p-adic Hodge theory and values of zeta functions of
modular forms}, Ast\'erisque 295 (2004), Theorem 14.2(1),
states finiteness of \(\operatorname{Sel}(L,T)\) for a stable
lattice in \(V_{F_\lambda}(f)(r)\), when \(L/\mathbf Q\)
is finite abelian, \(1\le r\le k-1\), and \(2r\ne k\).
It applies to every finite place \(\lambda\) of the eigenvalue
field, with no ordinarity assumption and no lower bound on its
residue characteristic.  The paper includes the CM case; the
separate treatment of that case does not remove it from the
statement.  The local Selmer convention is the one in \S14.1:
the local conditions on the discrete module are images of the
unramified rational conditions away from \(\ell\) and of the
\(B_{\rm cris}\) kernel at \(\ell\).
The primary text is
\url{https://www.numdam.org/item/AST_2004__295__117_0/}.

The original good-prime eigenrelations, normalization
\(a_1=1\), nebentype transformation, and exclusion of every
smaller conductor are the newform data used in RNC2's
adelization.  Strong multiplicity one fixes the cuspidal
eigensystem; the conductor condition selects its one-dimensional
newvector line.  Thus this is the classical normalized newform
to which Kato applies, with its finite eigenvalue field, including
the cases where its level is divisible by \(\ell\).

The cohomological realization \(V_{F_\lambda}(f)\) has geometric
Frobenius polynomial, by the explicit inverse-arithmetic-Frobenius
formula (14.10.3) on printed page 240,
\(X^2-a_p(f)X+\chi_f(p)p^{k-1}\).
This is exactly the polynomial in the source's
\texttt{Newform.\allowbreak hecke\allowbreak Poly}; it is not its reciprocal polynomial.
The source \texttt{Att} includes semisimplicity, unramifiedness
at all but finitely many primes, and this same polynomial for
every prolongation and every arithmetic Frobenius inverted.
The usual newform Galois realization and (KTR.2), after finite
coefficient descent, therefore agree by Chebotarev and the
semisimple character criterion (CAP6).  Hence
\[
          W\simeq V_{F_\lambda}(f)(e).
 \tag{KTR.3}
\]
The chosen \(\lambda\) is the one induced by the fixed
\(\iota^{-1}\) on the eigenvalue field.  If the actual
coefficient field is larger than \(F_\lambda\), restriction
of scalars is a finite direct sum, so the same vanishing follows
after extension.  There is no change of \(\iota\).
Now \(1\le e\le2e=k-1\) and \(2e\ne2e+1=k\).
Kato's theorem and KTR1 prove (KTR.3)'s rational Selmer vanishing.
This is an application at \(r=e\), rather than at the excluded
central value.  Its proof supplies the actual coboundary for
each original continuous cocycle.
% END KTR2

% BEGIN KTR3
\paragraph{KTR3. Finite-character descent and the full original KatoTwo.}
\label{claim:KTR3}
Under KTR2, for every continuous finite-order abelian character
\(\varphi:G\to\overline{\mathbf Q}_\ell^\times\),
\[
                       H^1_f(\mathbf Q,W\varphi)=0.
 \tag{KTR.4}
\]
Consequently the full original \texttt{Kato\allowbreak Use} and
\texttt{Kato\allowbreak Two} hold, with every coefficient closure, embedding,
level, weight, and allowed coefficient prime in their statements.

\emph{Proof.}
Let \(L\) be the finite abelian extension fixed by
\(\ker\varphi\).  On \(G_L\) the twist is \(W\), so
restriction of a Selmer cocycle is a coboundary by KTR2 and
KTR1.  Subtract that coboundary.  The remaining cocycle vanishes
on \(G_L\); its values are invariant under \(G_L\), and it
descends to \(\operatorname{Gal}(L/\mathbf Q)\).
Every cocycle of a finite group with characteristic-zero vector
coefficients is a coboundary.  Explicitly, for a cocycle \(a\)
on a finite group \(Q\), put
\(v=-|Q|^{-1}\sum_{q\in Q}a(q)\); then
\(a(g)=g v-v\).  Thus restriction is injective in degree one.
No condition \(\ell\nmid[L:\mathbf Q]\) is needed.

For \texttt{Kato\allowbreak Use}, the original condition
\texttt{Kills\allowbreak Mu}\((M,\varphi)\) makes \(\varphi\) factor
through \(\operatorname{Gal}(\mathbf Q(\mu_M)/\mathbf Q)\),
a finite abelian quotient.  It is continuous and satisfies the
condition just proved.  The original assumption \(M\ge1\)
and the original odd-prime quantifier are left intact.

For \texttt{Kato\allowbreak Two}, retain all fields of the source's
\texttt{Used} record: \(\ell\ne2\), \(e\ge1\),
\(\ell-1\nmid2e\), imaginary quadratic \(K\),
\(\ell\nmid|d_K|h_K\), the even positive level and its
coprimality, both character conditions, the complex-conjugation
condition, the exact \texttt{Level} and \texttt{Ord\allowbreak Poly},
determinant, inertial trace, residual trace, the specified flag,
and irreducibility on the infinite cyclotomic kernel.
The \texttt{Level} field includes continuity of \(W\).
For this conclusion the extra hypotheses are unnecessary:
the character \(\theta\) has kernel \(G_K\), hence is the
continuous quadratic character, and we may apply (KTR.4) to
\(\varphi=1\) and \(\varphi=\theta\).
This gives the source's two literal column representations
\(\operatorname{colRep}(\operatorname{tw}W1)\) and
\(\operatorname{colRep}(\operatorname{tw}W\theta)\).
The claim is the complete implication with the original
\texttt{Used} antecedent, not a demand for a stronger antecedent.
It includes \(\ell=3\) wherever the original antecedent is
satisfiable; the published theorem itself excludes no such prime.
% END KTR3

\subsection{The actual point and the smaller local calculation}

For KTR4--KTR13 fix all the original data in \texttt{Rig\allowbreak II}:
\(D:\texttt{Dat}(\Omega,K,\mathcal O,e,\psi)\),
\(U:D.\texttt{ULoc}\), the marked element \(g_1\), and
an actual \(r\) satisfying \(D.\texttt{Is\allowbreak Pt}(U,g_1,r)\).
Assume \(\texttt{Type\allowbreak II}(\eta,r)\).  Let \(R(g)\) be the
matrix part of \(r(g)\).  On \(H=\ker\eta\) it is an
ordinary representation; its only possible off-diagonal blocks
are the corner positions \(0,2\).  Write
\[
 R(h)=\begin{pmatrix}a_h&0&b_h\\0&\alpha_h&0\\c_h&0&d_h\end{pmatrix},
 \qquad B_h=\begin{pmatrix}a_h&b_h\\c_h&d_h\end{pmatrix}
 \quad(h\in H).
 \tag{KTR.5}
\]
The corner representation is irreducible.  The point is continuous,
its entries are integral and entrywise congruent to the original
\(r_0(\eta,\Psi)\), and \(R(c)=1\).
These are fields of \texttt{Is\allowbreak Pt} and \texttt{Normalised}.
Its away-from-\(\ell\) killing condition and its map through
the original local ideal \(I_t\) will both be used.

% BEGIN KTR4
\paragraph{KTR4. The original local ideal forces an ordered inertia flag on a finite-index subgroup.}
\label{claim:KTR4}
Choose the \(\ell\)-adic embedding stored in \(D\), let
\(I\subset G\) be its inertia group, let \(\tau\) be the
stored finite local character, and put \(J=I\cap\ker\tau\).
Then \(I\subset H\), \(\tau^2|_I=1\), and on \(J\)
the actual \(R\) has an ordered flag with characters
\[
                      (\epsilon^e,1,\epsilon^{-e}).
 \tag{KTR.6}
\]
Every element of \(I\) preserves this ordered flag.

\emph{Proof.}
The assumptions \(\ell\nmid d_K\) and \(K/\mathbf Q\)
quadratic put inertia in \(H\).  The original \texttt{GD}
field \((\operatorname{comp}_\ell\psi)^2=1\) gives the
assertion about \(\tau\), and the stored local character
identity is \(\Psi|_I=\epsilon^e\tau\).
The embedding comparison in \texttt{Dat\allowbreak Loc} follows at finite
Galois levels by extending the specified \(\ell\)-adic
valuation and taking its decomposition group; inverse limits
identify the stored local homomorphism, inertia, and cyclotomic
character with these groups.  This uses the actual embedding
equation \(s(jx)=j(\iota_s x)\) of \(D\), so it does not
choose a different local action.

Unfold the defining ideals.  The ordinary ideal is contained in
its saturation \(I_Y\) (use exponent zero), which is contained
in its radical \(I_{\rm cr}\).  Every minimal prime used in
\(I_t\) contains \(I_{\rm cr}\), so
\[
                    I_{\rm ord}\subset I_t.
 \tag{KTR.7}
\]
This also holds if that set of primes is empty.  The local map
in \texttt{Is\allowbreak Pt} kills \(I_t\); consequently for all
\(s,s_0,s_1,s_2\in J\) it gives
\[
 \det(T-R(s))=\prod_{a=0}^2(T-\epsilon(s)^{e(1-a)}),
 \qquad
 \prod_{a=0}^2\bigl(R(s_a)-\epsilon(s_a)^{e(1-a)}1\bigr)=0.
 \tag{KTR.8}
\]
The order of the three factors is as displayed.

The image of \(\epsilon\) on rational \(\ell\)-inertia is
all \(\mathbf Z_\ell^\times\).  Indeed each cyclotomic
\(\ell\)-power extension of \(\mathbf Q_\ell\) is totally
ramified, and compactness passes the finite surjections to the
inverse limit.  Choose \(u\in I\) with
\(\epsilon(u)=1+\ell\) and replace it by \(s_*=u^2\).
Then \(s_*\in J\) and \(z=\epsilon(s_*)^e\) has infinite
order.  The three eigenvalues \(z,1,z^{-1}\) are distinct.
Diagonalize \(R(s_*)\) in precisely that order, writing
\(D_* =\operatorname{diag}(z,1,z^{-1})\).

In that basis insert \((s,s_*,s_*)\),
\((s_*,s,s_*)\), and \((s_*,s_*,s)\) in the second
equation of (KTR.8).  If \(N=R(s)\) and
\(p=(\epsilon(s)^e,1,\epsilon(s)^{-e})\), the equations are
\[
 \begin{split}
 &(N-p_0)(D_*-1)(D_*-z^{-1})=0,\\
 &(D_*-z)(N-p_1)(D_*-z^{-1})=0,\\
 &(D_*-z)(D_*-1)(N-p_2)=0.
 \end{split}
 \tag{KTR.9}
\]
Here a scalar denotes its scalar matrix.  The first equation
gives \(N_{00}=p_0,N_{10}=N_{20}=0\); the second gives
\(N_{11}=p_1\); the third gives
\(N_{21}=0,N_{22}=p_2\).  The factors divided out are
differences among \(z,1,z^{-1}\), all nonzero.  Thus the
matrices on \(J\) are upper triangular with (KTR.6).

Finally \(J\) is normal in \(I\) and \(\epsilon\) is
unchanged by conjugation.  For \(g\in I\), the matrices of
\(s_*\) and \(g s_*g^{-1}\) are upper triangular with the
same distinct ordered diagonal.  Their intertwiner \(R(g)\)
is upper triangular: the \((2,0)\) equation first kills that
entry, then the \((2,1)\) and \((1,0)\) equations kill those
entries.  Their coefficients are again nonzero diagonal
differences.  This proves preservation of the flag by \(I\).
Only (KTR.7), rather than the stronger full local crystalline
component theorem, has been required.
% END KTR4

% BEGIN KTR5
\paragraph{KTR5. The actual middle and determinant characters are trivial on the quadratic field.}
\label{claim:KTR5}
For every \(h\in H\), the matrix in (KTR.5) satisfies
\[
                          \alpha_h=1,\qquad \det B_h=1.
 \tag{KTR.10}
\]

\emph{Proof.}
Both \(\alpha\) and \(\delta=\det R|_H\) are continuous
characters.  Integral entrywise residual closeness implies
\(\alpha(h)\equiv1\) and \(\delta(h)\equiv1\) modulo
the maximal ideal for every \(h\in H\): the residual diagonal
on \(H\) is \((\overline\Psi,1,\overline\Psi^{-1})\).
We shall prove both characters unramified everywhere over \(K\).

First consider the inertia group in KTR4.  The condition
\(\ell-1\nmid2e\) supplies a unit
\(a\in\mathbf Z_\ell^\times\) with
\(\bar a^{2e}\ne1\).  Choose \(s\in I\) with
\(\epsilon(s)=a\) and put \(t=s^2\in J\).
Then \(\overline{\epsilon(t)^e}\ne1\).
For any \(x\in J\), the middle-block eigenvalue
\(\alpha(x)\) is one of the roots in (KTR.8).
If \(\overline{\epsilon(x)^e}\ne1\), residual closeness
forces \(\alpha(x)=1\).  In the other case apply this
argument to \(xt\) and to \(t\), and divide their character
values.  Hence \(\alpha|_J=1\).  Also (KTR.8) gives
\(\delta|_J=1\) directly.  For \(x\in I\), its square
lies in \(J\).  Each of \(\alpha(x),\delta(x)\) is thus
a square root of unity congruent to \(1\), and \(\ell\ne2\)
makes it equal to \(1\).  The other places above \(\ell\)
are conjugate in \(G\).  Conjugation by \(H\) preserves
the characters \(\alpha,\delta\).  Conjugation by \(c\)
sends each to its inverse: \(R(c)=1\) and the disconnected
action is \(R(h)\mapsto J_3R(h)^{-t}J_3\), whose middle
entry and determinant are respectively
\(\alpha(h)^{-1},\delta(h)^{-1}\).
Triviality therefore transfers to every place above \(\ell\).

At a rational prime \(p\ne\ell\), restrict to
\(H\cap I_p\).  The cyclotomic character is unramified and
the residual model is the diagonal representation
\((\psi,1,\psi^{-1})\).  The finite character \(\psi\)
has order prime to \(\ell\), by \texttt{GD}.  The point's
\texttt{Kill} condition says exactly that the kernel of this
finite residual representation is killed by \(R\).
Thus \(R\) factors on that inertia through a finite group
of order prime to \(\ell\).  The same holds for \(\alpha\)
and \(\delta\).  A prime-to-\(\ell\) root \(b\) with
\(b\equiv1\) is \(1\): in
\(b^m-1=(b-1)(1+\cdots+b^{m-1})\) the second factor is a
unit if \(\ell\nmid m\).  Both characters are therefore
trivial on these inertia groups too.

For completeness, the residual-kernel description just used does
not lose characteristic-zero information.  Reduction is injective
on prime-to-\(\ell\) roots of unity by the preceding argument;
on \(H\cap I_p\) the kernel of the residual diagonal is
therefore the kernel of \(\psi\).  If desired, the whole
finite inertial action is conjugate to that diagonal action:
the averaging intertwiner
\(\sum_{q}R_0(q)^{-1}R(q)\) is congruent to \(|Q|1\)
and hence invertible.  Only the prime-to-\(\ell\) order,
not this stronger conjugacy, is needed above.

Class-field theory identifies the everywhere unramified abelian
quotient of \(G_K\) with the ideal class group, since \(K\)
has no real places.  Thus \(\alpha\) and \(\delta\)
factor through a group of order \(h_K\).  The assumption
\(\ell\nmid h_K\) and residual closeness again force them
to be trivial.  As \(\det R=\alpha\det B\), (KTR.10)
follows.  This uses the original \(K\), its class number, and
the original residual point, with no auxiliary coefficient-ring
or modularity hypothesis.
% END KTR5

% BEGIN KTR6
\paragraph{KTR6. The explicit original rank-two extension.}
\label{claim:KTR6}
There is a continuous representation \(W:G\to
\mathrm{GL}_2(\overline{\mathbf Q}_\ell)\) satisfying
the literal \(D.\texttt{Is\allowbreak Ext}(r,W)\), and
\[
 W|_H=B,\qquad W(c)=D_2:=\operatorname{diag}(-1,1),
 \qquad \det W=\eta.
 \tag{KTR.11}
\]
For every \(g\in G\), \(R(g)\) is a block embedding of
an \(\mathrm{SL}_2\) matrix with middle entry \(1\).

\emph{Proof.}
Every \(g\notin H\) is \(hc\) with \(h\in H\).
The semidirect multiplication and \(R(c)=1\) give
\(R(hc)=R(h)\), so KTR5 proves the last assertion.
For \(g\in G\), let \(B_g\) be this corner matrix and put
\[
 W(g)=B_gD_2^{(1-\eta(g))/2}.
 \tag{KTR.12}
\]
Write \(J_2=\left(\begin{smallmatrix}0&1\\1&0\end{smallmatrix}\right)\).
For every \(B\in\mathrm{SL}_2\), direct multiplication gives
\(J_2B^{-t}J_2=D_2BD_2\).  The disconnected action in the
original \(\mathrm{GG}_3\) restricts on the corner to its
left-hand side.  Its semidirect multiplication therefore makes
(KTR.12) a homomorphism.  The two cosets of \(H\) are open;
on each coset its entries are fixed signed entries of \(R\).
This proves continuity.  Formula (KTR.11) follows immediately,
including \(\operatorname{tr}W(c)=0\).
The definition of \texttt{Is\allowbreak Ext} asks exactly for the equality
with the corner on \(H\) and this trace at \(c\).
Thus this is an actual extension covered by the original universal
modularity hypothesis.  Its restriction to \(H\) is irreducible
by \texttt{Type\allowbreak II}, hence \(W\) is irreducible too.
% END KTR6

% BEGIN KTR7
\paragraph{KTR7. An elementary two-dimensional self-twist criterion.}
\label{claim:KTR7}
Let \(V\) be an irreducible two-dimensional representation of
a group \(\Gamma\) over an algebraically closed field of
characteristic zero.  If \(N\lhd\Gamma\), \(\Gamma/N\)
is abelian, and \(V|_N\) is reducible, then there is a
nontrivial character \(\xi:\Gamma/N\to\{\pm1\}\) with
\(V\simeq V\xi\).

\emph{Proof.}
An \(N\)-stable line cannot be the unique such line, for
normality would make it \(\Gamma\)-stable.  Two distinct
stable lines give a diagonal decomposition
\(V|_N=\nu_1\oplus\nu_2\).
If \(\nu_1\ne\nu_2\), these are the only stable lines.
The action of \(\Gamma\) permutes them nontrivially, giving
\(\xi\); the matrix acting as \(1,-1\) on them intertwines
\(V\) and \(V\xi\).

If \(\nu_1=\nu_2\), all elements of \(N\) act by scalars.
Commutators of elements of \(\Gamma\) lie in \(N\), so
their matrices commute projectively.  Choose a nonscalar matrix
\(A\) in the image, possible by irreducibility.  For each
\(g\), conjugation sends \(A\) to \(\xi(g)A\).
Determinants give \(\xi(g)^2=1\), and conjugation shows that
\(\xi\) is a character trivial on \(N\).
If it were trivial, \(A\) would commute with the whole
irreducible representation and would be scalar, a contradiction.
Finally \(A\) itself is an invertible intertwiner for this
self-twist.  This proves the assertion without a semisimplicity
assumption on the initial restriction to \(N\).
% END KTR7

% BEGIN KTR8
\paragraph{KTR8. Irreducibility on the full infinite cyclotomic kernel.}
\label{claim:KTR8}
The actual \(W\) constructed in KTR6 is irreducible on
\(G_{K(\mu_{\ell^\infty})}\).

\emph{Proof.}
On the group \(J\) of KTR4, the first flag line has character
\(\epsilon^e\).  It lies in the corner plane: the middle
line is trivial by KTR5, whereas \(\epsilon(s_*)^e\ne1\).
The ordered flag is preserved by \(I\).  Its first line is
therefore an inertia-stable line for \(W\), with character
\(\kappa\), and \(\det W|_I=1\) gives the other diagonal
character \(\kappa^{-1}\).  In particular
\[
 \operatorname{tr}W(s)=\kappa(s)+\kappa(s)^{-1}\quad(s\in I),
 \qquad \kappa|_J=\epsilon^e.
 \tag{KTR.13}
\]
The character \(\kappa\) has infinite image, so choose
\(u\in I\) with \(\kappa(u)^8\ne1\).

Set \(N=\ker(\epsilon|_H)\).  If the restriction to \(N\)
were reducible, KTR7 would give a nontrivial quadratic self-twist
\(\xi\) trivial on \(N\).  The image of \(I\) under
\(\epsilon\) is all \(\mathbf Z_\ell^\times\), and
\(I\subset H\), so \(I\) surjects onto \(H/N\).
There is consequently \(s\in I\) with \(\xi(s)=-1\).
Equality of traces under the self-twist and (KTR.13) imply
\(\kappa(s)^4=1\).  The same reasoning applies to
\(su^2\), because \(\xi(u^2)=1\).  Dividing the two
fourth-power equalities gives \(\kappa(u)^8=1\), a
contradiction.  Thus the restriction to \(N\) is irreducible.
ADK2 identifies this actual closed subgroup with
\(G_{K(\mu_{\ell^\infty})}\); no finite cyclotomic level
has been substituted.  We have not needed to identify the finite
part of \(\kappa\) with the original \(\tau\) on all of
\(I\), nor to invoke a stronger local deformation theorem.
% END KTR8

\subsection{The corner adjoint and the nine-dimensional decomposition}

% BEGIN KTR9
\paragraph{KTR9. The actual modular extension has a cuspidal quadratic base change.}
\label{claim:KTR9}
Suppose the \(W\) in KTR6 satisfies the original
\(D.\texttt{Is\allowbreak Modular}(\iota,W)\).  There is a regular
algebraic cuspidal representation \(\Pi\) of
\(\mathrm{GL}_2(\mathbf A_K)\), polarized over \(\mathbf Q\),
whose normalized Galois representation is the actual \(W|_H\).
At an unramified place \(w\) with norm \(q_w\), its
normalized Satake roots \(b_{w,i}\) satisfy
\[
       \operatorname{Eig}(W(F_w))
                  =\{\iota^{-1}(q_w^{1/2}b_{w,i})\}_{i=1,2}.
 \tag{KTR.14}
\]

\emph{Proof and precise published input.}
Modularity supplies an original newform \(f\) of weight
\(2e+1\) with attachment (KTR.2).  RNC2 constructs its actual
cuspidal adelization \(\pi^0\) over \(\mathbf Q\) and
the regular algebraic normalization
\(\pi=\pi^0|\det|^{1/2}\), whose normalized Galois
representation is \(W\).  Its infinitesimal parameters are
\(e+\tfrac12,\tfrac12-e\), distinct because \(e\ge1\).
This is a statement about the given \(f,W,\iota\), not an
existential replacement of its Frobenius system.

Apply quadratic base change for \(\mathrm{GL}_2\).
The precise primary input is Langlands, \emph{Base change for
GL(2)}, Annals of Mathematics Studies 96 (1980),
\S2(A)--(F), Lemma 11.3(b), and Proposition 11.4:
a cuspidal representation over a characteristic-zero number
field has a cuspidal lift to a given quadratic extension unless
it is induced from a character of that extension; the lift
matches local base change, including infinity and all spherical
places.  It suffices to apply the theorem to the unitary
\(\pi^0\), then apply the same norm twist.
The primary text is
\url{https://publications.ias.edu/sites/default/files/base-change-for-gl2_rpl_2.pdf}.
In the exceptional induced case, twisting by \(\eta\) fixes
the representation: the inducing character's induction is
unchanged because \(\eta|_H=1\).  This also follows directly
from its unramified parameters and strong multiplicity one.

That exceptional case cannot occur here.  If
\(\pi\simeq\pi\eta\), its unramified polynomials and
CAP6 give \(W\simeq W\eta\).  An intertwiner then commutes
with \(W(H)\), so is scalar by its irreducibility, while at
\(c\) it must anticommute with \(W(c)\).  An invertible
scalar cannot do so in characteristic zero.  Thus the base change
\(\Pi\) is cuspidal.  Its archimedean parameter is restriction
of the real parameter to \(\mathbf C^\times\); at each
complex embedding it has the same two distinct half-integral
parameters.  This is regular algebraicity in the original
Capelli/Harish--Chandra convention checked in ADK3.

For the normalization, let \(a_1,a_2\) be the spherical Satake
roots of \(\pi_p\).  At a place \(w\mid p\) of residue
degree \(f_w\), base change gives \(a_1^{f_w},a_2^{f_w}\).
Since \(q_w=p^{f_w}\),
\(q_w^{1/2}a_i^{f_w}=(p^{1/2}a_i)^{f_w}\).
This proves (KTR.14) for every unramified place outside one
finite set, including inert places.

Finally \(\det W|_H=1\) and \(W\) extends to \(G\).
Thus the normalized Galois roots at \(w^c\) equal those at
\(w\), and form an inverse pair.  Translating (KTR.14) back
to Satake roots gives
\[
 \{b_{w^c,i}\}=\{q_w^{-1}b_{w,i}^{-1}\}.
 \tag{KTR.15}
\]
Class-field theory supplies an algebraic Hecke character
\(\lambda_0\) over \(\mathbf Q\) whose pullback to \(K\)
has \(\ell\)-adic character \(\epsilon|_H\); it is the
norm character fixed by these geometric-Frobenius values.
If necessary multiply it by the quadratic Hecke character of
\(K/\mathbf Q\) so that its infinity value at \(-1\) is
\((-1)^2=1\).  This does not change its pullback or (KTR.15).
Strong multiplicity one now gives
\(\Pi^c\simeq\Pi^\vee\otimes(\lambda\circ N_{K/\mathbf Q}\circ\det)\).
It is exactly the polarized pair required by BLGGT \S2.1;
its normalized Galois multiplier on \(H\) is
\(\epsilon^{-1}\lambda_\ell=1\).
BLGGT Theorem 2.1.1 now constructs the normalized attached
representation.  Its polynomials agree with (KTR.14), so
Chebotarev and the semisimple character criterion identify it
with \(W|_H\).  Thus the pairing hypothesis is established
before invoking that existence theorem.
% END KTR9

% BEGIN KTR10
\paragraph{KTR10. Vanishing for the complete polarized corner adjoint.}
\label{claim:KTR10}
On \(M_2(\overline{\mathbf Q}_\ell)\), let \(A\) act by
\[
 A(h)X=W(h)XW(h)^{-1}\quad(h\in H),
 \qquad A(c)X=-J_2X^tJ_2.
 \tag{KTR.16}
\]
Then \(H^1_f(\mathbf Q,A)=0\) in the full original sense.

\emph{Proof.}
We verify the pairing rather than infer it from self-duality
alone.  Put
\(S=\left(\begin{smallmatrix}0&1\\-1&0\end{smallmatrix}\right)\).
Every two-by-two matrix satisfies
\(W(g)^t S W(g)=\det W(g)S=\eta(g)S\).
Since \(W(c)=D_2\) and \(S D_2=J_2\), it follows that
\[
             W(h)^tJ_2W(chc^{-1})=J_2\quad(h\in H).
 \tag{KTR.17}
\]
The pairing \(J_2\) is symmetric and nondegenerate.
KTR9 supplies the exact regular algebraic polarized cusp and
its attached representation.  BLGGT Theorem 2.1.1 supplies the
same normalized pairing; its ratio with (KTR.17) commutes with
the irreducible \(W(H)\), hence is scalar.  The disconnected
adjoint formula is therefore exactly (KTR.16), as in ADK4:
the differential of inverse transpose is minus transpose.
It acts on all of \(M_2\), including the scalar line.

KTR8 is precisely the full infinite-cyclotomic irreducibility
hypothesis of Thorne's Theorem 6.1, with \(F=K\),
\(F^+=\mathbf Q\), and \(n=2\).
That theorem states \(H^1_f(\mathbf Q,A)=0\) for this
polarized adjoint.  Its scope has no ordinary hypothesis and
no exclusion of the prime \(3\).  The exact published
statement, its full \(M_n\) action, and its BLGGT conventions
were checked in ADK3--ADK6; here (KTR.14)--(KTR.17) and KTR8
verify every representation-dependent hypothesis afresh.
KTR1 brings the result from finite coefficients and the classical
period ring to every cocycle and every local primitive in the
original \texttt{Af.\allowbreak Is\allowbreak BK} definition.  Thus the conclusion
is actual coboundary vanishing, not merely a vanishing dimension
for a different coefficient realization.
% END KTR10

% BEGIN KTR11
\paragraph{KTR11. The remaining quadratic scalar line.}
\label{claim:KTR11}
The one-dimensional representation with character \(\eta\)
has \(H^1_f(\mathbf Q,\overline{\mathbf Q}_\ell(\eta))=0\).

\emph{Proof.}
By KTR1 work over a finite coefficient field \(E\).
Restriction to \(K\) changes the representation into the
trivial representation.  For a trivial representation the local
finite condition at a place above \(\ell\) is the unramified
condition.  Explicitly the Bloch--Kato fundamental exact sequence
used in BX7 gives \(\dim_E H^1_f(K_v,E)=1\): the de Rham
tangent quotient is zero and the invariant space has dimension
one.  The unramified group is already one-dimensional and is
contained in the finite group, so they agree.  Away from
\(\ell\) this equality is the definition.  A global cocycle
for the trivial representation is a continuous additive
homomorphism.  If it is finite everywhere it therefore kills
every inertia subgroup, and factors through the finite ideal
class group of \(K\).  There is no nonzero homomorphism from
that finite group to the additive group of a characteristic-zero
field.  Thus \(H^1_f(K,E)=0\).  Restriction from \(\mathbf Q\)
is injective by the finite-group averaging argument in KTR3,
which proves the assertion.  This rational vanishing itself
does not require the stronger assumption \(\ell\nmid h_K\).
% END KTR11

% BEGIN KTR12
\paragraph{KTR12. Explicit equivariant decomposition of the nine-dimensional adjoint.}
\label{claim:KTR12}
The full original \(\operatorname{ad}r\) is the direct sum
of \(A\), \(W\), \(W\eta\), and \(\eta\), through
the following concrete continuous linear maps.  For
\(X\in M_3\), \(Y\in M_2\), and \(v=(v_0,v_1)\), set
\[
 \begin{aligned}
 p_A X&=\begin{pmatrix}X_{00}&X_{02}\\X_{20}&X_{22}\end{pmatrix},
 &s_A Y&=\begin{pmatrix}Y_{00}&0&Y_{01}\\0&0&0\\Y_{10}&0&Y_{11}\end{pmatrix},\\
 p_tX&=(X_{01}+tX_{12},\ X_{21}-tX_{10}),
 &s_tv&=\begin{pmatrix}0&v_0&0\\-t v_1&0&t v_0\\0&v_1&0\end{pmatrix},
 \quad t\in\{1,-1\},\\
 p_CX&=X_{11},&s_Cz&=zE_{11}.
 \end{aligned}
 \tag{KTR.18}
\]
Here \(p_1,s_1\) intertwine with \(W\), and
\(p_{-1},s_{-1}\) with \(W\eta\).  The identity is
\[
       X=s_Ap_AX+\tfrac12s_1p_1X+
                       \tfrac12s_{-1}p_{-1}X+s_Cp_CX.
 \tag{KTR.19}
\]

\emph{Proof.}
On \(H\), write the corner matrix as
\(B=\left(\begin{smallmatrix}a&b\\c&d\end{smallmatrix}\right)\)
with \(ad-bc=1\), and its three-dimensional embedding as
\(\widetilde B\).  The original adjoint action there is
\(X\mapsto\widetilde B X\widetilde B^{-1}\).
The four entries of \(p_A\) undergo \(Y\mapsto BYB^{-1}\).
For the off-diagonal entries, direct multiplication gives
\[
 \begin{split}
 p_t(\widetilde B X\widetilde B^{-1})_0
       &=a(X_{01}+tX_{12})+b(X_{21}-tX_{10}),\\
 p_t(\widetilde B X\widetilde B^{-1})_1
       &=c(X_{01}+tX_{12})+d(X_{21}-tX_{10}).
 \end{split}
\]
The same multiplication gives
\(\widetilde B s_t(v)\widetilde B^{-1}=s_t(Bv)\)
and fixes \(X_{11}\).

At \(c\), the source adjoint is
\(\Theta(X)=-J_3X^tJ_3\), with \(J_3\) the reverse
identity matrix; this follows from \(R(c)=1\) and
\(\eta(c)=-1\).  Its corner action is
\(-J_2(p_AX)^tJ_2\), and its middle entry is \(-X_{11}\).
For the other two maps one obtains
\[
       p_t\Theta(X)=\operatorname{diag}(-t,t)\,p_tX.
 \tag{KTR.20}
\]
For \(t=1\) this is \(D_2=W(c)\); for \(t=-1\) it
is \(-D_2=(W\eta)(c)\).  Substitution in \(s_t\)
gives the identical intertwining law in the other direction.
The groups \(H\) and \(\{c\}\) generate \(G\), so
all these maps intertwine the full disconnected action.
Finally (KTR.19) is an entrywise identity: the sum of the
\(t=1\) and \(t=-1\) terms cancels each unwanted cross
entry and retains the desired one.  Their pairwise projections
have zero cross compositions and \(p_ts_t=2\,1\), while
\(p_As_A=p_Cs_C=1\).  This proves the direct-sum assertion.
All maps have fixed rational coefficients and are continuous
for the original product \(\ell\)-adic topology.
% END KTR12

% BEGIN KTR13
\paragraph{KTR13. The full literal RigII producer.}
\label{claim:KTR13}
For every original datum and actual point of \texttt{Rig\allowbreak II},
if \(\texttt{Type\allowbreak II}(\eta,r)\) and
\[
 \forall W,\quad D.\texttt{Is\allowbreak Ext}(r,W)
                 \ \Longrightarrow\ D.\texttt{Is\allowbreak Modular}(\iota,W),
 \tag{KTR.21}
\]
then the source's
\(\texttt{Bloch\allowbreak Kato\allowbreak Selmer\allowbreak Vanishes}(\ell,\operatorname{ad}r)\)
holds.

\emph{Proof.}
KTR4--KTR6 construct the actual \(W\) from the point and
verify \texttt{Is\allowbreak Ext}; (KTR.21) gives its modular newform for
the original fixed \(\iota\).  KTR2--KTR3 give Selmer
vanishing for \(W\) and \(W\eta\).  KTR8--KTR10 give
vanishing for the corner adjoint \(A\); KTR11 gives it for
the remaining scalar character \(\eta\).

Let \(a:G\to M_3(\overline{\mathbf Q}_\ell)\) be any
continuous cocycle satisfying the original local conditions.
Apply \(p_A,p_1,p_{-1},p_C\) from KTR12.  KTR1 makes all
four images continuous Selmer cocycles for their stated actions.
Choose actual primitives \(Y,v_+,v_-,z\) supplied by the
four vanishings.  Equivariance and (KTR.19) show that
\[
 a(g)=(\operatorname{ad}r)(g)Q-Q,
 \qquad
 Q=s_A Y+\tfrac12s_1v_++\tfrac12s_{-1}v_-+s_Cz.
 \tag{KTR.22}
\]
This is the literal required global coboundary.

The argument retains the supplied \(\Omega,K,\ell,\mathcal O\),
the complete \(\texttt{GD}\) conditions, \(N\ne0\),
\(e\), \(\psi\), the universal local ring and its actual
point map, the marked \(g_1\), and the normalization at the
specified \(c\).  It uses only the subset of those hypotheses
identified in KTR4--KTR8, so it imposes no new level or weight
condition.  It retains the full infinite cyclotomic kernel,
the actual representation, every local \texttt{Is\allowbreak BK} witness,
and all nine adjoint coordinates.  The original odd-prime
scope includes \(3\); no division by \(\ell\) or omitted
small-prime case is introduced.  The ordinary-coefficient
automorphy producer and the remaining \texttt{One\allowbreak Char} length
argument are distinct tasks and are not conclusions of this
claim.
% END KTR13

\subsection{Source and foundation boundary}

The original endpoint statements are
\path{ModularityArgument/KatoQuadraticCharacters/KatoTwoOfUseN.lean},
lines 11--53.  The local ideal and its ordered products are in
\path{PadicHodgeTheory/IsInertiaInfiniteCyclotomic.lean},
lines 87--138.  The actual point and extension predicates are
\path{UnitaryThreeGroups/RhoKApply.lean}, lines 16--93.
The finite-dimensional source route appears in
\path{ModularityArgument/KatoQuadraticCharacters/ReducibleThreeDim.lean};
we have written its actual projection, section, and primitive
calculations in KTR12--KTR13.  KTR4--KTR8 give a smaller local
argument sufficient for this endpoint, without claiming the
stronger full local crystalline-component theorem.

The new published applications are Kato Theorem 14.2(1), with
its \S14.1 local conditions and the actual critical twist, and
quadratic \(\mathrm{GL}_2\) base change in Langlands's
Lemma 11.3(b) and Proposition 11.4.  BLGGT Theorem 2.1.1 and
Thorne Theorem 6.1 are the previously matched foundational
theorems; KTR9--KTR10 check this new representation's exact
hypotheses.  The finite-coefficient and tied-period comparisons
are inherited from ADK1--ADK5, BX1, BX7 and DX1; the actual
newform adelization is RNC2 and the dense Frobenius comparison
is CAP6.  All these dependencies and every exact source range
are recorded separately from the proof in the component
provenance.  No original Lean proof is used as the human proof.


\section{The complete original Kato input record}
\label{sec:kato-input-assembly-v47}

This short chapter performs the exact remaining record assembly.
The supplied statements are universal mathematical theorems, each
with its original antecedents.  They can therefore be used on the
same subsequently chosen arithmetic objects.  The result is the
literal rank-three input \texttt{Inp\allowbreak Z}; it does not yet prove its
one-character length consumer or that consumer's automorphy premises.
All source references are at commit
\texttt{7cbf7a0c32fca1badf05931879a37c3f8a2b03b4}.

% BEGIN KIZ1
\paragraph{KIZ1. The five fields of the original cohomological input.}
\label{claim:KIZ1}
There is a full original \(\mathrm{CohInp}\), with the single
constant \(C_K=31\) in its \(\mathrm{lemK}\) field.  This constant
is chosen before every coefficient prime, number field, algebraic
closure, Galois datum and positive integer weight.

\begin{proof}
Let \(\mathcal C\) be the cohomology record constructed in BX8.
The definition of \(\mathrm{CohomologyHyp}\) has the six fields
\(\mathrm{loc},\mathrm{pt},\mathrm{gf},\mathrm{ge},
\mathrm{uo},\mathrm{gw}\).  Its first four have exactly the
types of the first four fields of \(\mathrm{CohInp}\).
Keep these four proofs, without changing any of their data.
Here is their mathematical content and the required scope.

The local field \(\mathcal C.\mathrm{loc}\) contains all seven
original assertions.  Its finite-local-field finiteness, Tate
duality and Euler assertions are \(\mathrm{lfK},\mathrm{ltK},
\mathrm{leK}\); the same three assertions on the original
decomposition groups are \(\mathrm{lf},\mathrm{lt},\mathrm{le}\).
They use the original cyclotomic character and arbitrary finite
coefficient modules, including nonfree modules.  The last field
\(\mathrm{bk}\) is the entire Bloch--Kato local statement on
every original tied period presentation.  In particular, it includes
both integral annihilator directions on the actual cocycles and
their original local cup products, rather than just a rational
dimension equality.  These are precisely the seven fields installed
in BX8, so no new local theorem is assumed here.

For \(\mathrm{pt}\), retain the normalized invariant family
already used by \(\mathcal C\).  Its choice precedes the finite
set of places and the finite global module, as required by the
original \(\mathrm{PoitouTateNormalised}\).  The local restriction
and cup maps, uniformizer normalization and global exact sequence
are unchanged.  For \(\mathrm{gf}\) and \(\mathrm{ge}\), retain
the original global finite-cohomology and Euler assertions.  As
checked in DX1 and installed in BX8, they apply to every prime
allowed by those definitions.  The Euler expression at the prime
two is the original truncated degree-\(0,1,2\) expression; the
separate high-degree vanishing assertion keeps its original
odd-prime antecedent.  Thus extracting these four fields introduces
no restriction inherited from a later rank-three application.

It remains to give \(\exists C_K\in\mathbb R\,
\mathrm{LemmaK}(C_K)\).  KB8 supplies \(\mathrm{LemmaK}(31)\).
In ordinary notation, for every original prime \(\ell\), totally
complex number field \(L\), algebraic closure \(\Omega/L\),
and original datum \(D\), assume that \(L\) contains a primitive
\(\ell\)-th root of unity when \(\ell\ne2\), or a primitive
fourth root when \(\ell=2\).  Then for every \(e\ge1\), the
actual compact continuous cohomology group
\(H^2(D,\mathbb Z_\ell(e+1))\) is finite and
\[
 \log\bigl|H^2(D,\mathbb Z_\ell(e+1))\bigr|
 \le
 \exp\!\left((2e[L:\mathbb Q])^{31}
                    \log(2|\operatorname{disc}L|)\right).
 \tag{KIZ.1}
\]
This is the literal \(\mathrm{kBound}\) and the literal
\(\mathrm{HasMuQ}\) guard in the definition of \(\mathrm{LemmaK}\).
The KB proof fixes the exponent before all these variables; in
particular it is not a constant chosen after \(L\) or \(\ell\).
Use the real number \(31\) and that theorem as the existential
witness.  The four retained fields and this witness are all five
fields of \(\mathrm{CohInp}\), proving the assertion.
\end{proof}
% END KIZ1

% BEGIN KIZ2
\paragraph{KIZ2. The exact rank-three input, on the same arithmetic objects.}
\label{claim:KIZ2}
The full original \(\mathrm{UpperTheory.V.Z.Kato.InpZ}\) is
supplied.  Its \(\mathrm{toCohInp}\) is the record of KIZ1,
and its six additional fields are
\[
\begin{array}{c|c|l}
\text{field}&\text{proved supplier}&\text{original conclusion}\\ \hline
\mathrm{lemM}&\mathrm{SP6}&\text{two-sided orbit and cotangent bound}\\
\mathrm{lemM'}&\mathrm{SP1}&\text{additive-injection length bound}\\
\mathrm{lemP}&\mathrm{FP8}&\text{universal weighted framed presentation}\\
\mathrm{rigII}&\mathrm{KTR13}&\text{full type-II adjoint rigidity}\\
\mathrm{adK}&\mathrm{ADK6}&\text{full polarized CM adjoint vanishing}\\
\mathrm{number}&\mathrm{RNC16}&\text{the rank-three congruence integer}.
\end{array}
\tag{KIZ.2}
\]
Every antecedent, representation, scalar map, embedding and
quantifier order in these fields is the original one.

\begin{proof}
We verify the complete assertions stored in (KIZ.2), since a
statement restricted to one specially chosen application would not
be a proof of these universally quantified fields.

For \(\mathrm{lemM}\), let the complete characteristic-zero DVR
\(O\), algebraically closed valued coefficient field, graded
\(O\)-algebra \(A\), augmentation \(x\), equations \(Q\)
and element \(a\) be arbitrary data in the original statement.
Retain exactly its normalized valuation, weighted-ring and
\(M1\)--\(M5\) hypotheses.  SP6 constructs a finite cover of
the original two-sided orbits, proves that the weight-\(-1\)
cotangent quotient has finite length, and proves, for every such
finite cover \(T\),
\[
 \ell_O C_{-1}(Q.F)
             \le 2\sum_{\xi\in T}v(\xi(a)).
 \tag{KIZ.3}
\]
The quotient is the one formed using the saturated augmentation
square and the supplied equations.  No choice of a new chart,
orbit equivalence or valuation is made in this assembly.  Thus SP6
is the complete \(\mathrm{LemmaM}\) field, including its
existence clause and its quantifier over every finite cover.

For \(\mathrm{lemM'}\), SP1 uses an arbitrary DVR with finite
residue field, an arbitrary abelian grading group, the supplied
augmentation and weight, and the original nonzero \(c\in O\).
The supplied map from the module to
\(\operatorname{Hom}_O(A,O/(c))\) is only required to be
additive and injective.  SP1 retains its actual weighted derivation
lifts and its vanishing on the supplied equations, and concludes
\(\ell_O M\le\ell_O C_j(F)\), including the possibility that
the right side is infinite.  Consequently it supplies the full
field without adding linearity of that additive map or dropping
the finite-residue-field condition.

For \(\mathrm{lemP}\), apply FP8 with all the originally
specified data still arbitrary: the odd residual prime, complete
Noetherian local coefficient ring with finite residue field,
compact totally disconnected global group, original form and
weights, residual and integral representations, local conditions,
base ring and local frames.  Its original hypotheses are retained:
the finite-character condition on open subgroups, both normalization
conditions, residual triviality on the prescribed local kernels,
the local-condition axioms and base points, finite framed
obstruction spaces, and the given weighted local data.
FP8 constructs the actual universal ring and representations over
this very base, preserves the old variables and relations, and
identifies the new variables as original matrix or frame entries.
It proves the exact weightwise numbers of generators and equations,
respectively
\[
 \begin{aligned}
 \#\{\text{new variables of weight }j\}
      &=\dim_k H^1_F(\operatorname{ad}_j),\\
 \#\{\text{new equations of weight }j\}
      &=\dim_k\operatorname{Ob}_F(\operatorname{ad}_j).
 \end{aligned}
 \tag{KIZ.4}
\]
These are the full \(\mathrm{IsUnivFramed}\),
\(\mathrm{IsRelPres}\) and entry-weight conclusions of
\(\mathrm{LemmaP}\), not only numerical upper bounds.
In particular the same original theorem remains available at
\(\ell=3\).  Its application hypotheses are not asserted for
an unspecified later arithmetic datum merely by storing the theorem.

For \(\mathrm{rigII}\), take any original
\(D:\mathrm{Dat}(\Omega,K,c,e,\psi)\), its local datum
\(U\), marked element \(g_1\), actual point representation
\(r\), and fixed \(\iota:\overline{\mathbb Q}_\ell\simeq
\mathbb C\).  Suppose the exact source \(\mathrm{TypeII}\)
condition and its universal implication
\[
 \forall W,\quad D.\mathrm{IsExt}(r,W)
                  \Longrightarrow D.\mathrm{IsModular}(\iota,W)
 \tag{KIZ.5}
\]
hold.  KTR13 constructs the original rank-two extension from this
\(r\), applies (KIZ.5) to that extension, proves the four required
Selmer vanishings, and recombines actual primitives through the
explicit \(4+2+2+1\) decomposition of its nine-dimensional
adjoint.  Its conclusion is exactly
\(\mathrm{BlochKatoSelmerVanishes}(\ell,\operatorname{adRep}(r))\).
All original \(\mathrm{GD}\), local-point and marked-element
guards remain in force.  Thus this field does not assert the
modularity implication (KIZ.5); it proves the original theorem
which consumes that implication.

For \(\mathrm{adK}\), ADK6 is universal over the original
imaginary quadratic \(K\), every rank \(n\ge2\), every
coefficient prime and closure, and the same fixed \(\iota\).
It uses the given continuous conjugate-self-dual automorphic
representation, its irreducibility on the full infinite cyclotomic
kernel, and the given \(\mathrm{IsAdOf}\) witness with multiplier
\(\varepsilon^{1-n}\).  It kills every source cocycle in the
full matrix adjoint, including its scalar part, with the original
disconnected conjugation action and local finite conditions.  No
trace-zero replacement, finite-cyclotomic substitute, residual-image
condition or exclusion of the prime three is introduced.

Finally fix the original imaginary quadratic field \(K\).
RNC16 chooses \(c,A\) at this stage, before \(e,N,p,w\)
and the root-of-unity orbit.  For every \(e,N\ge1\), prime
\(p\ge7\) away from \(N\operatorname{disc}K\), place
\(w\mid p\), and original finite indexed orbit \(\mathcal O\),
it constructs one \(NN\in\mathbb Z\setminus\{0\}\) with
\[
 \log|NN|\le c\,|\mathcal O|\,[N(e+1)]^A\log p.
 \tag{KIZ.6}
\]
This integer is chosen before every later odd \(\ell\ne p\),
coefficient closure, scalar tower, embedding \(\iota\),
prolongation and Frobenius.  For each subsequently supplied family
of pairwise nonisomorphic admissible representations, it gives both
nonvanishing of every trace gap and the exact product inequality
for \(|NN|_\ell\).  Its trace is evaluated at geometric
\(\sigma^{-1}\) for the original arithmetic Frobenius \(\sigma\).
The predicates are precisely rank-three \(\mathrm{Adm3}(\iota,e,N)\)
and the rank-two-plus-middle-line
\(\mathrm{Adm2}(\iota,e,N|\operatorname{disc}K|)\); the
rank-four congruence statement is not used.

All six suppliers are therefore proofs of the complete predicates
appearing in the definition of \(\mathrm{InpZ}\).  Insert them
as (KIZ.2), with the exact \(\mathrm{CohInp}\) of KIZ1.
These are its eleven fields and there are no further compatibility
fields.  Their quantified inputs are not replaced by witnesses from
other applications: when a later proof fixes \(\Omega,K,\ell,c,
e,N,\psi\), one specializes each theorem to those same objects.
The period and normalized duality data are the retained ones of
KIZ1; the original Frobenius, cyclotomic, coefficient and adjoint
conventions of KTR13, ADK6 and RNC16 agree as just verified.
This constructs the required record.

It closes the exact record-installation task.  The separate
\(\mathrm{OneChar}\) argument must still construct and compare
its arithmetic rings and points, and its automorphy supplier remains
separate.  Neither is a field of \(\mathrm{InpZ}\), so neither
has been assumed or inferred in this construction.
\end{proof}
% END KIZ2

\paragraph{Source and foundation boundary.}
The exact records are \path{UnitaryThreeGroups/KatoUse.lean},
lines 19--24, and
\path{ModularityArgument/KatoQuadraticCharacters/KatoTwoOfUseN.lean},
lines 55--61.  The separately recorded source slices give every
stored predicate in full.  The numbered human proofs cited above
retain their previously identified classical foundations.  This
assembly uses no additional published theorem, no new automorphy
assertion, and no unreconstructed record field.


\section{Finite dimensionality of the original automorphic row spaces}
\label{sec:row-finiteness-v47}

This chapter proves the full original \texttt{Rows\allowbreak Fin\allowbreak Dim} assertion.
The test function need not be an idempotent, and its archimedean unit
vector need not be smooth or compact-finite.  There is no gap, rank,
weight, cuspidality or ordinary hypothesis beyond the literal statement.
We use classical Harish--Chandra finiteness in its precise published
scope; the transfer to the original quotient, level and operators is
proved here.  Source commit:
\texttt{7cbf7a0c32fca1badf05931879a37c3f8a2b03b4}.

Fix an original \(I:\texttt{Aut\allowbreak Inst}\).  Its group is the unitary
group of the supplied nonsingular Hermitian matrix over its CM
extension \(E/F\).  Write
\[
 G_A=G_\infty\times G_f,
 \quad \Gamma=G(F),
 \quad \mathcal H=L^2(\Gamma\backslash G_A,\nu).
 \tag{RFD.1}
\]
Here \(\nu\) is the quotient measure from the supplied Haar measure
\(\mu_f\), the original infinity Haar measure, and the supplied
fundamental domain.  No probability normalization of \(\nu\) is made.
The operator \(R(g)\) is the original right regular representation.

% BEGIN RFD1
\paragraph{RFD1. The exact classical finite-multiplicity input.}
\label{claim:RFD1}
Let \(G_\infty\) be a product of real unitary groups and let
\(\Delta\) be an arithmetic subgroup of a rational reductive group
whose real group is \(G_\infty\).  For every irreducible continuous
unitary representation \(\pi\), the vector space of bounded
intertwiners
\[
 \operatorname{Hom}^{\mathrm{bdd}}_{G_\infty}
       (\pi,L^2(\Delta\backslash G_\infty))
 \tag{RFD.2}
\]
is finite dimensional, with any positive Haar normalization.

\emph{Published input and hypothesis match.}
The inputs are Harish--Chandra admissibility of irreducible unitary
representations and finiteness of automorphic forms of a fixed finite
compact type and fixed finite-codimensional central ideal.  A primary
exposition is A.~Borel, \emph{Automorphic forms on reductive groups},
\S6.5, \S8.4 and \S10, printed pages 17, 23 and 27--28; in the latter section
the subparagraphs are printed as 8.5--8.6.  The same text explicitly
identifies compact-finite vectors in an irreducible \(L^2\) submodule
as automorphic forms, and derives finite multiplicity.  It is available at
\url{https://homepages.math.uic.edu/~rtakloo/papers/borel/borel3.pdf}.
These are named classical foundations, not new proofs of the
Harish--Chandra theorems.

For clarity, the deduction that will be used is as follows.  Fix a
nonzero smooth compact-finite vector \(v\) in \(\pi\); its existence
and its infinitesimal character follow from admissibility.  Let
\(\xi\) be a finite sum of compact-type idempotents fixing \(v\),
and let \(J\) be the kernel of that infinitesimal character.
For every bounded intertwiner \(A\), the vector \(A(v)\) has the
same compact-type and central-annihilator restrictions.  Its
automorphic representative is therefore in the one finite-dimensional
space \(\mathcal A(\Delta,J,\xi)\).  The map \(A\mapsto A(v)\)
is injective: the closed span of the translates of \(v\) is all of
\(\pi\), and a bounded map vanishing there is zero.  This proves
(RFD.2).  If \(A\ne0\), Schur's lemma gives \(A^*A=c\,1\),
\(c>0\), so its image is an actual closed irreducible submodule of
\(L^2\); the cited assertion about \(L^2\) vectors applies.
This deduction does not assume that the full regular representation
is discrete or that its residual spectrum vanishes.

The groups in question are connected reductive real matrix groups
with compact center, with arithmetic \(\Delta\).  These are precisely
the classical hypotheses.  Equivalently one can restrict scalars to
\(\mathbf Q\); there is no split central factor needing a separate
central-character quotient.  Multiplying Haar measure by a positive
constant gives an isometric rescaling of the Hilbert space and leaves
the dimension in (RFD.2) unchanged.  The trivial rank-zero group is
handled directly and has a one-dimensional regular space.
% END RFD1

% BEGIN RFD2
\paragraph{RFD2. One compact-open level for the complete original test function.}
\label{claim:RFD2}
Suppose \(e:G_f\to\mathbf C\) is locally constant and
\(e(x)\ne0\Rightarrow x\in I.Kset\).  There is a compact open
subgroup \(K_0\subset I.Kset\) such that
\[
             e(kx)=e(x)=e(xk)\qquad(k\in K_0, x\in G_f).
 \tag{RFD.3}
\]
In particular \(e\) is bounded and integrable.

\emph{Proof.}
The stabilizer \(I.Kset\) of the original full integral lattice is
compact open.  At each finite place choose an integral lattice basis;
its stabilizer is a closed subgroup of the compact matrix group,
and a sufficiently small principal congruence subgroup lies in it.
The restricted product has the same property.  This is the compact
open subgroup used in the original adelic topology, as in FV4.

Each point of \(I.Kset\) has an open right coset on which \(e\)
is constant.  Choose a finite such cover by compactness.  Intersect
the finitely many open subgroups, and then take its core in
\(I.Kset\).  The core is still open: an open subgroup of a compact
group has finite index, so there are only finitely many conjugates.
Call the resulting open normal subgroup \(K_0\).  Every chosen
coset is a union of \(K_0\)-cosets; hence \(e\) is constant on
right \(K_0\)-cosets in \(I.Kset\).  Normality gives left
invariance there.  Outside \(I.Kset\), multiplication by an element
of \(K_0\) stays outside, and both sides are zero.  The finite
coset cover bounds the values of \(e\), and its support lies in
the compact set \(I.Kset\), proving integrability.  The level is
chosen from \(e\) alone, before choosing any archimedean vector.
% END RFD2

% BEGIN RFD3
\paragraph{RFD3. The actual fixed-level quotient has finite real multiplicities.}
\label{claim:RFD3}
For a compact open \(K_0\subset G_f\), put
\(\mathcal H^{K_0}=\{F:R(1,k)F=F\ (k\in K_0)\}\).
Every irreducible unitary representation of \(G_\infty\) has
finite bounded-intertwiner multiplicity in \(\mathcal H^{K_0}\).

\emph{Proof.}
The coordinate argument in FV1--FV4 applies to the literal
\(I.E/I.F,I.J,I.L\): the unitary equation defines a connected
reductive group, its geometric form is \(\mathrm{GL}_n\),
and conjugation negates determinant characters.  Thus it has no
nontrivial rational character, and its real center is compact.
FV5--FV10 consequently apply to these same coordinates and give
finitely many representatives \(s\in G_f\) for
\(G(F)\backslash G_f/K_0\).  For each put
\[
 \Delta_s=\{\gamma_\infty:\gamma\in G(F),\
                         \gamma_f\in sK_0s^{-1}\}.
 \tag{RFD.4}
\]
The arithmeticity proof in FV10 uses the inverse image of
\(sK_0s^{-1}\), its finite-index intersection with an integral
stabilizer, and restriction of scalars.  It therefore proves
arithmeticity for these very subgroups; no strong approximation
or single-class assertion is being used.

A right \(K_0\)-invariant function is determined by the finitely
many functions \(F_s(g)=F(g,s)\).  They are left
\(\Delta_s\)-invariant.  Conversely these functions reconstruct
the original function using \(g_f=\gamma_fs k\); independence
of the representation follows exactly from (RFD.4).
Unfolding the original Haar measure gives the direct sum of the
\(L^2(\Delta_s\backslash G_\infty)\) norms with positive
constant factors.  These are the fiber volumes of \(K_0\) in
the original product Haar convention.  In particular the map is
a Hilbert equivalence after multiplying each component by the
square root of its factor.  Null-set identifications and the
choice of the original fundamental domain are legitimate by
RQC2--RQC3.  The equivalence commutes with the original right
\(G_\infty\)-action.

Apply RFD1 to each of the finitely many arithmetic components.
The space of intertwiners into their finite Hilbert direct sum
is the direct sum of the corresponding finite-dimensional spaces.
Transporting back proves the assertion with the original measure
and the original regular action.  No change of those data is
made in the conclusion.
% END RFD3

% BEGIN RFD4
\paragraph{RFD4. The original infinity tensor representation and coefficient.}
\label{claim:RFD4}
Let \((H_v,\sigma_v,d_v,w_v)\) be the complete original infinity
data in \texttt{Rows\allowbreak Fin\allowbreak Dim}.  On their Hilbert tensor product set
\[
 H_\sigma=\widehat\bigotimes_{v\mid\infty}H_v,
 \quad \sigma=\widehat\bigotimes_v\sigma_v,
 \quad w=\bigotimes_vw_v,
 \quad d=\prod_vd_v.
 \tag{RFD.5}
\]
This is an irreducible strongly continuous unitary representation,
\(\|w\|=1\), and, with inner products linear in the second variable,
\[
 m(g)=\langle w,\sigma(g)w\rangle,
 \qquad \phi(g)=d\,\overline{m(g)}
       =\texttt{ph\allowbreak Inf}(I,\sigma_v,w_v,d_v)(g).
 \tag{RFD.6}
\]
The function \(\phi\) is integrable for the original infinity measure.

\emph{Proof.}
There are finitely many real places.  Start with the algebraic
tensor product, give elementary tensors the product inner product,
and complete.  The tensor operators are isometries and extend to
that completion.  Continuity follows on elementary tensors from
the finite product, and then on all vectors by isometric
approximation.  The elementary tensors are dense.  Each factor
is irreducible by the original hypothesis; its bounded commutant
is scalar.  The bicommutant theorem therefore makes the strong
closure of its operator algebra all bounded operators.  Acting
one factor at a time shows that the commutant of the product
representation is scalar, so the product representation is
irreducible.  The empty tensor convention is \(\mathbf C\).

The inner-product identity on tensors gives (RFD.6), the precise
product used by the source.  Tonelli applied to absolute values
gives
\(\int|\phi|=\prod_v\int|\phi_v|<\infty\).
The original four-vector formal-degree identities tensor by
Fubini.  No compact-finiteness of the chosen \(w_v\) is used.
% END RFD4

% BEGIN RFD5
\paragraph{RFD5. Schur compression on every unitary Hilbert representation.}
\label{claim:RFD5}
Let \(U\) be any strongly continuous unitary representation of
\(G_\infty\), and let \(T=U(\phi)\), with \(\phi\) from RFD4.
Then
\[
 T^*=T,\qquad T^2=T,\qquad T U(g)T=m(g)T.
 \tag{RFD.7}
\]
These conclusions also hold for an operator given by the original
weak \texttt{Is\allowbreak Op\allowbreak Int} identity.

\emph{Proof.}
The original four-vector Schur identity, applied one place at a
time and then multiplied, gives the coefficient convolution laws
\[
 \phi(g^{-1})=\overline{\phi(g)},\qquad
 \int\phi(h)\phi(x^{-1}h^{-1}y)\,dh=m(x)\phi(y).
 \tag{RFD.8}
\]
To check the second equation, substitute
\(\phi(h)=d\overline{\langle w,\sigma(h)w\rangle}\), move
\(\sigma(x)\) and \(\sigma(y)\) across the inner product by
unitarity, and use the four-vector identity; its factor \(1/d\)
cancels one of the two factors \(d\).  The remaining coefficient
is \(m(x)\), with the displayed order of \(x^{-1}h^{-1}y\).
Putting \(x=1\) gives \(\phi*\phi=\phi\).

Since \(\phi\in L^1\), its integrated operator exists and has
norm at most \(\|\phi\|_1\).  The bounds for products of two
integrands are multiples of \(|\phi(h)\phi(k)|\), so Fubini
justifies operator multiplication through weak matrix coefficients.
Inversion and (RFD.8) prove all of (RFD.7).  A weakly specified
operator agrees with this integrated operator because its inner
product with every vector agrees.  This is also the direct
Schur-compression computation underlying ARO16; here its stated
formal-degree hypotheses, rather than \texttt{Gap\allowbreak Hyp}, are used.
% END RFD5

% BEGIN RFD6
\paragraph{RFD6. Every projector vector gives a bounded intertwiner.}
\label{claim:RFD6}
For every \(v\in\operatorname{ran}T\), there is a unique bounded
intertwiner \(A_v:H_\sigma\to H_U\) satisfying
\[
 A_v(\sigma(g)w)=U(g)v,
 \qquad A_v(w)=v,
 \qquad \|A_v\|=\|v\|.
 \tag{RFD.9}
\]
The map \(v\mapsto A_v\) is linear and injective.

\emph{Proof.}
For \(v,v'\in\operatorname{ran}T\), self-adjointness and
compression give
\[
 \langle U(g)v,U(h)v'\rangle
   =m(g^{-1}h)\langle v,v'\rangle.
 \tag{RFD.10}
\]
The same coefficient \(m(g^{-1}h)\) is
\(\langle\sigma(g)w,\sigma(h)w\rangle\).
For a finite linear combination, (RFD.10) therefore shows
\[
 \left\|\sum_j a_jU(g_j)v\right\|^2
 =\|v\|^2\left\|\sum_j a_j\sigma(g_j)w\right\|^2.
 \tag{RFD.11}
\]
Thus the prescribed map is well-defined on the algebraic orbit
span and is bounded.  That span is dense by irreducibility and
\(w\ne0\); completeness extends it uniquely.  The formula on
the orbit span proves intertwining, linearity in \(v\), and
uniqueness.  Evaluation at \(w\) proves injectivity and the
operator-norm equality.  This gives a direct argument even when
the ambient regular representation has continuous spectrum.
% END RFD6

% BEGIN RFD7
\paragraph{RFD7. The fixed-level archimedean projector has finite rank.}
\label{claim:RFD7}
For every compact open \(K_0\subset G_f\), the restriction of the
original \(T=R(\phi)\) to \(\mathcal H^{K_0}\) has finite-dimensional
range.

\emph{Proof.}
Finite and infinite right translations commute, so their weak
integrals commute, and \(T\) preserves \(\mathcal H^{K_0}\).
Apply RFD6 to the infinity representation on this closed Hilbert
subspace.  It injects \(\operatorname{ran}(T|\mathcal H^{K_0})\)
linearly into
\[
 \operatorname{Hom}^{\mathrm{bdd}}_{G_\infty}
       (H_\sigma,\mathcal H^{K_0}),
\]
which is finite dimensional by RFD3--RFD4.  This proves the
assertion, with the actual original \(w\); the auxiliary
compact-finite vector in RFD1 was used only to bound the space
of intertwiners and does not replace \(w\).
% END RFD7

% BEGIN RFD8
\paragraph{RFD8. Full original RowsFinDim, including arbitrary test functions.}
\label{claim:RFD8}
For every datum in the full original \texttt{Rows\allowbreak Fin\allowbreak Dim} statement,
\[
 \dim_{\mathbf C}\operatorname{ran}(R_eT)<\infty.
 \tag{RFD.12}
\]
Consequently the literal assertion
\(\texttt{Formalism.\allowbreak L2.\allowbreak SG.\allowbreak Rows\allowbreak Def.\allowbreak Rows\allowbreak Fin\allowbreak Dim}\) holds at the
classical foundations recorded in RFD1 and FV5--FV10.

\emph{Proof.}
Choose \(K_0\) by RFD2 and let
\(P_0=\mu_f(K_0)^{-1}\int_{K_0}R(1,k)\,dk\).
It is the orthogonal projection onto \(\mathcal H^{K_0}\).
Both sides of (RFD.3) and Haar change of variables give
\[
                 P_0R_e=R_e=R_eP_0.
 \tag{RFD.13}
\]
The finite and infinity actions commute.  Fubini is legitimate
because \(e,\phi\in L^1\); hence \(R_eT=TR_e\).
For any \(F\in\mathcal H\), the vector \(R_eF\) is in
\(\mathcal H^{K_0}\), so
\[
 R_eTF=T(R_eF)\in\operatorname{ran}(T|\mathcal H^{K_0}).
\]
RFD7 proves this containing space finite dimensional.  The
original \texttt{Is\allowbreak RConv\allowbreak F} and \texttt{Is\allowbreak Op\allowbreak Int} identities
uniquely identify \(R_e,T\), so this is precisely the original
\texttt{HH} range on the supplied \(L^2\) space.  Neither
\(e*e=e\) nor \(e^*=e\) has been used; neither is a hypothesis
of \texttt{Rows\allowbreak Fin\allowbreak Dim}.  All coefficient-field, signature, Haar
and fundamental-domain quantifiers remain those of the source.
% END RFD8

% BEGIN RFD9
\paragraph{RFD9. The finite-dimensionality input is removed from the row constructions.}
\label{claim:RFD9}
On every supplied original adelic \(\texttt{SG.\allowbreak Hida\allowbreak Setting}\),
the finite-multiplicity hypothesis retained in ARO18--ARO30 is
now supplied by RFD8.  Their actual row construction, including
\texttt{Rows\allowbreak Add}, \texttt{Rows\allowbreak Good}, \texttt{Rows\allowbreak Lab} and one
common \texttt{All\allowbreak Rows}, consequently requires only their other
stated inputs.

\emph{Proof.}
The specialization in ARO18 supplies the original infinity
formal-degree and integrability hypotheses from the recorded
\texttt{Gap\allowbreak Hyp}, and its event function is locally constant
with support in \(I.Kset\).  RFD8 therefore proves the exact
finite-dimensionality assertion for that event.  The original
evaluation bijection in ARO17 transfers it to the bounded
intertwiner space.  ARO19--ARO24 then construct the rows with
the same measures, projectors, labels and ties.  ARO25--ARO30
use these same rows and do not require a further finiteness
input.  Their original \texttt{Aut\allowbreak Tie}, \texttt{Row\allowbreak Dict\allowbreak Tie},
\texttt{Lab\allowbreak Tie}, arithmetic counting and common-family premises
are unchanged.  Thus this removes \texttt{Rows\allowbreak Fin\allowbreak Dim} as a
producer task without closing any of those distinct premises.
% END RFD9


\section{Actual lattice and adelic blocks for the original global setting}
\label{sec:adelic-setting-v47}

This chapter constructs the original setting blocks, rather than assuming
their existence as part of a completed \texttt{SG.\allowbreak Hida\allowbreak Setting}.
Throughout, \(D\) is the original \(D7\) record, \(E/F\) its CM
extension, \(V=\mathrm{VS}(D)\), \(\Lambda=V.L_0\),
\(s=D.R.rS\), \(k=D.R.k\), and \(n=2s+k+1\).
The Hermitian form is conjugate-linear in its first variable.
We retain the original \(\mathrm{Forms0}(D,D.u_0)\) precisely where
the lattice producer requires it, and the original \(\mathrm{GapHyp}\)
precisely where the discrete-series bundle requires it.
Neither condition is replaced by a new existence hypothesis.
The source is pinned at
\texttt{7cbf7a0c32fca1badf05931879a37c3f8a2b03b4}.

% BEGIN ASG1
\paragraph{ASG1. The actual unitary-group instance on the original vector space.}
\label{claim:ASG1}
Every original \(D\) has a \(\mathrm{UnitaryGroupInstance}(D)\)
with coefficient fields exactly \(E,F\), and with the original lattice
and form in the original frame coordinates.
\begin{proof}
Let \(b_i\) be \(V.\mathrm{frame}\), put
\(J_{ij}=h(b_i,b_j)\), and let
\(c:E^n\longrightarrow V\) send \(x\) to \(\sum_i x_i b_i\).
Set \(L=c^{-1}\Lambda\).  Expanding both variables gives
\[
                 \overline{x}^{T}Jy=h(c(x),c(y)).
\tag{ASG.1}
\]
Thus \(J\) is Hermitian.  If \(Jx=0\), (ASG.1) and
nondegeneracy give \(c(x)=0\), hence \(x=0\); consequently
\(\det J\ne0\).  Finite generation and full span of \(L\)
follow by transporting those properties of \(\Lambda\) through
\(c\).  The original unimodularity assertion becomes exactly
\[
 x\in L\quad\Longleftrightarrow\quad
       h_J(y,x)\in\mathcal O_E\text{ for every }y\in L.
\tag{ASG.2}
\]

Use the restricted-product topology on the finite adelic unitary group
in these coordinates.  Local integral stabilizers are compact open:
after choosing local lattice bases they are closed in integral general
linear groups and contain a congruence neighbourhood of the identity.
The resulting congruence subgroups give the exact topology fields of
\(\mathrm{AutInst}\), as in FV4.

The CM pairs of embeddings are in bijection with real places of \(F\).
Use the original chosen member \(\mathrm{embS}(D,\tau)\) of each
pair.  Its restriction is the corresponding real embedding.  The
original \(D.sp.Y\_\allowbreak sig\) field gives signature
\((k+s,s+1)\).  Hermitian diagonalization therefore gives a matrix
\(P_v\) with
\(P_v^*\phi_v(J)P_v=J_{\epsilon_v}\).
The original Boolean sequence \(\mathrm{epsS}\) has exactly
\(k+s\) positive and \(s+1\) negative entries; permuting a
diagonalizing basis puts the signs in that literal order.
This supplies \(\epsilon,\phi,P\) and their identities.

Take the coefficient isomorphism \(\iota_E\) to be the identity,
the basis to be \((b_i)\), and the place bijection to be the inverse
of the original \(\mathrm{tauPl}\).  Complex conjugation agrees
with the CM involution, and its fixed field is \(F\).
Equation (ASG.1), the definition of \(L\), the rank identity in
\(D\), and the preceding sign count verify every remaining field
of \(\mathrm{UnitaryGroupInstance}\).  No lattice is replaced
by a principal one.
\end{proof}
% END ASG1

% BEGIN ASG2
\paragraph{ASG2. The original trace-one element and prescribed fractional ideals.}
\label{claim:ASG2}
There is \(t\in\mathcal O_E\) with \(t+\bar t=1\).
Every prescribed ideal
\(\mathfrak a_i=\mathrm{aaD}(D,i)\) is nonzero and invariant
under conjugation.  At every prime above
\(g=D.host.g(0)D.host.g(1)\), both \(\mathfrak a_i\) and its
inverse are the unit local ideal.
\begin{proof}
The original \(\mathrm{cmUnram}\) field gives, at every maximal
ideal of \(\mathcal O_E\), an integral \(a\) for which
\(a-\bar a\) is outside that ideal.  At a ramified quadratic
place, conjugation would lie in inertia and fix every residue, which
contradicts this field.  Thus \(E/F\) is unramified at finite
places.  The trace on the corresponding integral local extension is
surjective: in the split case it is addition on two copies of the local
ring; in the unramified field case its reduction is the surjective
finite-field trace, and a lift of an element with nonzero trace has
unit trace.  The integral trace image is an ideal of \(\mathcal O_F\)
which is the unit ideal at every maximal ideal, hence is the unit ideal.
Choose \(t\) of trace one.  This argument includes residue
characteristic two.

The ideals \(\mathfrak a_i\) are the specified integral powers,
possibly negative, of the different of \(F/\mathbb Q\) extended
to \(E\).  The different is nonzero and comes from the fixed
field, so it and all its powers are conjugation invariant.
Each of the two original rational primes splits completely in \(E\):
the source counts \([E:\mathbb Q]\) distinct primes above it,
and the identity \(\sum e_w f_w=[E:\mathbb Q]\) forces every
ramification and residue degree to be one.  The primes are therefore
unramified in \(F\), so its different is a unit at all these
places.  This proves the final assertion without excluding either
original prime.
\end{proof}
% END ASG2

% BEGIN ASG3
\paragraph{ASG3. A split prime ideal in the required ideal class.}
\label{claim:ASG3}
For every nonzero fractional ideal \(\mathfrak c\) of a CM
field \(E\), and every nonzero \(N\in\mathcal O_E\), there
are an integral prime ideal \(\mathfrak p\) and \(\alpha\ne0\)
such that
\[
 \mathfrak p+\overline{\mathfrak p}=\mathcal O_E,
 \qquad \mathfrak p+(N)=\mathcal O_E,
 \qquad \mathfrak p\mathfrak c=(\alpha).
\tag{ASG.3}
\]
\emph{Classical input and proof of the exact application.}
Use the Hilbert class field and Chebotarev's theorem in their usual
unconditional forms.  A primary exposition is J.~S.~Milne,
\emph{Class Field Theory}, version 4.03, Chapter V, \S3
(especially Example 3.9), and Chapter VIII, Theorem 7.4:
\url{https://www.jmilne.org/math/CourseNotes/CFT.pdf}.
The two inputs say that the ideal-class Artin map identifies the
class group with the unramified abelian Galois group, and that a
conjugacy class in a finite Galois group occurs at infinitely many
unramified primes.  They are named classical foundations.

Let \(H\) be the Hilbert class field of \(E\), and let
\(F=E^{\bar{\ }}\).  Every \(F\)-automorphism of an algebraic
closure sends \(H\) to the maximal unramified abelian extension
of its image of \(E\), which is again \(E\).  Hence \(H/F\)
is Galois.  Let \(\sigma\in\operatorname{Gal}(H/E)\) be the
Artin image of \([\mathfrak c]^{-1}\).  Apply Chebotarev to
its conjugacy class in \(\operatorname{Gal}(H/F)\), avoiding
the finitely many primes below \((N)\).  The restriction of every
element in this class to \(E\) is the identity; the selected prime
of \(F\) therefore splits into two distinct primes in \(E\).
Choose a prime of \(H\) whose Frobenius is \(\sigma\) itself
(conjugate the chosen prime if necessary), and let \(\mathfrak p\)
be its restriction to \(E\).  Its relative residue degree over
\(F\) is one, so its Artin symbol in \(H/E\) is also
\(\sigma\).  Thus \([\mathfrak p]=[\mathfrak c]^{-1}\).
The two distinct conjugate primes are comaximal; avoidance gives the
second equality in (ASG.3), and the class equality gives the last.
The use of \(H\) proves existence of an ideal in the original
\(E\); the resulting lattice is not extended to \(H\).
% END ASG3

% BEGIN ASG4
\paragraph{ASG4. Primitive negative vectors in an integral hyperbolic plane.}
\label{claim:ASG4}
Suppose \(\Lambda\) is unimodular, \(e,f\) form a hyperbolic
pair, and \(\mathfrak c\ne0\) satisfies
\(\mathfrak c e\subset\Lambda\) and
\(h(f,\Lambda)\subset\mathfrak c\).  If an integral trace-one
element exists, then for every positive integer \(T\) there is
\(z\in(Ee+Ef)\cap\Lambda\) such that
\[
 h(z,z)=-T,\qquad
       \gamma z\in\Lambda\ \Longrightarrow\ \gamma\in\mathcal O_E.
\tag{ASG.4}
\]
\begin{proof}
Choose \(\mathfrak p,\alpha\) from ASG3 with \(N=T\).
Choose \(p\in\mathfrak p\), \(q\in\bar{\mathfrak p}\)
with \(p+q=1\).  If \(t+\bar t=1\), put
\(w=pt+\bar q\bar t\).  Then \(w\in\mathfrak p\)
and \(w+\bar w=1\).  Define
\[
             b=\overline{-Tw/\alpha},\qquad z=\alpha e+bf.
\tag{ASG.5}
\]
Since \(\alpha\in\mathfrak c\), the first summand lies in
\(\Lambda\).  For \(y\in\Lambda\),
\(h(bf,y)=-Tw h(f,y)/\alpha\) is integral because
\(w\mathfrak c\subset\mathfrak p\mathfrak c=(\alpha)\).
Unimodularity therefore puts the second summand in \(\Lambda\).
The two cross terms give \(h(z,z)=-T(w+\bar w)=-T\).

Suppose \(\gamma z\in\Lambda\).  Pairing with \(f\)
gives \(\gamma\alpha\in\mathfrak c\), equivalently
\(\gamma\mathfrak p\subset\mathcal O_E\).  In particular
\(\gamma w\) is integral.  Pairing with \(\alpha e\in\Lambda\)
gives \(-T\gamma\bar w\in\mathcal O_E\).
Choose \(a\in\mathcal O_E\), \(p_0\in\mathfrak p\)
with \(aT+p_0=1\).  Then
\[
        \gamma=\gamma w+\bar w(\gamma p_0)
                              +a(T\gamma\bar w)
\]
is integral.  This proves actual primitiveness, not merely a nonzero
reduction at the two distinguished primes.
\end{proof}
% END ASG4

% BEGIN ASG5
\paragraph{ASG5. The extra isotropic direction supplied by the actual Forms0 data.}
\label{claim:ASG5}
The original \(\mathrm{Forms0}(D,D.u_0)\) supplies a nonzero
isotropic vector orthogonal to every original hyperbolic pair in \(D.sp\).
\begin{proof}
Its \(\mathrm{val0.de}\) field gives an invertible matrix
\(C\) over \(E\) with \(C^*C\) equal to the original Gram
matrix on \(D.sp.X(0)\).  That space has positive rank by the
rank inequalities of \(D\).  The inverse image under \(C\)
of one standard unit vector gives \(x\) with \(h(x,x)=1\).
Its \(\mathrm{atT.piU}\) field gives \(\pi\in\mathcal O_E\)
with \(\pi\bar\pi=T_{u_0}\).  Let \(v\) be the image of
\(\pi x\) under the original isometric inclusion into \(V\).
The original vector \(z_m\) is orthogonal to this space and has
norm \(-T_{u_0}\).  Therefore
\(w=v+z_m\) has norm zero.  It is nonzero, since
\(h(w,z_m)=-T_{u_0}\ne0\).  All original hyperbolic vectors
are orthogonal to both summands by the defining \(D.sp\) fields.
This proves the assertion in the original space without an isotropy
or approximation theorem as an additional premise.
\end{proof}
% END ASG5

% BEGIN ASG6
\paragraph{ASG6. A saturated isotropic module with the prescribed ideal factors.}
\label{claim:ASG6}
There are independent isotropic vectors \(e_1,\ldots,e_{s+1}\)
and a nonzero fractional ideal \(\mathfrak c\) for which
\[
 \Lambda\cap U=\bigoplus_{i=1}^{s}\mathfrak a_i e_i
                             \ \oplus\ \mathfrak c e_{s+1},
 \qquad U=\operatorname{span}_E(e_1,\ldots,e_{s+1}).
\tag{ASG.6}
\]
\begin{proof}
Adjoin the vector of ASG5 to the original \(s\) isotropic
vectors.  Pairing a linear relation with their original dual vectors
kills the first \(s\) coefficients, and the last vector is nonzero.
They are independent and their span \(U\) is totally isotropic.
The module \(M=\Lambda\cap U\) is saturated in \(\Lambda\):
if \(ay\in U\), \(0\ne a\in\mathcal O_E\), then \(y\in U\).
It has rank \(s+1\), since denominators in a basis of \(U\)
can be cleared into the full lattice \(\Lambda\).

Finite torsion-free modules over a Dedekind domain are projective and
are classified, at fixed positive rank, by their determinant ideal
class.  In the usual Steinitz normal form they are a direct sum of
free rank-one modules and one fractional ideal.  In particular
\(I\oplus J\simeq\mathcal O_E\oplus IJ\) for invertible
fractional ideals.  Choose the class of \(\mathfrak c\) to be
\([\det M]\prod_i[\mathfrak a_i]^{-1}\).
The normal form gives an isomorphism
\(\bigoplus_i\mathfrak a_i\oplus\mathfrak c\simeq M\).
Extend it to \(E\); the images of the standard coordinate vectors
are the desired \(e_i\).  This is an equality of the actual
embedded module after choosing its basis, as in (ASG.6).
All vectors remain in the same totally isotropic \(U\).
The statement also covers \(s=0\).
\end{proof}
% END ASG6

% BEGIN ASG7
\paragraph{ASG7. Integral dual hyperbolic vectors without division by two.}
\label{claim:ASG7}
Write \(c_i=\mathfrak a_i\) for \(i\le s\) and
\(c_{s+1}=\mathfrak c\).  There are isotropic dual vectors
\(f_1,\ldots,f_{s+1}\) with
\[
       h(e_i,f_j)=\delta_{ij},\qquad h(f_i,\Lambda)\subset c_i.
\tag{ASG.7}
\]
\begin{proof}
The saturated quotient \(\Lambda/M\) is finite torsion-free,
hence projective.  Split its exact sequence to obtain a retraction
\(\rho:\Lambda\to M\).  The \(i\)-th coefficient of
\(\rho(y)\) lies in \(c_i\).  Extend this coefficient map
linearly to \(V\); nondegeneracy represents it uniquely as
\(h(f_i^0,-)\).  Thus
\(h(f_i^0,e_j)=\delta_{ij}\) and
\(h(f_i^0,\Lambda)\subset c_i\).
Put \(B_{ij}=h(f_i^0,f_j^0)\).  With the trace-one element
\(t\) of ASG2 set
\[
 C_{ij}=\begin{cases}-B_{ij}&i<j,\\-tB_{ii}&i=j,\\0&i>j,
              \end{cases}
 \qquad f_j=f_j^0+\sum_i C_{ij}e_i.
\tag{ASG.8}
\]
Since \(B\) is Hermitian,
\(B_{ij}+C_{ij}+\overline{C_{ji}}=0\).
Expanding the pairing proves \(h(f_i,f_j)=0\); pairings with
the \(e_i\) are unchanged.

Here is the integrality check for the correction, including the last
ideal \(\mathfrak c\), which need not be conjugation invariant.
For \(y,l\in\Lambda\),
\(h(e_i,y)\in\bar c_i^{-1}\) because \(c_i e_i\subset\Lambda\).
Consequently
\(h(h(e_i,y)f_i^0,l)=\overline{h(e_i,y)}h(f_i^0,l)\)
is integral.  Unimodularity implies
\(h(e_i,y)f_i^0\in\Lambda\).  Pairing this vector with
\(f_j^0\) gives \(h(e_i,y)B_{ji}\in c_j\).
Each correction \(\overline{C_{ij}}h(e_i,y)\) is this element
times \(-1\), \(-\bar t\), or zero.  Hence it lies in
\(c_j\), proving \(h(f_j,\Lambda)\subset c_j\).
No integral division by two has occurred.
\end{proof}
% END ASG7

% BEGIN ASG8
\paragraph{ASG8. The actual orthogonal lattice splitting and all negative vectors.}
\label{claim:ASG8}
Keep the first \(s\) hyperbolic pairs from ASG7 and put
\(\Lambda_b=\Lambda\cap(\sum_i Ee_i+Ef_i)^\perp\).
Then
\[
 \Lambda=\bigoplus_{i=1}^{s}\mathfrak a_i e_i
              \ \oplus\ \Lambda_b
              \ \oplus\ \bigoplus_{i=1}^{s}\mathfrak a_i^{-1}f_i.
\tag{ASG.9}
\]
For every original word \(u\), there is a primitive
\(z_u\in\Lambda_b\) of norm \(-T_u\).
\begin{proof}
By ASG2 the ideals \(\mathfrak a_i\) are conjugation invariant.
Pairing with \(\mathfrak a_i e_i\) shows
\(h(e_i,\Lambda)\subset\mathfrak a_i^{-1}\).
Conversely, if \(b\in\mathfrak a_i^{-1}\), then
\(h(bf_i,\Lambda)\subset\bar b\mathfrak a_i\subset\mathcal O_E\);
unimodularity gives \(\mathfrak a_i^{-1}f_i\subset\Lambda\).
For \(y\in\Lambda\), subtract
\[
            \sum_i h(f_i,y)e_i+\sum_i h(e_i,y)f_i.
\]
Both sums are integral and the remainder is orthogonal to every
\(e_i,f_i\).  This gives (ASG.9); uniqueness follows by pairing
with the hyperbolic vectors and then taking the orthogonal remainder.

The reserved last pair is orthogonal to all the retained pairs.
It satisfies exactly ASG4 with \(\mathfrak c\), by ASG6--ASG7.
Apply ASG4 for each positive \(T_u\) and choose the resulting
\(z_u\) once for that word.  Every vector is in the same
\(\Lambda\), orthogonal to the retained pairs, and primitive in
\(\Lambda\).  This supplies the entire word-indexed family without
changing the common isotropic flag or any ideal \(\mathfrak a_i\).
\end{proof}
% END ASG8

% BEGIN ASG9
\paragraph{ASG9. Every original FixedHyp field and the full lattice block.}
\label{claim:ASG9}
The preceding data give the complete original
\(\mathrm{UnitaryLatticeData}(D)\), including
\(\mathrm{FixedHyp}\) for every word's flag.
\begin{proof}
Use \((e_i)_{i\le s},(f_i)_{i\le s},\Lambda_b,z_u\) from
ASG8 and let \(F_u\) have precisely that isotropic frame and
negative vector.  Its independence, isotropy, orthogonality and norm
are ASG6--ASG8.  The lattice, unimodularity and membership fields are
the original \(\Lambda\) and ASG8.

The integral intersection with \(U_0=\sum_{i\le s}Ee_i\)
is exactly \(\bigoplus\mathfrak a_i e_i\), either by (ASG.6)
or by the projection in (ASG.9).  At every prime \(P\) choose
\(c_{i,P}\in\mathfrak a_i\) generating
\((\mathfrak a_i)_P\).  Such a global member exists by choosing
a nonzero class in \(\mathfrak a_i/P\mathfrak a_i\); its
local coefficient is a unit.  At primes above \(g\), take
\(c_{i,P}=1\), as ASG2 says that these negative powers of the
different contain one and have unit localization there.
Thus localization of the intersection is generated by
\(c_{i,P}e_i\), with exactly the source's prescribed value one
at every distinguished prime.

The remaining primitive condition is slightly stronger in notation
than primitiveness in \(\Lambda\), and can be checked directly.
If \(\gamma z_u=l+x\), with
\(l\in\Lambda\cap U_0^\perp\) and \(x\in U_0\), subtract
\(\sum_i h(f_i,l)e_i\) from \(l\).  This subtraction is
integral, by ASG7, and the remainder is orthogonal to both frames.
Pairing the displayed equality with each \(f_i\) shows that
\(x+\sum_i h(f_i,l)e_i=0\).  Hence \(\gamma z_u\in\Lambda\),
and ASG8 gives \(\gamma\in\mathcal O_E\).
This is precisely primitiveness in the source's \(Lsub\).
All \(\mathrm{FixedHyp}\) fields have now been supplied.
The remaining fields \(\mathrm{Fu\_\allowbreak e}\) and
\(\mathrm{flagZVector}\) hold by the definitions of \(F_u\).
\end{proof}
% END ASG9

% BEGIN ASG10
\paragraph{ASG10. The finite additive character with its exact infinity normalization.}
\label{claim:ASG10}
For every original \(\mathrm{AutInst}\), and hence for ASG1,
there is the complete original \(\mathrm{AdditiveCharacterData}\).
\begin{proof}
On \(\mathbb Q_p\) let \(\{x\}_p\in\mathbb Z[1/p]/\mathbb Z\)
be its negative-power part.  Define
\[
 \psi_{\mathbb Q,f}((x_p))
       =\exp\!\left(-2\pi i\sum_p\{x_p\}_p\right).
\tag{ASG.10}
\]
Only finitely many terms are nonzero, and their classes are additive.
This is a continuous unitary additive character, trivial on
\(\prod_p\mathbb Z_p\).  For \(x\in\mathbb Q\), subtracting
its finite principal parts leaves an integer, so
\(\psi_{\mathbb Q,f}(x)e^{2\pi ix}=1\).
Identify \(E\otimes_F\mathbb A_{F,f}\) with
\(\mathbb A_{E,f}\) and compose (ASG.10) with the adelic trace
\(\operatorname{Tr}_{E/\mathbb Q}\).
The trace of an integral local element is integral, proving
\(\psi_{E,f}=1\) on the actual \(\mathrm{IntAf}\).

The original infinity factor is
\[
 \prod_v \exp\bigl(2\pi i(\phi_v(x)+\overline{\phi_v(x)})\bigr)
               =\exp(2\pi i\operatorname{Tr}_{E/\mathbb Q}x),
\]
because the \(\phi_v\)'s select exactly one member of each CM
pair.  The rational product identity just proved gives the literal
\(\mathrm{psiEf\_\allowbreak rat}\), with its stated sign and without an
extra factor of two.  Additivity and integral triviality give the other
two fields.  Thus this also proves the entire generic
\texttt{Sc\allowbreak Psi\allowbreak Exists.\allowbreak Target}.
\end{proof}
% END ASG10

% BEGIN ASG11
\paragraph{ASG11. Split local lattices and the literal congruence predicates.}
\label{claim:ASG11}
At every finite place \(v\) of \(F\) split in \(E\), a
self-dual integral Hermitian lattice is a pair of dual free
\(\mathcal O_{F_v}\)-modules.  For positive integers \(g,m\)
and finite adelic vectors, the original predicate
\(\mathrm{InAt}(g,m)\) means precisely membership
in \(m\widehat\Lambda_v\) at every place dividing \(g\).
\begin{proof}
Write \(R=\mathcal O_{F_v}\).  The integral quadratic algebra
is \(R\times R\), with conjugation interchanging its factors.
The two idempotents decompose the lattice as \(M^+\oplus M^-\).
The Hermitian form is zero within either summand and gives a perfect
pairing between them, by self-duality.  Both summands are finite free
over the discrete valuation ring \(R\).  Thus either choice of a
basis on one side determines a dual basis on the other.

The adelic lattice \(I.Lf\) is the product of the local completed
lattices.  One can verify this from its definition as the integral
adelic span of a finite set of global generators: locally these
generators span, and their finitely many coefficient coordinates can
be chosen independently in the product of local rings.
Consequently \(\mathrm{mL}(m)=m\widehat\Lambda\).
If \(c\) is coprime to \(g\), multiplication by \(c\) is a
unit at every place dividing \(g\), proving one implication in
the asserted interpretation of \(\mathrm{InAt}\).
For the converse, clear the finitely many remaining denominator and
divisibility defects by a positive rational integer \(c\) involving
only rational primes not dividing \(g\).  The definition of a
finite adele makes this a finite operation.  Then
\(c x\in m\widehat\Lambda\) and \(\gcd(c,g)=1\), exactly
the required witness.  Similarly \(\mathrm{ZMod'}(g,a)\)
means \(a\in g\widehat{\mathcal O}_E\), and
for integral vectors and integral supplied generators,
\(\mathrm{InMod}\) is the stated residue span at all places
dividing \(g\).  Integrality supplies the conditions away from
\(g\), where \(g\) is a unit.  These interpretations also
cover \(g=1\).
\end{proof}
% END ASG11

% BEGIN ASG12
\paragraph{ASG12. The full original reference frame by finite adelic construction.}
\label{claim:ASG12}
The generic \texttt{Sc\allowbreak Ref\allowbreak Exists.\allowbreak Target} holds.  In particular,
ASG1 and ASG9 have an actual \(\mathrm{ReferenceFrame}\).
\begin{proof}
Use exactly the hypotheses of that target: \(g\ge1\), splitting
of \(E/F\) at every \(v\mid g\), \(n=2r+k+1\), a
hyperbolic frame \(e_i,f_i\) integral at these places, and the
given self-dual lattice.  At such a place the projection
\(y\mapsto\sum_i h(f_i,y)e_i+h(e_i,y)f_i\) is integral.
It splits off the hyperbolic pairs and leaves a self-dual orthogonal
complement of rank \(k+1\) over \(R\times R\).
Choose bases \(u_0,\ldots,u_k\) and \(v_0,\ldots,v_k\)
of its two summands with
\(h(v_i,u_j)=\delta_{ij}(1,0)\), and put
\[
 e_{0,i}=e_i,\quad f_{0,i}=f_i,\quad z_0=u_k-v_k,
 \quad \beta_i=u_i,\quad\beta_i^d=v_i\ (i<k),
 \quad e_w=(1,0).
\tag{ASG.11}
\]
These are local choices at finitely many places.  At all other places
take all these vectors and \(e_w\) to be zero.  They are genuine
integral finite adeles.  ASG11 proves \(\mathrm{basisApprox}\):
the differences from the rational frames vanish at every place above
\(g\), and their other denominators can be cleared away from \(g\).

At the selected places the Gram equalities for the two frames are
exact.  Also \(h(z_0,z_0)=-1\); every required cross pairing
vanishes; \(e_w^2=e_w\), \(e_w\bar e_w=0\),
\(\beta_i=e_w\beta_i\),
\(\beta_i^d=\bar e_w\beta_i^d\), and
\(h(\beta_i^d,\beta_j)=\delta_{ij}e_w\).
Off these places \(g\) is a unit, so all required congruences
are automatic for integral vectors.  The exact idempotent and support
equalities also hold there because both sides are zero.
The condition \((1-e_w-\bar e_w)y\in g\widehat\Lambda\)
holds place by place: its left side is zero above \(g\) and
the modulus is a unit elsewhere.

Finally the displayed vectors span modulo \(g\).
The first pairs span their hyperbolic part; \(\beta_i,\beta_i^d\)
give the first \(k\) coordinates of the complement; and
\(e_w z_0=u_k\), \(\bar e_w z_0=-v_k\) give its last two
components.  This proves the actual \(\mathrm{latticeSpan}\)
field, including \(k=0\).  Every field of \(\mathrm{RefData}\)
has been verified.  For \(g=1\) the empty local construction gives
zero vectors and zero idempotent, and every congruence is modulo a unit,
so it is an allowed instance as well.
The original \(D\)-specialization follows from ASG2 and ASG9.
\end{proof}
% END ASG12

% BEGIN ASG13
\paragraph{ASG13. Transport to any original reference frame, with exact support.}
\label{claim:ASG13}
The full generic \texttt{Sc\allowbreak Xu\allowbreak Exists.\allowbreak Target} holds for any supplied
reference frame, not just the particular frame of ASG12.
It supplies \(\mathrm{TransportData}\) for every word in ASG9.
\begin{proof}
At a place \(v\mid g\), split off the exact integral hyperbolic
frame as in ASG12 and call its orthogonal complement \(W\).
The given vector \(z\) lies in \(W\) and has norm \(-t\),
where \(t\equiv1\pmod g\).  Reduce modulo \(gR\).
The reference vector has norm \(-1\) and is orthogonal to the
reduced hyperbolic frame, by \(\mathrm{basisApprox}\) and
\(\mathrm{formsAtZero}\).  Its projection to \(W\) therefore
has the same reduction.  In the split description of ASG11 a vector
is a pair \((u,\lambda)\) with \(\lambda(u)=-1\).

Over the local Artinian ring \(R/gR\), general linear matrices act
transitively on such pairs by
\((u,\lambda)\mapsto(Au,\lambda A^{-1})\).
To see this, \(u\) is unimodular and
\((R/gR)^{k+1}=(R/gR)u\oplus\ker\lambda\).
Choose a basis of the free complement; it identifies the pair with
\((e_1,-e_1^*)\).  Do this for both pairs and compose the basis
changes.  Thus choose \(\bar A\) taking the reference pair to
the reduction of \(z\).  Lift its entries to \(R\); its
determinant remains a unit, so the lift is in \(\mathrm{GL}_{k+1}(R)\).
Together with its inverse dual it is an integral unitary operator
on \(W\).  Extend it by the identity on the hyperbolic part.

Use these operators at \(v\mid g\), and the identity elsewhere,
to obtain \(x\in G_f\).  It preserves the entire adelic lattice
and fixes every original \(e_i,f_i\) exactly.  Its inverse sends
\(z\) to the reference \(z_0\) modulo \(g\) at all selected
places, which is the literal \(\mathrm{InAt}(g,g)\) by ASG11.
If \(m\) is coprime to \(g\), then \(x-1\) is zero at
every place above \(m\); elsewhere it preserves integral lattices
and \(m\) is a unit.  Hence \((x-1)y\in m\widehat\Lambda\)
for all \(y\in\widehat\Lambda\), proving \(\mathrm{OneOff}\)
with clearing integer one.  For \(g=1\), use \(x=1\) and
clear the remaining rational denominators as in ASG11.
All four conclusion fields of the target are proved.

For the original specialization, \(t=T_u\) has the required
congruence by the original word datum.  ASG8--ASG9 give its norm,
orthogonality and local integrality.  Choose the resulting \(x_u\)
for every word with one already fixed reference frame.  These are
exactly \(\mathrm{xu\_\allowbreak fix},\mathrm{xu\_\allowbreak z},\mathrm{xu\_\allowbreak one}\)
and \(\mathrm{xu\_\allowbreak K}\).
\end{proof}
% END ASG13

% BEGIN ASG14
\paragraph{ASG14. The actual finite-level character idempotent is admissible.}
\label{claim:ASG14}
For the original two-prime modulus, any reference frame as above, the
character of ASG10, and any positive Haar measure \(\mu_f\), the
literal \(\mathrm{eOf}\) is an integrable
\(\mathrm{IsAdmIdem}(\mu_f,g,g_0)\).
\begin{proof}
Here the exact finite calculation of ARO7--ARO10 uses only the lattice,
reference frame, character and modulus, all now available before any
\(\mathrm{HidaSetting}\).  We recall the argument to make this
independence explicit.  At each conjugate pair of residue places the
nonzero parts of the frames give dual bases, since their pairing matrix
is the identity and \(\mathrm{latticeSpan}\) gives spanning.
Order one basis as \((e;\beta;z;f_{\rm reversed})\); in the
other orientation use the reversed active dual \(\beta\)-frame.
The \(K_2\) congruences make the direct \(e,\beta,z\) columns
upper unitriangular, with the middle entries of the \(z\) column
zero.  Applying the conjugate conditions to the inverse dual matrix
does the same for the remaining columns.  Conversely that matrix
shape gives all defining congruences, at both orientations.

Such matrices form a subgroup.  Every coordinate used in the original
argument
\[
 C(x)=\sum_a h(f_{0,a},x\,\mathrm{next}(e_0,z_0)_a)
                    +\sum_{i<k-1}h(\beta_i^d,x\beta_{i+1})
\tag{ASG.12}
\]
is additive under multiplication, modulo the modulus.  For successive
entries this is ordinary unitriangular matrix multiplication.  For the
link from the last \(e\) to \(z\), the possible intermediate
terms are zero because the intervening entries in the \(z\) column
are zero.  At the conjugate orientation inverse transpose reverses the
chains and changes their signs; these are still additive coordinates.
Thus \(C(xy)-C(x)-C(y)\in g\widehat{\mathcal O}_E\).
It follows that the literal \(\chi(x)=\psi_{E,f}(C(x)/g)\)
is a character of \(K_2\), trivial on \(K_g\).
It has modulus one and factors through the finite residue quotient.

The subgroup \(K_2\) contains \(K_g\) and lies in the compact
open lattice stabilizer.  It is therefore compact open, and
\(v=\mu_f(K_2)\) is positive and finite.  Define exactly
\(e(x)=v^{-1}\overline{\chi(x)}\) on \(K_2\), zero elsewhere.
This is locally constant, compactly supported and integrable.
Convolution on \(K_2\) gives
\(e*e(x)=v\,v^{-2}\overline{\chi(x)}=e(x)\);
outside \(K_2\) no pair of nonzero factors can occur.
Inversion gives \(e(x^{-1})=\overline{e(x)}\), and the kernel
condition gives the prescribed two-sided \(K_g\)-invariance.

For full positive type, group finitely many \(x_i\)'s by their
right cosets \(x_iK_2\).  In a group write \(x_i=x_C k_i\).
Then its contribution to
\(\sum_{i,j}\bar c_i c_j e(x_i^{-1}x_j)\) is
\(v^{-1}|\sum_i c_i\overline{\chi(k_i)}|^2\).
Different groups contribute zero.  This proves nonnegativity for all
finite lists and coefficients, not only positivity of one operator.
Finally every supporting residue matrix is unitriangular, so
\((x-1)^n\) is zero modulo each prime dividing \(g_0\).
This is precisely \(\mathrm{UnipMod}(g_0,x)\).
All seven clauses of \(\mathrm{IsAdmIdem}\) are proved for the
original \(\mathrm{eOf}\), with its actual Haar denominator.
\end{proof}
% END ASG14

% BEGIN ASG15
\paragraph{ASG15. The full automorphic bundle under the original gap.}
\label{claim:ASG15}
For the actual blocks constructed above and the original
\(\mathrm{GapHyp}(D)\), there exists a complete
\(\mathrm{AutomorphicBundle}\).  The separate weighted local
coefficient statement used by one source route is unnecessary for
this bundle's unweighted integrability field.
\begin{proof}
Choose the Haar measure of \(G_f\) with \(\mu_f(K)=1\).
It is sigma finite because the group is second countable and locally
compact.  FV6 supplies a measurable fundamental domain of the actual
rational subgroup; RQC2--RQC6 supply its quotient Hilbert space and
the original strongly continuous unitary right regular representation
\(R\).  These choices fix the measure and fundamental domain for
all fields of this bundle.

At each real place choose the actual discrete-series representation of
DC15 at \(\epsilon_v,\lambda_v=\mathrm{lamOf}(A,v)\).
DC14 verifies its hypotheses from the literal rank bound \(1801\)
and gaps \(16n^{19}\); all casts and sign orders are those of
ASG1.  Choose a nonzero vector in its lowest compact type and normalize
it to \(w_v\).  The coefficient identity of DC11 and decay estimate
of DC6 give
\(g\mapsto\langle w_v,\sigma_v(g)w_v\rangle\in L^1\)
for the original infinity Haar measure.  These statements apply to the
chosen actual representations, so no separate assumption about
\(\mathrm{CoefWeightedLoc}\) is needed.  DC12 supplies the complete
four-vector Schur identity and identifies its positive degree with the
literal \(\mathrm{fdegOf}\).

Take the completed finite Hilbert tensor product of these spaces.
Its elementary tensor map has inner product the product of local inner
products; elementary tensors span densely.  The product representation
acts componentwise and is strongly continuous by density and isometry.
This is the explicit tensor construction also used in RFD4, and proves
every field of \(\mathrm{IsHilbTensor}\).
With \(d_v=\mathrm{fdegOf}(\sigma_v,w_v)\), the exact function is
\[
 \phi_\infty(g)=\left(\prod_v d_v\right)
            \overline{\prod_v\langle w_v,\sigma_v(g_v)w_v\rangle}.
\tag{ASG.13}
\]
Tonelli on the product of absolute values proves its \(L^1\)
integrability, including the degree constants.

ASG14 gives the exact integrable finite test function and every
admissibility field.  For any integrable scalar kernel \(a\) and
unitary representation \(U\), the sesquilinear form
\(\int a(g)\langle x,U(g)y\rangle\,dg\) has absolute value
at most \(\|a\|_1\|x\|\|y\|\).  Riesz representation
therefore gives its unique bounded integral operator, of norm at most
\(\|a\|_1\).  Apply this to the finite restriction of \(R\)
with \(a=e\), and to its infinity restriction with
\(a=\phi_\infty\).  These are the required \(R_e,R_\phi\)
and precisely their weak-integral identities.
We have now supplied \(\mu_f,\mathrm{fund},H_v,\sigma_v,w_v\),
the tensor space and map, \(R\), \(e\)'s admissibility,
\(R_e\), all discrete-series and lowest-vector fields,
\(R_\phi\), and \(\phi_\infty\)'s integrability.
These are all fields of \(\mathrm{AutomorphicBundle}\).
\end{proof}
% END ASG15

% BEGIN ASG16
\paragraph{ASG16. The actual finite representation block on every stabilizer.}
\label{claim:ASG16}
For the lattice block of ASG9 there is the complete original
\(\mathrm{RepresentationData}\), with all original highest weights.
\begin{proof}
At a real place the hyperbolic pairs span a nondegenerate space of
signature \((s,s)\).  Their orthogonal complement has signature
\((k,1)\).  Since \(h(z_u,z_u)=-T_u<0\), the further
orthogonal complement of \(z_u\) is positive definite of dimension
\(k\).  Gram--Schmidt gives an orthonormal basis \(w_{u,v,i}\)
there.  These are exactly the \(\mathrm{wfr\_\allowbreak on}\) and
\(\mathrm{wfr\_\allowbreak perp}\) equations after the original real
coordinate map \(\mathrm{rv}\).

Let \(U\) be the span of the isotropic frame.  The real stabilizer
\(S\) in the source acts on \(U^\perp/U\), fixes the class
of \(z_u\), and acts unitarily on its positive orthogonal
complement.  In the chosen basis this gives a continuous homomorphism
\(q_v:S\to U(k)\).  The subgroup \(N\) of the source is
exactly its kernel: its defining last condition says that the action
on \(U^\perp/U\) is the identity, and the class of \(z_u\)
is already fixed by \(S\).

Write \(\nu_i=D.hcW(u,v).u_i\), with indices starting at zero.
These are the original Harish--Chandra parameters, not the integral
highest weights.  For the all-positive sign sequence,
\(\rho_N=0\) and \(\rho_C(i)=(k-1)/2-i\).  Thus the
highest weight required by the literal \(\mathrm{IsHCTypeU}\) is
\[
       \alpha_i=\nu_i-\frac{k-1}{2}+i.
\]
The original rank \(k\) is even and each \(\nu_i\) is
half-integral, so \(\alpha_i\) is an integer.  The original
interlacing has an integral parameter strictly between
\(\nu_i\) and \(\nu_{i+1}\); hence
\(\nu_i-\nu_{i+1}\ge1\), and
\(\alpha_i-\alpha_{i+1}\ge0\).
Apply the Schur-functor and determinant-power construction of DC1
to this nonincreasing integral sequence \(\alpha\).
The resulting finite-dimensional unitary representation
\(\pi_{u,v}\) of \(U(k)\) has precisely the required
\(\mathrm{lowK}(\mathrm{allPos},\nu)\) highest weight.
It therefore satisfies the full original
\(\mathrm{IsHCTypeU}\) with parameter \(\nu\), including
possibly negative last highest weight through the determinant power.
The finite Hilbert tensor product of their spaces is \(V_u\).
On \(S\) define \(\pi_u(s)=\bigotimes_v\pi_{u,v}(q_v(s))\);
extend the function to all of \(G_\infty\) by zero outside \(S\).
The source imposes its representation and continuity laws only on
\(S\).  There they follow from the homomorphisms \(q_v\),
unitarity, and the finite tensor construction; on \(N\) the
operator is the identity.  This proves every \(\mathrm{IsRepS}\)
clause, including continuity on \(S\).

The subgroup \(H\) fixing all \(e_i,f_i,z_u\) acts only on
the positive complement, so its \(q_v\) matrix is exactly
\(\mathrm{matW}(w_{u,v},v,s)\), by its defining pairings.
Thus elementary tensors satisfy the literal
\(\mathrm{tensU\_\allowbreak equiv}\) formula.  Their inner products are
the products of local inner products and they span the full tensor
space (which is finite dimensional).  This proves both other tensor
fields and completes the original representation block for every word.
\end{proof}
% END ASG16

% BEGIN ASG17
\paragraph{ASG17. The original subgroup and quotient measures.}
\label{claim:ASG17}
The same lattice block has a \(\mathrm{MuBlocks}\) object with
the literal infinity Haar measure in its quotient integral identity.
\begin{proof}
The source's groups \(S,N,H\) are closed real matrix subgroups:
their defining flag inclusions are closed linear conditions and their
unitary conditions are closed equations.  The map \(q\) of ASG16
has kernel \(N\), and its restriction to \(H\) is an isomorphism
onto \(\prod_v U(k)\).  Extending an operator on the positive
complement by the identity on \(e_i,f_i,z_u\) proves surjectivity.
Consequently \(S=N\rtimes H\), with multiplication a homeomorphism
from \(N\times H\) onto \(S\): the inverse sends \(s\) to
\((s q(s)^{-1},q(s))\), using this section.

In a basis adapted to \(U\subset U^\perp\), every element of
\(N\) is upper unitriangular.  Indeed the flag action has diagonal
one on \(U\), its middle action is the identity by definition,
and unitarity gives the inverse dual action on the last block.
Thus \(N\) is a real unipotent group and is unimodular.
For completeness its Lie algebra is nilpotent; every adjoint
endomorphism has trace zero, so the modular character is one on its
identity component.  The exponential and finite logarithm polynomials
identify a real unipotent group with its nilpotent Lie algebra, proving
connectedness and the assertion on the whole group.
The group \(H\) is compact.  Its action on a Haar measure of \(N\)
has a positive scalar modulus, a continuous homomorphism
\(H\to\mathbb R_{>0}\).  A compact subgroup of
\(\mathbb R_{>0}\) is trivial, so this action preserves that Haar
measure.  The pushforward of \(\mu_N\times\mu_H\) is therefore
a left and right Haar measure \(\mu_S\) on \(S\).

Regard \(\mu_N,\mu_H\) as measures on the ambient group by
closed inclusion.  They are sigma finite, nonzero, supported on their
subgroups and invariant under their left translations, which are the
four \(\mathrm{IsHaarOnSet}\) fields.
The classical quotient Haar theorem applies to the unimodular closed
subgroup \(S\) of the unimodular real reductive group \(G_\infty\).
Choose its quotient measure relative to this \(\mu_S\) and the
original \(\mu_\infty\).  The second-countable quotient has a
Borel section; push that quotient measure forward to a measure
\(\nu_Q\) on \(G_\infty\).  Weil's identity is then
\[
 \int_{G_\infty} f\,d\mu_\infty
       =\int_{G_\infty}\int_S f(sx)\,d\mu_S(s)\,d\nu_Q(x)
\tag{ASG.14}
\]
for every nonnegative measurable \(f\).  This standard form follows
first for continuous compactly supported functions and then for all
nonnegative measurable functions by equality of the induced Radon
measures and monotone approximation.  It is precisely
\(\mathrm{IsQuotMeasS}\).  The ambient measure has not been
rescaled: the normalization freedom was used in choosing the quotient
measure.  Choose these measures for each word.  This supplies every
\(\mathrm{MuBlocks}\) field, without yet asserting a global
Siegel--Weil measure comparison.
\end{proof}
% END ASG17

% BEGIN ASG18
\paragraph{ASG18. Exact supplied blocks and the remaining common-setting boundary.}
\label{claim:ASG18}
For one original \(D\), its supplied \(\mathrm{Forms0}(D,D.u_0)\),
and \(\mathrm{GapHyp}(D)\), the preceding proofs supply,
simultaneously and on the same objects, the original blocks
\(A,L,P,R_f,X_u,B_u,\Pi,M_\mu\).
They also supply \(\mathrm{RowsFinDim}\) for its actual bundle
through RFD8.  This does not alone assert the full common
\(\mathrm{ConcreteSettingExists}\) or any entire remaining global leaf.
\begin{proof}
Choose \(A\) by ASG1, then \(L\) by ASG9, then \(P\)
by ASG10.  Choose one \(R_f\) by ASG12 and use it for all the
\(X_u\)'s of ASG13 and the bundle of ASG15.  Choose the
representation and measure blocks by ASG16--ASG17 on that same
\(A,L\).  Each construction retains the preceding choices, so
these are dependent blocks of the original constructor, rather than
unrelated existential witnesses.  All word quantifiers occur after
the common lattice and frame choices.  Apply the universal RFD8 to
the actual measure, fundamental domain, finite kernel and infinity
representations now constructed.  It supplies the exact
finite-dimensionality premise used by ARO18 and subsequent row proofs.

The constructor still requires the original nonzero Bessel-function
block \(W\) with its normalized quotient integral and decay, and
the original measured-field block.  The global endpoint additionally
requires the original Siegel--Weil rescaling comparison and its
same-object measure and character ties.  The spectral endpoint also
uses its Bessel uniqueness and local-tie inputs.  These are retained
as their actual source statements; none is inferred from the algebraic
frame or from the existence of Haar measures.  The independent
\(\mathrm{CountIdempotentsClass}\) producer and a common family
threshold likewise remain outside this chapter.
Thus the exact setting-building inputs listed above are now supplied;
substituting them into the original conditional constructors gives the
claimed reduction of those constructors' outstanding inputs, and no
stronger assertion.
\end{proof}
% END ASG18


\section{The original canonical cusp connection over arbitrary coefficient algebras}
\label{sec:cnc-v47}

We construct the connection required by the original predicate
\(\mathrm{IsNab}\).  In particular, the connection hypothesis retained in
ICL11 is supplied on each original moduli object.  The construction uses
the given universal abelian scheme and its line bundle; it does not
choose a different moduli space, complex embedding, Hermitian form,
metric, measure, or tuple phase.  Principal analytic charts and the
finite theta system remain separate producers.

Fix a CM number field \(E\), its original host \(H\), a datum
\(D_w\), a level \(l\), a moduli object \(X_o:\mathrm{Moduli}(D_w,l)\),
and the original coefficient field \(K=c.F\).  Write
\[
 \begin{gathered}
 O=\mathcal O_E,\qquad A=X_o.\mathrm{univ}.t.G,
 \qquad L=X_o.\mathrm{univ}.t.PP.L,\\
 I=\coprod_{\tau:E\hookrightarrow\mathbb C}
            \{0,\ldots,D_w.\mathrm{mult}(\tau)-1\}.
 \end{gathered}
                                                               \tag{CNC.1}
\]
The involution on \(O\) is \(a^*=H.\mathrm{bar}(a)\).  It is an
involution because its image in \(E\) is complex conjugation.  The
element \(w_0\) is the original \(H.w_0\).
For a commutative \(K\)-algebra \(S\), no finiteness, reducedness,
locality, or domain hypothesis is imposed.  Put
\[
 J_m(S,n)=S[t_{b,i}:1\leq b\leq n,\ 1\leq i\leq d]
                    /(t)^{m+1}.
\]
All infinite series used below are formal series.  Their substitutions,
exponentials, logarithms, and integrals are computed coefficientwise;
at each fixed degree every calculation takes place in a finite jet
algebra.  There is no analytic convergence assertion about \(S\).

Two general geometric foundations enter.  A smooth morphism with a
section has locally power-series formal coordinates along that section;
the conormal graded algebra is its symmetric algebra.  Also, a proper
flat morphism of finite presentation with geometrically reduced connected
fibres has universally only base functions.  Precise primary references
are the Stacks Project, tags
\href{https://stacks.math.columbia.edu/tag/05D9}{05D9} and
\href{https://stacks.math.columbia.edu/tag/0E0L}{0E0L}, respectively.
The passage from local coordinates to the globally framed coefficient
algebra, the normalization of all line generators, and the formal gauge
equations are proved here.  The supplied cubical, adjoint, and primitive
line-bundle identifications are fields of the original \(\mathrm{SAB}\)
object, not additional existence hypotheses.

% BEGIN CNC1
\paragraph{CNC1. The given family and uniqueness of normalized line isomorphisms.}
\label{claim:CNC1}
The universal group \(A\to X_o.X\) and each of its finite powers are
abelian schemes.  After every base change \(T\to X_o.X\),
\[
 \mathcal O_T\ \xrightarrow{\ \sim\ }\ (A_T^r\to T)_*
                       \mathcal O_{A_T^r}\qquad(r\geq0).       \tag{CNC.2}
\]
Consequently two isomorphisms between the same two line bundles on
\(A_T^r\) agree if their ratio is one on the zero section.  This assertion
holds for nonreduced \(T\).
\begin{proof}
The original \(\mathrm{SAB}.\mathrm{semi}\) supplies smoothness,
commutativity and geometrically connected fibres.  The field
\(\mathrm{Obj}.\mathrm{proper}\), applied to \(X_o.\mathrm{univ}\),
supplies properness.  These are exactly the four fields of the original
\(\mathrm{IsAbelianScheme}\), as is also visible in the source proof of
\texttt{Obj\allowbreak Dim.\allowbreak is\allowbreak Abelian\allowbreak Scheme}.  Smoothness gives flatness, local finite
presentation and geometrically reduced fibres.  Properness supplies
quasi-compactness and quasi-separatedness, so the morphism is of finite
presentation.  These properties persist under base change and finite
products; the product of the geometrically connected smooth fibres is
geometrically connected and reduced.  Apply the universal base-functions
theorem just stated.  For \(r=0\) the map is the identity of \(T\).

The ratio of two line-bundle isomorphisms is a global invertible
function: this follows in a local basis and the ratios agree on overlaps.
By (CNC.2) it is pulled back from a unique invertible function on \(T\).
The zero section is a left inverse of this pullback, so ratio one there
forces ratio one everywhere.  This also proves uniqueness of a line
isomorphism after its zero-section normalization.  The argument uses
actual functions and their pullbacks, not tests on geometric points.
\end{proof}
% END CNC1

% BEGIN CNC2
\paragraph{CNC2. The original eigenframes form the whole tangent basis.}
\label{claim:CNC2}
For every original framed point \(y\in\mathrm{FrPt}(X_o,c,S)\), its
vectors \(y.u_{\tau,i}\), indexed by \(I\), form an \(S\)-basis of
the whole tangent module of \(A_S\) at zero.
\begin{proof}
The definition of \(\mathrm{Coeff}\) puts every embedding of \(E\)
inside \(K\).  Hence
\[
 E\otimes_{\mathbb Q}K\ \longrightarrow\prod_{\tau:E\hookrightarrow K}K,
       \qquad a\otimes b\longmapsto(\tau(a)b)_\tau             \tag{CNC.3}
\]
is an isomorphism.  Here is an explicit verification sufficient for
arbitrary \(S\).  Choose an integral basis \(b_1,\ldots,b_e\) of
\(O\).  The matrix \((\tau(b_j))_{\tau,j}\) is invertible over
\(K\): its squared determinant is the nonzero discriminant, after the
chosen embeddings are ordered.  Its inverse expresses every coordinate
idempotent of the product in the basis \(b_j\).  It also proves
multiplicativity and bijectivity of (CNC.3), by checking on the displayed
coordinates.

The differential of the given \(O\)-action extends uniquely
\(\mathbb Q\)-linearly to \(E\).  Combining it with the scalar
\(K\)-action gives an action of the left side of (CNC.3) on the tangent
module.  Its coordinate idempotents \(p_\tau\) are mutually orthogonal
and sum to one.  Their images are exactly
\[
 M_\tau=\{v:\ a v=\tau(a)v\text{ for every }a\in O\}.
\]
Indeed the coordinate relation proves the inclusion into \(M_\tau\);
conversely the inverse embedding matrix shows that \(p_\sigma v\)
is zero for \(\sigma\ne\tau\) and equals \(v\) for
\(\sigma=\tau\).  Thus the whole tangent module is
\(\bigoplus_\tau M_\tau\), over the given \(S\), without making
any assertion about integral eigenspaces at ramified primes.
The fields \(y.\mathrm{tan}\), \(y.\mathrm{eig}\), and
\(y.\mathrm{basis}\) say precisely that the indicated vectors are
bases of these summands.  Their union is a basis.  The proof remains
valid if \(S\) has nilpotents or is the zero ring.
\end{proof}
% END CNC2

% BEGIN CNC3
\paragraph{CNC3. Compatible formal coordinates with the specified tangent vectors.}
\label{claim:CNC3}
There exist \(d\), a commutative formal group law \(F\) over \(S\),
a formal \(O\)-action \(a\mapsto[a]\), coordinate vectors
\(u_\iota\in S^d\), and maps
\(x_m:\operatorname{Spec}J_m(S,1)\to A_S\) which satisfy every
field of the original \(\mathrm{IsChart}\) at \(y\).
\begin{proof}
Take \(d=|I|\) and enumerate the basis of CNC2.  Its dual is a basis
of the conormal module \(\mathcal I/\mathcal I^2\) of the zero
section.  The section is closed since \(A_S\to\operatorname{Spec}S\)
is separated.  Its order-\(m\) infinitesimal neighbourhood is affine:
a scheme which is a nilpotent thickening of an affine scheme is affine.
Write its ring as \(B_m\), with augmentation \(B_m\to S\).
The structural map \(S\to B_m\) splits this augmentation.

The smooth-section coordinate theorem gives, locally on the base,
\(\operatorname{gr}_{\mathcal I}\mathcal O_{A_S}
 \simeq\operatorname{Sym}_S(\mathcal I/\mathcal I^2)\).
This isomorphism is intrinsic and so glues.  Choose lifts of the given
conormal basis to the augmentation ideal of \(B_1\), then inductively
lift them through \(B_{m+1}\twoheadrightarrow B_m\).
Surjectivity on these global sections follows because the neighbourhoods
are affine and the kernels are quasi-coherent.  Subtracting the
augmentation preserves the chosen lift and places it in the augmentation
ideal.  The resulting maps
\[
 S[t_1,\ldots,t_d]/(t)^{m+1}\longrightarrow B_m              \tag{CNC.4}
\]
are isomorphisms: their maps on the degree-zero term are identities and
on every positive associated graded term are the symmetric powers of
the chosen conormal basis.  Induction on the finite filtration proves
injectivity and surjectivity of (CNC.4).  They commute with truncation.
This proves the needed global chart over the affine base, rather than
assuming that one local neighbourhood contains the whole zero section.

Transport the actual group multiplication and its action maps through
these formal coordinates.  They give series \(F(x,z)\) and \([a](x)\)
with zero constant term.  The identities of the group and action hold
at every jet order, hence coefficientwise: \(F\) is associative and
commutative with zero as identity; \([a+b](x)=F([a]x,[b]x)\),
\([ab]=[a]\circ[b]\), and \([1]=\mathrm{id}\).
The maps \(x_m\) are the inverse charts followed by the actual
inclusion into \(A_S\).  Their zero and truncation identities, their
group-law identity and their action identity are the four corresponding
fields of \(\mathrm{IsChart}\).

The differential of (CNC.4) is the chosen dual basis.  Therefore its
coordinate tangent vectors are the prescribed vectors of \(y\).
Take \(u_\iota\) to be the corresponding standard coordinate vector.
This proves the fields \(\mathrm{tan}\) and \(\mathrm{basis}\),
including uniqueness of every linear combination.  No Noetherian
assumption on \(S\) has been used.
\end{proof}
% END CNC3

% BEGIN CNC4
\paragraph{CNC4. Nilpotent lifting and normalized global primitive data.}
\label{claim:CNC4}
An invertible module over a ring with nilpotent ideal admits a generator
lifting any generator of its reduction.  The given primitive and adjoint
line-bundle isomorphisms consequently admit the normalizations used in
the original definitions \(\mathrm{IsPrim}\) and \(\mathrm{IsFrame}\).
\begin{proof}
Let \(R\) be a ring, \(J\) a nilpotent ideal and \(P\) an invertible
\(R\)-module.  Lift the prescribed generator of \(P/JP\) to \(p\in P\).
At every prime of \(R\), choose a basis of the rank-one free module
\(P_\mathfrak p\).  The coefficient of \(p\) in this basis is a
unit modulo \(J_\mathfrak p\), hence is a unit, since \(J\) is
nilpotent.  Thus \(R\to P\), \(1\mapsto p\), is an isomorphism
at every prime and therefore an isomorphism.  This proves the lifting
statement without a local-ring assumption.

Write the cube-defect line as
\[
 B_L(x,z)=L_{x+z}\otimes L_x^{-1}\otimes L_z^{-1}\otimes L_0.
                                                               \tag{CNC.5}
\]
The field \(\mathrm{SAB}.\mathrm{prim}\) is the actual isomorphism
\(L\simeq B_L(\mathrm{id},[w_0])\).  Equivalently it gives a
global generator of
\(L\otimes B_L(\mathrm{id},[w_0])^{-1}\).
At zero the second factor has its canonical trivialization.  The ratio
of this generator with the given rigidification of \(L_0\) is a
unit on the original base.  Multiply by the inverse pullback of that
unit.  The resulting generator \(w_P\) has exactly the zero-section
value required by \(\mathrm{IsPrim}\).

Similarly \(\mathrm{SAB}.\mathrm{adj}(a)\) supplies an isomorphism
between \(B_L([a]x,z)\) and \(B_L(x,[a^*]z)\).  Choose its
generator \(w_a\) in
\(B_L(x,[a^*]z)\otimes B_L([a]x,z)^{-1}\), and normalize its
value at \((0,0)\) to one.  This is possible by a base unit and
unique by CNC1.  All these choices are made for the original line
bundle on the original family.  In particular the primitive identity
is not asserted for an arbitrary symmetric line bundle lacking the
original \(\mathrm{prim}\) field.
\end{proof}
% END CNC4

% BEGIN CNC5
\paragraph{CNC5. Compatible normalized frames and their four formal ratios.}
\label{claim:CNC5}
For the chart of CNC3 there are compatible generators \(e_m\) of
the pullbacks of \(B_L\) to \(J_m(S,2)\), equal to the canonical
generator on either zero axis.  They define formal units
\(P(x,z,w),Q(x,z,w),s(x,z),a_b(x,z)\) which give all the original
\(\mathrm{IsFrame}\) equations, with \(P=\mathrm{ph1}\),
\(Q=\mathrm{ph2}\), \(s=\mathrm{sy}\), and
\(a_b=\mathrm{ad}(b)\).
\begin{proof}
Pull \(L\) back along the formal sum of the two universal points.
Starting from its rigidification at zero, repeatedly apply CNC4 to the
surjective jet truncations.  This gives compatible generators \(g_m\)
of this line at every order.  Divide by their restrictions to the two
axes and multiply by their value at their intersection.  The result is
the compatible generator
\[
 e_m(x,z)=g_m(x,z)\,g_m(x,0)^{-1}g_m(0,z)^{-1}g_m(0,0)
                                                               \tag{CNC.6}
\]
of (CNC.5).  The formula is an equality of line tensors, so it is
independent of local trivializations.  It gives one on both axes.

Let \(c(x,z,w)\) denote the generator of the triple cube defect
specified by the original cubical structure.  At each jet order take
the following ratios of generators of the same line:
\[
\begin{aligned}
 P(x,z,w)&=c(x,z,w)\,
              \frac{e(x,w)e(z,w)}{e(F(x,z),w)},\\
 Q(x,z,w)&=c(z,w,x)\,
              \frac{e(x,z)e(x,w)}{e(x,F(z,w))},\\
 s(x,z)&=\frac{e(z,x)}{e(x,z)},\\
 a_b(x,z)&=w_b(x,z)\,
                    \frac{e([b]x,z)}{e(x,[b^*]z)}.
\end{aligned}                                                   \tag{CNC.7}
\]
Here and below a fraction is a tensor product with inverse generators;
each whole expression is an actual scalar unit.  Ratios commute with
truncation, so these units define unique formal series over \(S\).
Their constant terms are one, by the axes of \(e\), the normalized
cube, and CNC4.  Conversely their truncations are exactly the units in
the source \(\mathrm{IsFrame}\) equations: multiplying (CNC.7)
by its displayed denominator recovers, in order, its first partial
law, second partial law with the permutation \((z,w,x)\), symmetry
law, and adjoint law.  This also checks the direction of the adjoint
ratio and the inverse in the source symmetry equation.
\end{proof}
% END CNC5

% BEGIN CNC6
\paragraph{CNC6. All formal biextension identities on the original frames.}
\label{claim:CNC6}
The series from CNC5, together with \(F\) and \([b]\), are a
literal object \(B:\mathrm{FB}(S,d,O)\).
\begin{proof}
For brevity in this proof only, write \(x+z=F(x,z)\) and
\(b x=[b](x)\).  The ten zero-axis conditions are
\[
 P(0,z,w)=P(x,0,w)=P(x,z,0)=1,
 \quad Q(0,z,w)=Q(x,0,w)=Q(x,z,0)=1,
\]
\[
 s(0,z)=s(x,0)=a_b(0,z)=a_b(x,0)=1.                            \tag{CNC.8}
\]
The remaining identities are the following; this list includes every
non-axis field of the original \(\mathrm{FB}\):
\[
\begin{aligned}
 P(x,z,t)P(x+z,w,t)&=P(z,w,t)P(x,z+w,t),\\
 P(z,x,w)&=P(x,z,w),\\
 Q(x,z,w)Q(x,z+w,t)&=Q(x,w,t)Q(x,z,w+t),\\
 Q(x,w,z)&=Q(x,z,w),\\
 P(x,z,w)P(x,z,t)Q(x+z,w,t)
   &=Q(x,w,t)Q(z,w,t)P(x,z,w+t),\\
 s(x,z)s(z,x)&=1,\\
 Q(x,z,w)s(x,z)s(x,w)&=s(x,z+w)P(z,w,x).
\end{aligned}                                                   \tag{CNC.9}
\]
For every \(b,c\in O\), the adjoint identities are
\[
\begin{aligned}
 P(bx,bz,w)a_b(x+z,w)
   &=a_b(x,w)a_b(z,w)P(x,z,b^*w),\\
 Q(bx,z,w)a_b(x,z+w)
   &=a_b(x,z)a_b(x,w)Q(x,b^*z,b^*w),\\
 a_1(x,z)&=1,\\
 a_{bc}(x,z)&=a_b(cx,z)a_c(x,b^*z),\\
 P(bx,cx,z)a_{b+c}(x,z)
   &=a_b(x,z)a_c(x,z)Q(x,b^*z,c^*z),\\
 a_b(x,z)a_{b^*}(z,x)s(bx,z)&=s(x,b^*z).
\end{aligned}                                                   \tag{CNC.10}
\]

We verify these as line-tensor identities before reading their
coefficients.  The cubical structure gives a symmetric normalized
trivialization of the triple cube defect.  Its cocycle equation is the
agreement of the two ways to split a sum of three points; its
interchange equation is the agreement of splitting in the first and
second variables.  Inserting (CNC.7), the \(e\)-factors cancel from
the first, third and fifth equations of (CNC.9), leaving respectively
these two cocycle equations and the interchange equation for the cube.
The second and fourth use its symmetry.  The sixth cancels immediately.
In the seventh both sides reduce to the same cube at \((z,w,x)\)
and the same ratio
\(e(z,x)e(w,x)/e(x,z+w)\).  The faces of the cube and (CNC.6)
give the first eight axis equations.

For completeness, the normalized adjoint isomorphisms used here have
all the compatibility diagrams which these cancellations require.
The two ways to combine the adjoint isomorphism with the first or
second cubical law are isomorphisms of the same line bundles on
\(A^3\).  The two composites for a product \(bc\), for a sum
\(b+c\), and for transposing the variables are isomorphisms of
the same line bundles on \(A^2\).  They agree at the zero tuple:
every cube face is one and every \(w_b\) was normalized there.
Their ratios are therefore one by CNC1.  The same argument on \(A\)
shows that either zero-axis restriction of \(w_b\) is one.  The
isomorphism for \(1\) is likewise the identity.  Substitution of
(CNC.7) into these six diagrams gives exactly (CNC.10), including its
factors \(P(bx,cx,z)\) and \(Q(x,b^*z,c^*z)\).
This proves the last two equations of (CNC.8) as well.

All line identifications just used are between tensor expressions in
the same \(L\), with the indicated group and \(O\)-action maps;
associativity, commutativity, and \([bc]=[b][c]\) identify their
source and target.  Thus no biextension compatibility of an arbitrarily
chosen adjoint isomorphism was presumed: it follows from its
normalization and CNC1.  Pull the identities to each universal jet.
Equality of every truncation is equality of formal series, so
(CNC.8)--(CNC.10) prove the asserted \(\mathrm{FB}\) structure
over the entire \(S\).
\end{proof}
% END CNC6

% BEGIN CNC7
\paragraph{CNC7. Formal logarithms over a general rational algebra.}
\label{claim:CNC7}
Let \(R\) be a commutative \(\mathbb Q\)-algebra and \(F\) a
commutative \(d\)-dimensional formal group law over \(R\).
There is a unique formal coordinate change \(\ell(x)=x+O(x^2)\)
with
\(\ell(F(x,z))=\ell(x)+\ell(z)\).  In these coordinates every
formal \(O\)-action is linear.
\begin{proof}
Set \(J(x)=\partial_zF(x,0)\), an invertible matrix with
\(J(0)=1\).  Associativity, differentiated in the third variable,
gives
\[
 \partial_zF(x,z)J(z)=J(F(x,z)).                              \tag{CNC.11}
\]
The vector fields
\(X_i=\sum_jJ_{ji}(x)\partial_{x_j}\) commute.  Indeed
commutativity and associativity give
\(F(F(x,se_i),te_j)=F(F(x,te_j),se_i)\); the coefficient of
\(st\) is exactly \([X_i,X_j]=0\).
The rows of the one-form \(\omega=J(x)^{-1}dx\) are dual to
these fields.  For each row,
\(d\omega(X_i,X_j)=X_i(\omega(X_j))-X_j(\omega(X_i))
-\omega([X_i,X_j])=0\).  Since the \(X_i\) form a basis of
derivations, \(d\omega=0\).

Every closed formal one-form over a rational algebra has its radial
primitive with zero constant term: for a row \(\sum_i b_i(x)dx_i\),
it is
\(\int_0^1\sum_i x_i b_i(tx)\,dt\).
Here integration means dividing a coefficient of \(t^m\) by
\(m+1\).  Differentiation and the equality
\(\partial_jb_i=\partial_ib_j\) verify the primitive directly.
Apply this construction to \(\omega\) to get \(d\ell=\omega\),
\(\ell(0)=0\), and linear part the identity.
Equation (CNC.11) shows that the derivative in \(z\) of
\(\ell(F(x,z))-\ell(z)\) is zero.  A formal series over a
\(\mathbb Q\)-algebra with all first derivatives zero is constant,
coefficient by coefficient.  Evaluation at \(z=0\) gives the desired
addition identity.

A coordinate series with invertible linear part has a unique inverse,
obtained one homogeneous degree at a time.  Any two logarithms differ
by an additive coordinate series with linear part the identity.
For an additive series \(q\), differentiating
\(q(x+z)=q(x)+q(z)\) at \(z=0\) shows that its derivative is
constant.  It is therefore linear.  This proves uniqueness, and the
same argument applies to
\(\ell\circ[b]\circ\ell^{-1}\), giving a matrix \(M(b)\).
The formal action laws show that \(b\mapsto M(b)\) is a unital
ring homomorphism to \(\operatorname{Mat}_d(R)\).
\end{proof}
% END CNC7

% BEGIN CNC8
\paragraph{CNC8. A normalized symmetric additive two-cocycle is a coboundary.}
\label{claim:CNC8}
Suppose \(c(x,z)\in R[[x,z]]\), with \(R\) a commutative
\(\mathbb Q\)-algebra, satisfies
\[
 c(x,0)=c(0,z)=0,\quad c(x,z)=c(z,x),\quad
 c(x,z)+c(x+z,w)=c(z,w)+c(x,z+w).
\]
Then a series \(h(x)\), with zero constant and linear parts, satisfies
\(c(x,z)=h(x)+h(z)-h(x+z)\).
The construction is linear in the coefficients of \(c\), so it also
works with formal parameter variables.
\begin{proof}
Put \(b_i(x)=\partial_{z_i}c(x,0)\).  Differentiating the cocycle
identity in \(w_i\) at zero gives
\[
 \partial_{z_i}c(x,z)=b_i(x+z)-b_i(z).                        \tag{CNC.12}
\]
By symmetry,
\(\partial_{x_j}c(x,z)=b_j(x+z)-b_j(x)\).
Equality of the two mixed derivatives therefore gives
\(\partial_jb_i=\partial_ib_j\).  Also \(b_i(0)=0\).
Define
\[
 h(x)=-\int_0^1\sum_i x_i b_i(tx)\,dt.                       \tag{CNC.13}
\]
The same elementary radial calculation as in CNC7 gives
\(\partial_i h=-b_i\).  Thus the derivative in each \(z_i\)
of \(h(x)+h(z)-h(x+z)\) is the right side of (CNC.12).
The difference from \(c\) is independent of \(z\) and vanishes
at \(z=0\); it is zero.  The origin and linear terms of \(h\)
vanish because the \(b_i\) have zero constant term.  Formula
(CNC.13) uses only derivatives, rational divisions and linear
operations.  It respects any linear conditions in other parameter
variables, including vanishing and linear dependence on those
variables.  For parameters we apply it coefficient by coefficient;
each resulting bidegree depends on finitely many coefficients.
\end{proof}
% END CNC8

% BEGIN CNC9
\paragraph{CNC9. The precise gauge change.}
\label{claim:CNC9}
For any formal \(\mathrm{FB}\) over a \(\mathbb Q\)-algebra
and any \(h(x,z)\) vanishing on both axes, define
\[
\begin{aligned}
 \delta_1h(x,z,w)&=h(x,w)+h(z,w)-h(F(x,z),w),\\
 \delta_2h(x,z,w)&=h(x,z)+h(x,w)-h(x,F(z,w)),\\
 \delta_sh(x,z)&=h(z,x)-h(x,z),\\
 \delta_bh(x,z)&=h([b]x,z)-h(x,[b^*]z).
\end{aligned}                                                   \tag{CNC.14}
\]
Replacing the four units \((P,Q,s,a_b)\) by
\[
 (P\exp(-\delta_1h),\ Q\exp(-\delta_2h),\
       s\exp(-\delta_sh),\ a_b\exp(-\delta_bh))               \tag{CNC.15}
\]
again gives a formal \(\mathrm{FB}\), with the same group law and
action.  Successive gauge changes add their \(h\)'s.
\begin{proof}
Each series in (CNC.14) has zero constant term, so its exponential is
defined.  Formal exponential is multiplicative on sums in a commutative
rational algebra, as follows from the binomial formula at each finite
degree.  Substitute (CNC.15) into (CNC.8)--(CNC.10).  The old unit
identity cancels, leaving an additive identity in \(h\).
For the first partial associativity identity this is
\[
 \delta_1h(x,z,t)+\delta_1h(F(x,z),w,t)
 =\delta_1h(z,w,t)+\delta_1h(x,F(z,w),t),
\]
where all terms cancel by associativity of \(F\).  The other
partial associativity is the same calculation in the second variable;
the two commutativities and interchange follow by permuting or
regrouping the four displayed summands.  The symmetry identities use
\(\delta_sh(z,x)=-\delta_sh(x,z)\).

The adjoint compatibility equations follow by applying
\([b]F=F([b],[b])\) in the two variables.  For the product equation
the needed cancellation is
\[
 \delta_{bc}h(x,z)=\delta_bh([c]x,z)
                         +\delta_ch(x,[b^*]z).
\]
For the sum equation use
\([b+c]x=F([b]x,[c]x)\) and the corresponding identity for the
involution.  Finally
\(\delta_bh(x,z)+\delta_{b^*}h(z,x)+\delta_sh([b]x,z)
=\delta_sh(x,[b^*]z)\), which is immediate on expanding four
differences.  This verifies every adjoint law, and the axis laws follow
from those of \(h\).  Addition of gauges follows from linearity of
all four \(\delta\)'s and multiplicativity of exponential.
\end{proof}
% END CNC9

% BEGIN CNC10
\paragraph{CNC10. Killing the two partial cocycles.}
\label{claim:CNC10}
For every formal \(\mathrm{FB}\) over a commutative
\(\mathbb Q\)-algebra there is an axis-zero gauge \(h_{12}\)
after which \(P=Q=1\).
\begin{proof}
Pass through CNC7 to additive coordinates.  Formal logarithm is inverse
to exponential between the ideal of series with zero constant term and
units with constant term one.  This follows by the usual one-variable
formal identities, whose substitutions are valid in the positive-degree
ideal.  Taking logarithms turns all unit identities into additive
identities over the same coefficient algebra.

Write \(p=\log P\).  For each formal parameter \(w\), the
first two identities in (CNC.9) say that \(p(x,z,w)\) is a
normalized symmetric two-cocycle in \(x,z\).  CNC8 constructs
\(h_1(x,w)\) with \(\delta_1h_1=p\).  It has zero first-variable
origin, and vanishes at \(w=0\) because \(p\) does.  Gauge by
\(h_1\).  The new \(P\) is one, and all identities remain true
by CNC9.

Let \(q=\log Q\) for the new system.  Interchange in (CNC.9)
now says
\(q(x+z,w,t)=q(x,w,t)+q(z,w,t)\).
Differentiating in \(z\) at zero shows that \(q\) is linear
in \(x\).  The remaining partial associativity and symmetry say
that it is a normalized symmetric two-cocycle in its last two variables.
Apply CNC8 there, coefficient by coefficient in the first variable,
to obtain \(h_2(x,z)\) with \(\delta_2h_2=q\).
It is linear in \(x\) and vanishes on both axes, so
\(\delta_1h_2=0\).  A second gauge kills \(Q\) while preserving
\(P=1\).  Thus \(h_{12}=h_1+h_2\) has the required property
in additive coordinates.  Substituting the formal logarithm back into
both variables gives the assertion in the original coordinates.
\end{proof}
% END CNC10

% BEGIN CNC11
\paragraph{CNC11. Killing every adjoint cocycle by one separability calculation.}
\label{claim:CNC11}
Suppose that \(O\) is a commutative reduced ring finite and free over
\(\mathbb Z\), with an involution.  After CNC10, one bilinear gauge
kills \(a_b\) simultaneously for every \(b\in O\), and preserves
\(P=Q=1\).
\begin{proof}
Remain in additive coordinates and write the action as the matrices
\(M(b)\) of CNC7.  The first two equations of (CNC.10), with
\(P=Q=1\), say that \(\log a_b(x,z)\) is additive in each
variable.  It is therefore bilinear, so write it as \(x^TD(b)z\).
The remaining equations give
\[
 D(b+c)=D(b)+D(c),\quad D(1)=0,\quad
 D(bc)=M(c)^TD(b)+D(c)M(b^*).                                 \tag{CNC.16}
\]
Commutativity of \(O\) permits exchanging \(b,c\) in the last
equation.  With
\(\mathcal L(b)Z=M(b)^TZ\) and
\(\mathcal R(b)Z=ZM(b^*)\), it is consequently the ordinary
bimodule derivation identity
\(D(bc)=\mathcal L(b)D(c)+D(b)\mathcal R(c)\).
The two actions commute.  The maps extend \(\mathbb Q\)-linearly
to \(B=O\otimes\mathbb Q\); this is allowed because the target
coefficient algebra is rational.  The identity extends as well.

The algebra \(B\) is finite and reduced over \(\mathbb Q\), hence
a product of finite separable fields.  Let
\(e=\sum_i u_i\otimes v_i\in B\otimes_{\mathbb Q}B\)
be its separability element.  Explicitly it is the sum of a trace-dual
basis tensor in each field factor.  Its two identities are
\[
 \sum_i u_iv_i=1,\qquad
 (b\otimes1)e=e(1\otimes b).                                 \tag{CNC.17}
\]
One can check them after a splitting-field extension: under
\(B\otimes\overline{\mathbb Q}\simeq
\overline{\mathbb Q}^{\operatorname{Hom}(B,\overline{\mathbb Q})}\),
the tensor is the diagonal idempotent, so both identities are coordinate
identities.  Faithful scalar extension brings them back to \(B\).
This also covers a zero algebra by the empty sum on its zero module.

Set \(Z=\sum_i\mathcal L(u_i)D(v_i)\).  Applying the second
identity of (CNC.17) to the linear map
\(u\otimes v\mapsto\mathcal L(u)D(v)\), and then using the
derivation identity, gives
\[
 \mathcal L(b)Z
 =\sum_i\mathcal L(u_i)D(v_ib)
 =D(b)+Z\mathcal R(b).
\]
Thus \(D(b)=M(b)^TZ-ZM(b^*)\).
The bilinear series \(h_3(x,z)=x^TZz\) has
\(\delta_bh_3=\log a_b\) for every \(b\), while its two
partial coboundaries are zero.  CNC9 therefore kills every adjoint
unit in a single gauge change and preserves the already trivial
partial cocycles.  The separability calculation takes place over
\(\mathbb Q\); it does not require the possibly nonreduced
coefficient algebra to be a field or an etale algebra.
\end{proof}
% END CNC11

% BEGIN CNC12
\paragraph{CNC12. The full original canonical-gauge theorem.}
\label{claim:CNC12}
For every commutative ring \(R_0\), every \(d\), every commutative
reduced finite free \(\mathbb Z\)-algebra \(O\) with involution,
every commutative \(\mathbb Q\)-algebra \(R\) which is also an
\(R_0\)-algebra, and every \(B:\mathrm{FB}(R_0,d,O)\), there is
\(h\in R[[x,z]]\) with
\[
 h(0,z)=h(x,0)=0,\qquad
 (P,Q,s,a_b)_R
   =(\exp\delta_1h,\exp\delta_2h,
      \exp\delta_sh,\exp\delta_bh)\quad(b\in O).              \tag{CNC.18}
\]
This is the complete original \(\mathrm{CanonGaugeSt}\), hence in
particular supplies \(B.\mathrm{IsCanonGauge}(S,h)\) for CNC6.
\begin{proof}
Map the coefficients of \(B\) to \(R\); all its equations persist.
Use CNC10 and CNC11 in additive coordinates.  At this stage
\(P=Q=a_b=1\) for every \(b\).  Put \(t=\log s\).
The last two symmetry equations of (CNC.9) give
\[
 t(x,z)=-t(z,x),\qquad t(x,z+w)=t(x,z)+t(x,w).
\]
They imply that \(t\) is bilinear.  The last equation of (CNC.10)
now gives \(t(M(b)x,z)=t(x,M(b^*)z)\).
Take \(h_4=-t/2\).  Then
\(\delta_sh_4=t\); its two partial coboundaries vanish by
bilinearity, and every \(\delta_bh_4\) vanishes by the displayed
adjoint relation.  This last gauge makes all four units equal to one.

Add the four gauges of CNC10--CNC11 and this last step, and substitute
\(\ell(x),\ell(z)\) into the resulting additive-coordinate
series.  CNC9 shows that the original units are exactly the four
exponentials in (CNC.18), with the original formal group law and
original action.  This is precisely the six fields \(\mathrm{ax0}\),
\(\mathrm{ax1}\), \(\mathrm{ph1}\), \(\mathrm{ph2}\),
\(\mathrm{sy}\), and \(\mathrm{ad}\) of
\(\mathrm{IsCanonGauge}\), including the sign of the symmetry
difference and the involution in the adjoint difference.  Since \(O_E\)
is a reduced finite free \(\mathbb Z\)-algebra, CNC6 is an allowed
instance.  Every construction is valid in all finite jet orders over
the same \(R\), so no reducedness or continuity condition has entered.
\end{proof}
% END CNC12

% BEGIN CNC13
\paragraph{CNC13. The actual line generator and the unique germ of every section.}
\label{claim:CNC13}
Let \(w_P\) and \(e_m\) be the generators of CNC4--CNC5.
For every \(k\geq0\) and every original section
\(F\in H^0(A_S,L^{\otimes k})\) there is a unique series
\(f\in S[[x_1,\ldots,x_d]]\) satisfying the original
\(\mathrm{IsGerm}\).  Its coefficient at every jet order is
evaluation of the given section in the given line generator.
\begin{proof}
Define the generator of \(L\) at the formal point by
\[
 e_{L,m}(x)=w_P(x)\,e_m(x,[w_0]x).                           \tag{CNC.19}
\]
The factors belong to \(L_x\otimes B_L(x,[w_0]x)^{-1}\)
and \(B_L(x,[w_0]x)\), respectively, so their product belongs
to the actual \(L_x\).  Their generator properties show that it is
a generator.  It commutes with truncation, by CNC3 and CNC5.
At the zero jet it equals the original rigidification, by the
normalization of \(w_P\) and the axes of \(e_m\).

Pull the section \(F\) to the actual family over each jet algebra,
then evaluate it at \(x_m\) in the basis \(e_{L,m}^{\otimes k}\).
This gives a unique element \(f_m\in J_m(S,1)\).
Compatibility of pullback, evaluation and (CNC.19) gives
\(f_{m+1}\bmod(t)^{m+1}=f_m\).
The inverse limit of these rings is the formal series ring over \(S\):
the coefficient of a monomial is its coefficient in any truncation
large enough to contain it.  Let \(f\) be this series.

The pullback section just used is the \(F'\) in the definition of
\(\mathrm{IsGerm}\), the tensor power of (CNC.19) is its
\(\tau\), and the equality \(f\bmod(t)^{m+1}=f_m\) is its
\(\mathrm{IsValueAt}\) clause.  Conversely any witnesses in that
definition have the same pullback section and the same generator under
its prescribed correspondence.  The scalar coefficient in a generator
is unique.  Thus every such series has all truncations \(f_m\),
which proves uniqueness.  This uses equality in \(S\), not equality
on its points, and therefore also proves the assertion for nonreduced
coefficient algebras.
\end{proof}
% END CNC13

% BEGIN CNC14
\paragraph{CNC14. A complete original chart at every framed point.}
\label{claim:CNC14}
For every commutative \(K\)-algebra \(S\), with its rational-algebra
structure, and every \(y:\mathrm{FrPt}(X_o,c,S)\), there exists
an original \(C:\mathrm{ChartAt}(X_o,c,y)\) with compatible frames
and a germ for every \(k\geq0\) and every original section.
\begin{proof}
Use the involution of (CNC.1).  CNC2--CNC3 supply
\(d,F,[b],u,x_m\) and all six chart fields.  CNC4 supplies the
global primitive generator \(w_P\) and its \(\mathrm{IsPrim}\)
normalization.  CNC5 supplies \(e_m\) and all four frame equations;
CNC6 packages their series as the original formal \(\mathrm{FB}\).
CNC12 supplies its gauge \(h\).  Formula (CNC.19) at zero is the
field \(\mathrm{rig0}\), by CNC13.  These are all the data and
all the proof fields in \(\mathrm{ChartAt}\).  Frame truncation,
although not separately demanded in that structure, is available from
CNC5 and is exactly \(\mathrm{FrTruncC}\) used by the source
germ producer.  Finally CNC13 supplies, for every natural \(k\) and
every section, the required germ.

The result thus matches the entire conclusion of
\texttt{Chart\allowbreak Ring.\allowbreak exists\_\allowbreak chart\allowbreak At\_\allowbreak with\_\allowbreak germs} on the original
moduli family.  The large source \(\mathrm{Laws}\) package is not
being taken as a new premise: of its geometric content this construction
needed only the base-functions theorem of CNC1, whose hypotheses were
checked there.  The smooth formal chart and all gauge equations have
been constructed above.
\end{proof}
% END CNC14

% BEGIN CNC15
\paragraph{CNC15. Construction of the literal original \(\mathrm{IsNab}\).}
\label{claim:CNC15}
For every original \(E,H,D_w,l,X_o,c\) there exist the star-ring
structure \(a^*=H.\mathrm{bar}(a)\) on \(O_E\) and
\(\nabla:\mathrm{NabV}(X_o,c)\) such that
\(\mathrm{IsNab}(X_o,c,\nabla)\).
This has the full coefficient-algebra and section quantifiers of
\texttt{Global\allowbreak Theory.\allowbreak Canon\allowbreak Conn.\allowbreak exists\_\allowbreak is\allowbreak Nab\_\allowbreak uncond}.
\begin{proof}
For every coefficient algebra and framed point, choose once a chart
\(C\) supplied by CNC14.  For its formal group law define the
invariant differential operator in direction \(v\in S^d\) by
\[
 D_v f=\sum_{j=1}^d\left(\sum_{i=1}^d
       v_i\,\partial_{z_i}F_j(x,0)\right)\partial_{x_j}f.
                                                               \tag{CNC.20}
\]
For a multi-index \(\alpha:I\to\mathbb N\) of finite support,
take the product of \(D_{u_\iota}^{\alpha(\iota)}\) in the
original finite-index order.  This is exactly
\(\mathrm{DpolyOf}(F,u,X^\alpha)\) in the source.  The operators
also commute, by the commuting invariant fields in CNC7, so the value
would be independent of the order.

For \(k\geq0\), let \(f\) be the unique germ of \(F\) from
CNC13 and set
\[
 \nabla_y(k,F)(\alpha)=
 \left[\prod_{\iota\in I}D_{u_\iota}^{\alpha(\iota)}
    \left(f(x)\exp\{k\,h(x,[w_0]x)\}\right)\right]_{x=0}.
                                                               \tag{CNC.21}
\]
This is an element of \(S\); every requested coefficient involves
only finitely many coefficients of \(f,F,h\).
For a negative integer \(k\), set the corresponding value to zero.
This defines the function on every integer required by
\(\mathrm{NabV}\); \(\mathrm{IsNab}\) constrains only its
nonnegative values.

A \(K\)-algebra has its canonical rational-algebra structure, since
\(K\) is a number field.  Any other rational-algebra structure appearing
in the \(\mathrm{IsNab}\) quantifier has the same structure map:
a unital ring map from \(\mathbb Q\) is determined by the images
of integers and their inverses.  Thus the choice of chart and the value
in (CNC.21) can be made for the \(\mathrm{NabV}\) data alone,
and work for every such additionally specified rational structure.
The star equation is true by construction.  For every \(S,y\), the
chosen chart has all the required germs by CNC14.  If \(f'\) is
\emph{any} series satisfying its \(\mathrm{IsGerm}\), CNC13
gives \(f'=f\).  Formula (CNC.21) is then exactly
\[
 C.B.\mathrm{nabVal}(S,C.h,H.w_0,k,f',C.uvec,X^\alpha),
\]
because the source \(\mathrm{hL}\) is \(h(x,[w_0]x)\) and
its \(\mathrm{nabVal}\) is the constant coefficient after
(CNC.20).  This proves the equality for every germ, not merely a selected
representative.  It is the final field of \(\mathrm{IsNab}\),
so all its quantifiers and fields have been verified.
\end{proof}
% END CNC15

% BEGIN CNC16
\paragraph{CNC16. Exact removal of the connection premise from the lattice route.}
\label{claim:CNC16}
On every supplied original geometric family to which ICL11--ICL21
apply, CNC15 supplies their original \(\mathrm{IsNab}\) producer
at every row and for both original Boolean level choices.  The
conditional lattice-spanning implications of that chapter consequently
require no separate connection-existence premise.
\begin{proof}
CNC15 quantifies over every \(D_w,l,X_o,c\), without a rank,
weight, prime, or embedding exclusion.  Apply it to each row's actual
\(Q.X_o(j)\), with that row's original coefficient field and level.
Choosing these connections gives exactly the data assumed at the start
of ICL11; its identity cover and all subsequent displayed maps use those
same objects.  If two Boolean choices use different original levels,
apply the same universally quantified theorem to each.  No change in
the \(\mathrm{FWTie}\) maps of ICL21 is needed.

This supplies one producer in that conditional route.  CNC15 asserts
neither construction of the common moduli family nor integral
valuation bounds for (CNC.21).  In particular the principal
same-Gram charts, actual finite theta system, Gaussian identity and
absolute integrability, global short-rule dictionary, and ramified
short-value estimates retained by ICL16--ICL21 are still required.
The higher-cover and all-tuple class-sum obligations are not consequences
of \(\mathrm{IsNab}\) alone.  Keeping those hypotheses gives exactly
the stated implication and makes no additional whole-obligation closure.
\end{proof}
% END CNC16


\section{The original principal cusp lattice and its exact normalization}
\label{sec:principal-cusp-lattices-v047}

\noindent\textbf{Scope.}
This chapter constructs the Hermitian coordinates and the self-dual
lattice that the principal-chart part of ICL retains.  Its generic
endpoint is the full original \(\mathrm{DataX}\) record, on the same
\(\mathrm{Data}\), with the same Hermitian matrix, scalar, arithmetic
group and principal modulus.  On a supplied original uniformization
it gives the exact host denominators and their logarithmic bounds.
Existence of that uniformization and of its principal congruence
inclusion remains an explicit hypothesis.  In particular this chapter
does not assert \(\mathrm{exists\_\allowbreak unf\_\allowbreak hasPrinc}\), construct
\(\mathrm{LawsCore}\), or complete the five remaining ICL spanning
producer groups.

Write \(O=\mathcal O_E\), where \(E\) is a CM number field, and put
\(\bar a=\operatorname{cc}(a)\).  The original convention is
\[
 h_G(x,y)=x^tG\bar y,\qquad
 B_{\xi,G}(x,y)=\operatorname{Tr}_{E/\mathbb Q}
                                      (-\xi h_G(x,y)).       \tag{PCL.1}
\]
Thus \(h_G\) is linear in its first variable.  We write
\(\mathfrak D_E\) for the absolute different and \(aL\) for the
\(O\)-span of products \(\alpha v\), \(\alpha\in a,v\in L\),
where \(a\) is a fractional ideal.  A full lattice is a finitely
generated \(O\)-submodule with the whole \(E\)-span; it need not
have an \(O\)-basis.  The algebraic foundations used here are the
Dedekind ideal laws, the nondegenerate trace pairing for a number
field, and the different as the inverse trace-dual ideal.  Their
specific consequences are derived below.  No local freeness or
principal-ideal hypothesis is imposed on the actual lattices.

% BEGIN PCL1
\paragraph{PCL1. Hermitian coordinates for the original alternating datum.}
\label{claim:PCL1}
Let \(D_w:\mathrm{Datum}(H)\) be the original datum, let
\(r=D_w.p+D_w.q\), and let \(0\ne\xi\in E\) satisfy
\(\bar\xi=-\xi\).  There exist a matrix \(G\in M_r(E)\)
and an \(O\)-linear injection \(\phi:D_w.\mathrm{La}\to E^r\)
such that its image spans \(E^r\), \(\bar G=G^t\),
\(\det G\ne0\), and
\[
 D_w.\mathrm{EF}(x,y)
       =\operatorname{Tr}_{E/\mathbb Q}
                            (\xi h_G(\phi x,\phi y)).       \tag{PCL.2}
\]
This is the literal \(\mathrm{exists\_\allowbreak herm\_\allowbreak coord\_\allowbreak of\_\allowbreak datum}\),
including its sign and its original lattice domain.

\emph{Proof.}
Put \(L=D_w.\mathrm{La}\) and \(b=D_w.\mathrm{EF}\).
The supplied datum says that \(b\) is a perfect alternating
\(\mathbb Z\)-pairing and
\(b(ax,y)=b(x,\bar a y)\) for \(a\in O\).
Perfection makes \(L\) torsion-free over \(\mathbb Z\).
It is also torsion-free over \(O\): a nonzero \(a\in O\)
divides a nonzero rational integer in \(O\), by its integral
equation, so \(ax=0\) would give integer torsion.  Hence
\(L\to V=E\otimes_O L\) is injective.  The original finite
generation and rank fields give \(\dim_E V=r\).

For fixed \(x,y\in L\), the additive functional
\(a\mapsto b(ax,y)\) on \(O\) extends uniquely to a
\(\mathbb Q\)-linear functional on \(E\).  Nondegeneracy
of the trace gives a unique \(t(x,y)\in E\) with
\[
 \operatorname{Tr}(a\,t(x,y))=b(ax,y)\quad(a\in O).
                                                               \tag{PCL.3}
\]
Uniqueness in (PCL.3), additivity, and the adjoint law show
\(t(ax,y)=a\,t(x,y)\),
\(t(x,ay)=\bar a\,t(x,y)\), and additivity in both variables.
Alternation and invariance of the trace under conjugation give
\(t(y,x)=-\overline{t(x,y)}\).  For example, pairing both
sides with every \(a\in O\) reduces this last assertion to
\(b(ay,x)=-b(x,ay)=-b(\bar a x,y)\).
These laws extend \(t\) uniquely to a sesquilinear form on
\(V\), by clearing denominators in each argument.  Then
\(h=\xi^{-1}t\) is Hermitian.  If \(h(v,V)=0\), choose
\(0\ne a\in O\) with \(av\in L\).  Equation (PCL.3) and
perfection give \(av=0\), hence \(v=0\).  Thus \(h\) is
nondegenerate.

Choose an \(E\)-basis of \(V\), take its coordinate map
\(V\simeq E^r\), and let \(G\) be the Gram matrix of \(h\).
Its restriction is \(\phi\); Hermitian symmetry and
nondegeneracy give the matrix assertions.  Taking \(a=1\)
in (PCL.3) proves (PCL.2).  This also covers \(r=0\), with
the empty matrix of determinant one.  For the chart convention
(PCL.1), the input to this theorem is \(-\xi_0\).
\hfill\(\square\)
% END PCL1

% BEGIN PCL2
\paragraph{PCL2. Ideal-valued duals without a principal lattice basis.}
\label{claim:PCL2}
For an \(O\)-lattice \(L\subset E^r\), put
\[
 L^\#=\{x:h_G(x,L)\subset O\}.
\]
For every nonzero fractional ideal \(a\),
\[
 \{x:h_G(x,L)\subset a\}=aL^\#,
 \qquad (aL)^\#=\bar a^{-1}L^\#.                    \tag{PCL.4}
\]
Moreover \(aL\) is a full finitely generated lattice,
\(a(a^{-1}L)=L\), and the product notation is exactly the
span of the products in the original source.

\emph{Proof.}
Since \(aa^{-1}=O\), choose finitely many
\(\alpha_i\in a,\beta_i\in a^{-1}\) with
\(\sum_i\alpha_i\beta_i=1\).  If \(h_G(x,L)\subset a\),
then \(h_G(\beta_i x,L)\subset O\), so
\(x=\sum_i\alpha_i(\beta_i x)\in aL^\#\).
The reverse inclusion follows term by term from linearity.
Next, \(x\in(aL)^\#\) is equivalent to
\(\bar\alpha h_G(x,y)\in O\) for all
\(\alpha\in a,y\in L\), which is equivalent to
\(h_G(x,L)\subset\bar a^{-1}\).  This proves (PCL.4).

Products of finite generating sets generate \(aL\), and any
nonzero \(\alpha\in a\) shows that its \(E\)-span contains
that of \(L\).  Expanding finite sums proves
\((ab)L=a(bL)\).  Together with \(aa^{-1}=O\), this gives
the cancellation identity and the source's displayed span
formula.  All these arguments use invertible ideals, not
generators for them. \hfill\(\square\)
% END PCL2

% BEGIN PCL3
\paragraph{PCL3. The exact different in the perfect trace pairing.}
\label{claim:PCL3}
Assume \(\xi\ne0\) and \(\det G\ne0\).  Define the
trace dual by
\(L^B=\{x:B_{\xi,G}(x,L)\subset\mathbb Z\}\).
Then
\[
 L^B=((\xi)\mathfrak D_E)^{-1}L^\# .              \tag{PCL.5}
\]
If \(L\) is full, the following conditions are equivalent:
the restriction of \(B_{\xi,G}\) is integral and represents
every additive functional \(L\to\mathbb Z\); \(L^B=L\);
and
\[
                         L^\#=(\xi)\mathfrak D_E L. \tag{PCL.6}
\]

\emph{Proof.}
The defining property of the inverse different is
\[
 z\in\mathfrak D_E^{-1}
       \ \Longleftrightarrow\
             \operatorname{Tr}(az)\in\mathbb Z\quad(a\in O).
\]
Because \(\bar a y\in L\) for \(a\in O,y\in L\), the
condition \(x\in L^B\) is equivalent to
\(-\xi h_G(x,y)\in\mathfrak D_E^{-1}\) for every
\(y\in L\).  The sign is a unit and does not change the
ideal.  PCL2 now gives (PCL.5).

The pairing \(B_{\xi,G}\) on the underlying rational vector
space is nondegenerate.  Indeed, if all its values on
\((x,y)\) vanish, also replace \(y\) by \(ay\) for every
\(a\in E\); trace nondegeneracy gives
\(\xi h_G(x,y)=0\), and then \(x=0\).
A full \(O\)-lattice is a full free \(\mathbb Z\)-lattice:
it is finitely generated and torsion-free, and its rational
span is all of \(E^r\).  The latter follows by replacing
an \(O\)-denominator by a nonzero integer multiple of it.
Every additive functional on \(L\) therefore extends to a
rational linear functional, represented uniquely by
\(B_{\xi,G}(x,-)\) for some \(x\in E^r\).  Its integer
values say exactly that \(x\in L^B\).

Consequently \(L^B=L\) gives both integrality and the stated
representation property.  Conversely integrality gives
\(L\subseteq L^B\); for \(x\in L^B\), the representation
property supplies \(x'\in L\) with the same pairing on
\(L\), and nondegeneracy on its rational span gives
\(x=x'\).  Finally cancel the nonzero fractional ideal in
(PCL.5), using PCL2, to obtain (PCL.6).
\hfill\(\square\)
% END PCL3

% BEGIN PCL4
\paragraph{PCL4. The actual fractional-ideal untwisting is self-dual.}
\label{claim:PCL4}
Let \(\Lambda\) carry the perfect integral trace pairing
of PCL3.  Suppose nonzero fractional ideals
\(\mathfrak A,\mathfrak A_b\) satisfy
\[
 \mathfrak A_b=\overline{\mathfrak A},\qquad
             (\xi)\mathfrak A\mathfrak A_b\mathfrak D_E=O.
                                                               \tag{PCL.7}
\]
Then the particular lattice
\[
               L'=\mathfrak A^{-1}\Lambda            \tag{PCL.8}
\]
is full, finitely generated, integral and self-dual for
\(h_G\), and \(\Lambda=\mathfrak A L'\).

\emph{Proof.}
PCL2 gives finite generation, fullness and the last
identity.  Equations (PCL.4), (PCL.6) and (PCL.7) give
\[
 (L')^\#=\overline{\mathfrak A}\,\Lambda^\#
       =\overline{\mathfrak A}(\xi)\mathfrak D_E\Lambda
       =\mathfrak A^{-1}\Lambda=L'.                  \tag{PCL.9}
\]
The inclusion \(L'\subseteq(L')^\#\) is precisely
integrality; its reverse is precisely the source's
\(\mathrm{L'\_\allowbreak dual}\) field.  Thus this is an equality
of the original lattices, with no saturation, denominator
enlargement or change of Hermitian form.
\hfill\(\square\)
% END PCL4

% BEGIN PCL5
\paragraph{PCL5. Construction of the complete original \(\mathrm{DataX}\).}
\label{claim:PCL5}
Let \(X:\mathrm{Data}(E,r)\), and assume exactly the inputs
of \(\mathrm{exists\_\allowbreak dataX\_\allowbreak of\_\allowbreak perfect\_\allowbreak traceForm}\):
(PCL.7) with \(\xi=X.\xi_0\), integrality and perfection
of \(B_{\xi,G}\) on \(X.\Lambda\), a real \(k\le1\)
with \(|\tau(\xi)|=k\) for every complex embedding, and
a positive integer \(g\) whose principal congruence
subgroup in the original unitary stabilizer lies in
\(X.\Gamma\).  There is \(X_x:\mathrm{DataX}(E,r)\) such
that
\[
 X_x.\mathrm{toData}=X,\quad X_x.g=g,\quad
 X_x.L'=\mathfrak A^{-1}X.\Lambda,\quad
 X_x.\Lambda=\mathfrak A X_x.L'.                     \tag{PCL.10}
\]

\emph{Proof.}
The ideal equation forces \(\mathfrak A\) and
\(\mathfrak A_b\) to be nonzero; their given
conjugate-membership condition says exactly that
\(\mathfrak A_b=\bar{\mathfrak A}\).
Every nonzero fractional ideal \(a\) contains a
positive rational integer.  Choose \(0\ne u\in a\),
multiply it by a nonzero integer so that \(u\in O\), and
use the finite additive group \(O/(u)\): its positive
order \(N((u))\) kills the class of \(1\), so
\(N((u))\in(u)\subset a\).  Choose in this way
\(N_A>0\) in \(\mathfrak A^{-1}\) and \(N_C>0\) in
\(\mathfrak A\).  Then
\[
                N_A\Lambda\subset L',\qquad
                            N_C L'\subset\Lambda.    \tag{PCL.11}
\]
There is a complex embedding of \(E\); its injectivity
and \(\xi\ne0\) give \(k>0\).  Set \(k_0=2k\).
It satisfies \(0<k_0\le2\) and the exact field
\(|\tau(\xi_0)|=k_0/2\) on the source CM type.

Keep the entire record \(X\).  Supply \(\mathrm{EF\_\allowbreak int}\)
and \(\mathrm{EF\_\allowbreak perf}\) by the two given trace-form
properties.  Supply \(L'\), its finite generation,
\(\mathrm{L'\_\allowbreak int}\), and \(\mathrm{L'\_\allowbreak dual}\) by PCL4.
Supply the two positive integers and both comparison
inclusions by (PCL.11).  Set its modulus to the given
\(g\); positivity and the \(\mathrm{cong}\) field are
exactly the assumed principal congruence inclusion.
These are every additional field of \(\mathrm{DataX}\).
The last equality in (PCL.10), expanded using PCL2, is
the literal span of \(\alpha u\) in the original endpoint.
In particular the construction changes no scalar,
matrix, character, group, domain or pre-existing metric.
\hfill\(\square\)
% END PCL5

% BEGIN PCL6
\paragraph{PCL6. Perfection on every component of the supplied uniformization.}
\label{claim:PCL6}
For any original \(U:\mathrm{Unf}(X_o,c,\sigma,r)\)
satisfying \(\mathrm{IsUnf}(U)\), each component
\(X=U.X(k)\) has the integral perfect trace pairing
in PCL5.  This proves the full
\(\mathrm{traceForm\_\allowbreak integral\_\allowbreak perfect\_\allowbreak of\_\allowbreak isUnf}\).

\emph{Proof.}
Unpack the actual \(\mathrm{IsUnf}\) witness.  It supplies
a family of original alternating data \(D_i\), a component
index \(i=i(k)\), and an \(O\)-linear equivalence
\(\phi_k:D_i.\mathrm{La}\simeq\Lambda_k\) such that
\[
        D_i.\mathrm{EF}(a,b)
             =B_{\xi_k,G_k}(\phi_k a,\phi_k b).       \tag{PCL.12}
\]
For two vectors in \(\Lambda_k\), apply \(\phi_k^{-1}\)
and use the integer-valued left side.  For an additive
\(\psi:\Lambda_k\to\mathbb Z\), perfection of
\(D_i.\mathrm{EF}\) represents \(\psi\circ\phi_k\)
by some \(a\in D_i.\mathrm{La}\).  Equation (PCL.12)
then represents \(\psi\) by \(\phi_k a\in\Lambda_k\).
This argument uses the lattice equivalence appearing in
the supplied uniformization; no independently chosen
lattice or rational isometry is substituted.
\hfill\(\square\)
% END PCL6

% BEGIN PCL7
\paragraph{PCL7. The two actual quadratic ideal identities.}
\label{claim:PCL7}
For the original host, let \(K=\mathbb Q(s)\),
\(s=\sqrt{-d}\), let \(C=\mathrm{ccK}\ne0\), and
let \(C_E=CO_E\).  The original prime and congruence
conditions give \(d>0\) squarefree and \(d\equiv3\pmod4\).
Then
\[
       \mathfrak D_K O_E=(s),\qquad
          C_E\overline{C_E}=(N_K C).                \tag{PCL.13}
\]
These are the literal \(\mathrm{DiffKPrincipal}\) and
\(\mathrm{CcNorm}\) suppliers on the original host.

\emph{Proof.}
The host conjugation is the CM conjugation used in
(PCL.1): the host's embedding law and the CM embedding
law both send it under \(\phi_0\) to complex conjugation,
so injectivity identifies them on all of \(E\).
The integer ring of \(K\) is
\(\mathbb Z[\omega]\), \(\omega=(1+s)/2\).
Here is the elementary integral-coordinate check.  If
\(x=a+bs\) is integral, its trace gives \(2a\in\mathbb Z\)
and its norm gives \(a^2+db^2\in\mathbb Z\).
For a rational denominator of \(b\) in lowest terms,
these conditions imply its square divides \(4d\);
squarefreeness of \(d\) makes that denominator divide
\(2\).  Writing \(x=(A+Bs)/2\), the norm condition is
\(A^2+dB^2\equiv0\pmod4\), so \(A\equiv B\pmod2\).
Conversely \(\omega\) satisfies the monic equation
\(T^2-T+(d+1)/4=0\).  This proves the stated integer ring.

Its trace dual is \(s^{-1}O_K\).  Indeed the trace matrix
of \(1,\omega\) has determinant \(-d\), so the trace dual
has index \(d\) over \(O_K\).  Directly,
\(\operatorname{Tr}(1/s)=0\),
\(\operatorname{Tr}(\omega/s)=1\), and
\(\operatorname{Tr}(\omega^2/s)=1\); hence
\(s^{-1}O_K\) lies in the trace dual.  It also has index
\(|N_{K/\mathbb Q}(s)|=d\), proving equality.
Inverting gives \(\mathfrak D_K=(s)\).

For the norm identity, factor \(C\) into prime ideals.
Over a split rational prime the two primes are exchanged
by conjugation and their product is \((p)\); each has
norm \(p\).  At an inert prime the single prime is
\((p)\), with norm \(p^2\), so its product with its
conjugate is \((p^2)\).  At a ramified prime its square
is \((p)\) and its norm is \(p\).  These three cases
give \(P\bar P=(NP)\) for every prime factor.  Products
give \(C\bar C=(N_K C)\).  Extending both identities
to \(E\) gives (PCL.13); host conjugation restricts to
the nontrivial conjugation of \(K\), since it sends
\(s\) to \(-s\). \hfill\(\square\)
% END PCL7

% BEGIN PCL8
\paragraph{PCL8. The original principal ideal and scalar normalization.}
\label{claim:PCL8}
Write \(\mathfrak b\) for the stored half-different, so
\(\mathfrak D_{E/K}=\mathfrak b\bar{\mathfrak b}\), and put
\[
 N=N_K C,\qquad
 \xi_H=-\frac1{Ns},\qquad
 \mathfrak A=C_E\mathfrak b^{-1},\qquad
 k_H=\frac1{N\sqrt d}.                              \tag{PCL.14}
\]
These are exactly the source's \(\xi_0,\mathrm{AA}\) and
\(\xi R\), with \(\mathrm{AAbar}=\bar{\mathfrak A}\).
They satisfy
\[
 (\xi_H)\mathfrak A\bar{\mathfrak A}\mathfrak D_E=O,
 \qquad \bar\xi_H=-\xi_H,
 \qquad |\tau(\xi_H)|=k_H\quad\hbox{for every }\tau,
 \qquad 0<k_H\le1.                                 \tag{PCL.15}
\]

\emph{Proof.}
The different in a tower gives
\(\mathfrak D_E=\mathfrak b\bar{\mathfrak b}(s)\).
Both half-different factors are nonzero because the
different is nonzero.  Apply PCL7 and cancel these
invertible ideals:
\[
 (\xi_H)(C_E\mathfrak b^{-1})
       (\bar C_E\bar{\mathfrak b}^{-1})
       \mathfrak b\bar{\mathfrak b}(s)
       =(\xi_H Ns)=(-1)=O.                          \tag{PCL.16}
\]
This is a proof of the actual \(\mathrm{FactPrincipalDatum}\),
not an additional principal-ideal assumption on
\(\mathfrak A\).  Conjugation of the definition gives
\(\bar\xi_H=-\xi_H\).  Every complex embedding sends
\(s\) to a root of \(z^2=-d\), whose absolute value is
\(\sqrt d\).  Hence its absolute value on \(\xi_H\)
is the same positive \(k_H\).  Since \(N\ge1,d\ge1\),
one has \(k_H\le1\).  At the source's fixed embedding
this equality is the definition of \(\xi R\).
\hfill\(\square\)
% END PCL8

% BEGIN PCL9
\paragraph{PCL9. The exact two host denominators.}
\label{claim:PCL9}
For any lattice \(\Lambda=\mathfrak A L'\) with the
original \(\mathfrak A\) of (PCL.14), set
\[
             N_A=N_E\mathfrak b,
             \qquad N_C=N_K C.                     \tag{PCL.17}
\]
Both are positive and
\[
               N_A\Lambda\subset L',\qquad
                           N_C L'\subset\Lambda.    \tag{PCL.18}
\]
These integers work for all components and embeddings;
they are the literal \(\mathrm{NA}\) and \(\mathrm{NC}\)
in the quantitative original endpoint.

\emph{Proof.}
For a nonzero integral ideal \(I\), the additive group
\(O/I\) has \(NI\) elements; thus \(NI\in I\).
In particular \(N_A\in\mathfrak b\), so
\(N_A\mathfrak A\subset C_E\subset O\).  Multiplying
\(\Lambda=\mathfrak A L'\) proves the first inclusion.
Also \(N_C\in C\), hence \(N_C\in C_E\).
Since \(\mathfrak b\subset O\),
\(N_C\mathfrak b\subset C_E\); equivalently
\(N_C\in C_E\mathfrak b^{-1}=\mathfrak A\).
This proves the second inclusion.  Nonzero ideals have
nonzero finite norm, so both integers are positive.
No ideal is principalized, and no integer depending on
the component has replaced either value.
\hfill\(\square\)
% END PCL9

% BEGIN PCL10
\paragraph{PCL10. The full half-different logarithmic bound.}
\label{claim:PCL10}
For the original host,
\[
 0<N_E\mathfrak b,\qquad
 (N_E\mathfrak b)^2\le|\operatorname{disc}(E)|,
 \qquad
 2\log N_E\mathfrak b\le\log|\operatorname{disc}(E)|.
                                                               \tag{PCL.19}
\]
This supplies \(\mathrm{HalfDiffNormSt}\), including its
positivity.  With the exact source definitions
\(f_E=[E:\mathbb Q]/2\) and
\(\operatorname{rd}(E)=|\operatorname{disc}(E)|^{1/[E:\mathbb Q]}\),
it gives
\[
        \log N_A\le f_E\log\operatorname{rd}(E).     \tag{PCL.20}
\]

\emph{Proof.}
Conjugation gives a ring isomorphism
\(O_E/\mathfrak b\simeq O_E/\bar{\mathfrak b}\), so their
finite cardinalities agree.  Multiplicativity of ideal
norm therefore gives
\(N_E\mathfrak D_{E/K}=(N_E\mathfrak b)^2\).
The different-in-a-tower identity writes
\(\mathfrak D_E\) as \(\mathfrak D_{E/K}\) times a
nonzero integral ideal.  The norm of that latter ideal
is a positive integer, and
\(N_E\mathfrak D_E=|\operatorname{disc}(E)|\).  This
proves the first two assertions.  Take the increasing
logarithm of the positive quantities to obtain the third.

A CM field has its complex embeddings paired by
fixed-point-free conjugation; thus its positive degree
is even and the source's natural-valued division by two
is exactly \([E:\mathbb Q]/2\).  Taking the logarithm
of the positive real power in the definition of
\(\operatorname{rd}\) shows
\(f_E\log\operatorname{rd}(E)=\tfrac12\log|\operatorname{disc}(E)|\).
This proves (PCL.20), with no hidden factor of two.
It also implies the weaker original bound
\(N_A\le|\operatorname{disc}(E)|\), since the latter
is a positive integer.
\hfill\(\square\)
% END PCL10

% BEGIN PCL11
\paragraph{PCL11. The literal normalized \(\mathrm{DataX}\) on every supplied chart.}
\label{claim:PCL11}
Fix the original \(D\) and \(T:\mathrm{Core}(D)\),
\(h:\mathrm{LawsCore}(T)\), a row \(j\) with
\(\mathrm{InR}(D,j)\), and any good embedding \(\sigma\).
Assume the original
\(\mathrm{HasPrinc}(T.U(j,\sigma),T.Q.l.M)\).
For every component \(k\), there is \(X_x\) with
\[
 \begin{gathered}
 X_x.\mathrm{toData}=(T.U(j,\sigma)).X(k),\qquad
 X_x.g=T.Q.l.M,\qquad X_x.k_0=2\xi R(D),\\
 X_x.L'=\mathrm{AA}^{-1}X_x.\Lambda,\qquad
 X_x.\Lambda=\mathrm{AA}\,X_x.L',\\
 X_x.N_A=N_E\mathfrak b,\qquad
 X_x.N_C=N_K\mathrm{ccK},\qquad
 X_x.N_A\le|\operatorname{disc}(E)|.
 \end{gathered}                                                   \tag{PCL.21}
\]
The fractional-ideal equality is the original span of
\(\alpha u\).  In particular this proves the entire
\(\mathrm{exists\_\allowbreak dataX\_\allowbreak of\_\allowbreak hasPrinc}\), with equality
in its advertised \(N_C\) bound.

\emph{Proof.}
The actual \(h.\mathrm{unf}\) field supplies PCL6, so
the trace pairing on this precise component is integral
and perfect.  Write its scalar as \(\xi\).
The field \(h.\mathrm{xiAbs}\) says
\(|\tau(\xi)|=\xi R(D)=|\tau(\xi_H)|\) at every
embedding, while \(\bar\xi=-\xi\) is part of its
\(\mathrm{Data}\).  At the one fixed embedding \(\phi_0\),
both values are purely imaginary and have equal absolute
value; hence one is the other or its negative.
Injectivity of \(\phi_0\) gives \(\xi=\xi_H\) or
\(\xi=-\xi_H\).  Their principal ideals coincide,
so PCL8 gives (PCL.7) for this actual component.

Apply PCL5 with \(\mathfrak A=\mathrm{AA}\),
\(k=\xi R(D)\), and \(g=T.Q.l.M\).
Its positive-modulus assumption follows from the
original level field \(M\ge3\); its congruence assumption
is precisely the given \(\mathrm{HasPrinc}\), with the
matrix acting on vectors on the left as in (PCL.1).
Use the prescribed integers in PCL9 instead of the
arbitrary integers in the proof of PCL5.  All record
fields still hold, now with (PCL.21); PCL10 gives the
claimed norm bound.  No change to \(T.U\), its component
index, \(G\), \(\xi\), \(\Gamma\), character or modulus
has occurred.  The argument works for every row and
every good embedding quantified in \(\mathrm{LawsCore}\),
not merely its distinguished complex reading.
\hfill\(\square\)
% END PCL11

% BEGIN PCL12
\paragraph{PCL12. The exact all-row model-lattice endpoints and remaining boundary.}
\label{claim:PCL12}
For the original core, retain \(\mathrm{GamAt}(T)\) and
\(\mathrm{LawsCore}(T)\).  Then the full
\(\mathrm{dataXOfModel\_\allowbreak of\_\allowbreak gamAt}\) holds: for every
row \(j\) satisfying \(\mathrm{InR}(D,j)\), every good
embedding and every component, there is
the unchanged \(\mathrm{Data}\) extended to \(\mathrm{DataX}\),
with \(g=T.Q.l.M\) and the same \(\mathrm{AA}\)-span equality.
For the original top record, the complete
\(\mathrm{DataXOfModelSt}(T,j)\) also holds.  In particular,
for every original \(\mathrm{large}_2\) and row weight \(w\),
under its \(\mathrm{SupBoundPageRows}\) and
\(2\le\mathrm{large}_2.\mathrm{sc}.L\) hypotheses, the
same chosen \(X_x\) satisfies
\[
 \log X_x.N_A\le f_E\log\operatorname{rd}(E),\qquad
 \log X_x.N_C\le2\,\mathrm{large}_2.\mathrm{sc}.L.     \tag{PCL.22}
\]

\emph{Proof.}
The source has a core version and a top version of
\(\mathrm{GamAt}\).  Each says, literally, for every
row, good embedding, component and unitary stabilizer
element \(\gamma\),
\[
       (\gamma-1)\Lambda\subset (T.Q.l.M)\Lambda
                      \ \Longrightarrow\ \gamma\in\Gamma.
                                                               \tag{PCL.23}
\]
Thus for the specified row and embedding it supplies
the exact \(\mathrm{HasPrinc}\) required by PCL11.
That claim gives the first endpoint, with all its
quantifiers.  Passing from \(\mathrm{Top}\) to its
\(\mathrm{Core}\) changes none of these data.

Choose the lattice by (PCL.8) and the integers by
(PCL.17), before any weight or \(\mathrm{large}_2\).
PCL10 gives the first inequality in (PCL.22).
The fourth original conjunct of
\(\mathrm{SupBoundPageRows}\) is exactly
\(\log N_K\mathrm{ccK}\le2\,\mathrm{large}_2.\mathrm{sc}.L\),
which gives the second.  All other conclusions of
\(\mathrm{DataXOfModelSt}\) are already (PCL.21).
Its extra size and weight hypotheses may therefore
be retained verbatim without changing the construction.
This also shows that the conclusion is uniform in
every positive original weight, and in both original
level choices whenever those supplied cores are used.
\hfill\(\square\)

The new result is the actual lattice and normalization
producer on a supplied principal chart.  The actual
existence of a common \(\mathrm{IsUnf}\) family with
\(\mathrm{HasPrinc}\), its framed-point normalization,
and its fixed metric dictionary remains separate.
So do the finite \(\mathrm{IsSystem}\), Gaussian
reproduction, short-rule gluing and ramified-jet inputs
retained by ICL.  This chapter does not infer full
\(\mathrm{LatFull}\), a higher \(\mathrm{CoversAt}\),
all-tuple class-sum or phase identities, or any compound
common-family closure from the constructed \(\mathrm{DataX}\).
% END PCL12


\section{Actual local growth, admission, and the canonical mean}
\label{sec:local-growth-suppliers-v047}

\noindent\textbf{Scope.}
This chapter supplies the local-growth and mean hypotheses that were
retained in GLG18, and assembles the complete admitted family used
there.  The actual common comparison record
\(\mathrm{GroupCompatibilityFamily}(\mathcal F)\) and the actual
selected-top predicate \(\mathrm{TopsF}(\mathcal F)\) remain explicit
inputs.  The conclusion concerns the same supplied
\(\mathcal F:\mathrm{FamP}(\epsilon)\), its lattices, degrees, support,
and local integer lengths.  We do not construct that common record.

The auxiliary chain introduced below is a variable in the local
counting argument.  It is not a replacement for the geometric theta
chain in the original arithmetic datum.  We verify precisely which
records and conclusions are independent of this auxiliary variable.

\paragraph{Notation and coefficient choices.}
Write \(m=D.R.m\) for the global chain length, \(k=D.R.k\) for the
coefficient rank, \(f=D.R.f\), \(n=D.R.n\), and \(n_v'\) for a local
Frobenius period.  The polynomial rank parameter is \(k-1\), not \(m\).
The source scales are
\[
 \begin{gathered}
 0<\epsilon<1/52,\qquad
 a=xe^{x^\epsilon},\quad X'=e^{x^{\epsilon/3}},
 \quad b=e^{x^\epsilon},\quad b'=a/X',\\
 \delta_1=x^{-\epsilon/8},\quad h=2+8\epsilon,\quad
 r=\lceil8(h+2)/\epsilon\rceil_++1,\\
 G'=\mathrm{Gp},\quad N_1=\mathrm{N1},\quad
 \Lambda=K.\mathrm{La}\,N_1^2\log a,\quad
 p_0=k^2x^{26},\quad c=4\mathrm{Sc}.
 \end{gathered}                                                   \tag{LGS.1}
\]
All ceilings with subscript \(+\) are natural-valued ceilings.
One sufficient original coefficient region is
\[
 \begin{gathered}
 K.\mathrm{La},K.\mathrm{kII},K.\mathrm{tII},
 K.\mathrm{DII},K.\mathrm{szII},K.\mathrm{kII}'\ge4129,\\
 K.\mathrm{kI},K.\mathrm{DI}\ge2800001,\qquad
 K.\rho\ge5600002,\\
 K.\kappa_1\ge604228,\qquad K.\kappa_2\ge302114,
 \qquad K.\mathrm{CN}\ge C_{\mathrm{orb}}(\epsilon).
 \end{gathered}                                                   \tag{LGS.2}
\]
Here \(\kappa_1,\kappa_2\) denote the two original
\texttt{kappa\allowbreak Threshold} coefficients and
\[
 C_{\mathrm{orb}}(\epsilon)
   =40(6r+1)\bigl((h+3)10^{6r+1}+6\bigr)>0.             \tag{LGS.3}
\]
The remaining coefficients are nonnegative.  In the original
\(\mathrm{FamP}\), the fixed fields are
\(K.\mathrm{Pr1}=1\), \(K.\mathrm{dl1}=5\),
\(K.\mathrm{dl2}=4\), and \(\mathcal F.cX.C_e=1\).
We retain them.  The constants in (LGS.2) are chosen before the
variable arithmetic data.  Thresholds in \(x\) follow the fixed
family and all its coefficients, and precede every \(i,u,w,p\),
orbit, coefficient extension, and degree assignment.

% BEGIN LGS1
\paragraph{LGS1. The integral polynomial sandwich at the actual places.}
\label{claim:LGS1}
For every original \(\mathrm{HidaSetting}\), every dominant integral
weight \(\lambda=(\lambda_1,\ldots,\lambda_k)\), and every coefficient
lattice \(A\mathcal O^k\), put
\[
 \begin{split}
 h_\lambda&=\lambda_1-\lambda_k,\qquad
 d_p(\lambda)=\binom{k}{2}\lfloor\log_p(h_\lambda+k-1)\rfloor,\\
 H_\lambda(A\mathcal O^k)
   &=\det(A)^{-\lambda_k}R(A^{-1})
                 \nabla_{\lambda-\lambda_k}(\mathcal O),\\
 D_\lambda(A\mathcal O^k)
   &=\det(A)^{-\lambda_k}R(A^{-1})
                 \Delta_{\lambda-\lambda_k}(\mathcal O).
 \end{split}                                                     \tag{LGS.4}
\]
For rank one the exponent is zero.  If \(h_\lambda=0\), the
inclusion below already holds with exponent zero; the displayed
nonnegative original exponent also works.  Then
\[
       p^{d_p(\lambda)}H_\lambda(A\mathcal O^k)
       \ \subseteq\ D_\lambda(A\mathcal O^k)
       \ \subseteq\ H_\lambda(A\mathcal O^k).            \tag{LGS.5}
\]
This supplies the literal \(\mathrm{ModelSandwichG}\) and
\(\mathrm{WeylSandwichG}\) on the original coefficient DVRs.

Here is the complete matching to the previously proved integral
comparison.  A Hida setting makes \(p\) a nonzero nonunit in
\(\mathcal O\).  Its embedded unramified DVR has characteristic zero:
in positive characteristic \(p\) would be either zero or a unit.
Thus the coefficient fraction field has characteristic zero.
The integer in the source is exactly
\(\mathrm{wExp}(p,k-1,\lambda)
 =N_0(k)\mathrm{deF}(p,k,h_\lambda)=d_p(\lambda)\).
The highest polynomial has coefficient one at its diagonal
monomial.  IW3 therefore applies to the actual coefficient
\(\nabla\) and the actual divided-power span \(\Delta\), with no
rescaling of either lattice or of their highest vector.
It proves \(p^{d_p(\lambda)}\nabla\subseteq\Delta\subseteq\nabla\).

The common invertible fraction-field map in (LGS.4) commutes with
the scalar \(p^{d_p(\lambda)}\), so it transports these inclusions.
Its determinant factor can have negative valuation; it occurs on
both sides and cancels.  Integral changes of basis preserve both
lattices, by CB2 and CR1.  Consequently the same inclusions hold
in the definitions made by intersection over all basis matrices.
For an input that has no presenting lattice basis, the
intersection in each model definition is over an empty
index set, so both are the ambient module and the required
inclusions are automatic.  The separate equivariant-function
definition requires zero values on nonlattice tuples.
For a genuine lattice and a nondominant weight, both
coefficient models are zero.  Thus none of these branches
requires an unverified basis choice.

For the odd Frobenius positions the weight is
\(-\operatorname{rev}\lambda\).  Its spread is again
\(h_\lambda\); (LGS.5) applies with exactly the original exponent.
Tensor products use disjoint polynomial-variable blocks.
Multiplication of polynomials identifies the tensor products
injectively, as may be checked on independent polynomial bases.
Thus exponents add over all positions, without any denominator
from the tensor realization.
% END LGS1

% BEGIN LGS2
\paragraph{LGS2. The original exponent really pays for every factor.}
\label{claim:LGS2}
Let the actual supplied records be
\[
       P:\mathrm{ParameterChoice}'(D,L),\qquad
       T:\mathrm{GatTies}(D).
\]
Suppose
\[
 e_p=X.\mathrm{ep}(p)
       =\left\lceil C_e f k^2\log_p a\right\rceil_+,
       \qquad C_e\ge1.                                  \tag{LGS.6}
\]
Then both fields of the original \(\mathrm{adm9}(D,L,X,C_e)\)
hold:
\[
 \begin{split}
 p\text{ prime},\ p\le a
   &\ \Longrightarrow\
       e_p\le(C_e+1)fk^2\log_p a,\\
 p^{e_p}G.M_{\mathrm{na}}
   &\ \subseteq G.M\subseteq G.M_{\mathrm{na}}
       \quad\text{at every original }(u,w,\mathrm{GIdx}).
 \end{split}                                                     \tag{LGS.7}
\]
The second line holds even when \(p>a\).

We first bound the weight spread on these very factors.
At a selected top the exact endpoint formulas in \(L\) give
\(\lambda_1-\lambda_k=a_t+b_t\).  The top window and
\(k\ge5,\ x^{13}\ge1\) give \(a_t\le6a/5\).
The first nonnegative summand of \(P.\rho_{\mathrm{le}}\)
and \(x^{-2}\le1\) give
\[
 \frac{N_*x^\epsilon\Lambda}{k}\le\frac14,\qquad
 b_t\le\frac54 b'\le\frac54a.
 \]
Hence the top spread is at most \(49a/20<3a\).
Away from tops, the original cloud endpoints give
\(\lambda_1-\lambda_k\le3\mathrm{Sc}\): subtract the lower
last coordinate from the upper first coordinate in the
\(\mathrm{cloud\_\allowbreak place\_\allowbreak bound}\),
\(\mathrm{uk\_\allowbreak extreme}\), and \(\mathrm{cloud\_\allowbreak v}\) formulas;
the \(\rho\)-shift contributes \(-(k-1)\).
The parameter inequality \(k\mathrm{Sc}\le w''\le a\)
therefore also gives a bound \(3a\).
The \(\mathrm{gat\_\allowbreak lam}\) and \(\mathrm{gat\_\allowbreak k}\) identities
transport these statements to every actual \(v,s\).
Odd reversal preserves the spread.

In particular \(0\le h_{\lambda_{v,s}}\le3a\).
The same parameter record gives \(a\ge4\), \(k\le a\), so
\[
       h_{\lambda_{v,s}}+k-1\le4a-1\le a^2.
 \]
Monotonicity of the logarithm and
\(\binom{k}{2}\le k^2/2\) imply
\[
 d_p(\lambda_{v,s})\le k^2\log_p a.                   \tag{LGS.8}
\]
The actual position bijection in \(T.\mathrm{gat\_\allowbreak pos}\) gives
\(\sum_v n_v'=f\), not an inequality with an unspecified
multiplicity.  Summing (LGS.8) yields
\[
 \sum_{v,s}d_p(\lambda_{v,s})
       \le fk^2\log_p a\le C_e fk^2\log_p a\le e_p.
 \]
LGS1 tensorized over precisely these positions now gives the
second line of (LGS.7) at every lattice tuple.  Taking equivariant
functions preserves the inclusions.  The unused exponent is an
additional nonnegative integral power of \(p\).

For \(p\le a\), set \(B=fk^2\log_p a\).  Since \(f\ge1\),
\(k\ge5\), and \(\log_p a\ge1\), one has \(B\ge1\).
The ceiling in (LGS.6) is at most \(C_eB+1\le(C_e+1)B\).
This proves the first line with its original coefficient.
In a \(\mathrm{FamP}\), \(\mathrm{ex.ep\_\allowbreak def}\) gives (LGS.6)
with \(C_e=1\) after the family start.  Thus this is an actual
supplier of \(\mathrm{Adm9F}\) and of
\(\mathrm{LatticeSandwich}\); neither is assumed in the
family application below.
% END LGS2

% BEGIN LGS3
\paragraph{LGS3. Uniform orbit, window, and logarithmic suppliers.}
\label{claim:LGS3}
Every supplied \(\mathcal F:\mathrm{FamP}(\epsilon)\) has the
following original outputs, with \(C_{\mathrm{orb}}\) as in
(LGS.3):
\[
 \begin{gathered}
 \mathrm{Adm1F}\text{ with }C_{31}=C_{\mathrm{orb}}>0,\\
 K.\mathrm{CN}\ge C_{\mathrm{orb}}
       \ \Longrightarrow\ \mathrm{FamilyWindowLogBound},\\
 K.\mathrm{La}>0
       \ \Longrightarrow\ \mathrm{Adm5F}.
 \end{gathered}                                                   \tag{LGS.9}
\]
The first assertion concerns every good orbit of the actual
chosen subset, not a new probabilistic choice.

To check the literal fields, GII1--GII3 give, for \(x\ge3\),
\[
 G'\delta_1\le C_{\mathrm{orb}}\log x,\quad
 N_1\le C_{\mathrm{orb}}\log x,\quad
 \frac{n'\delta_1}{C_{\mathrm{orb}}\log x}
       \le N_v\le40n'\delta_1.                         \tag{LGS.10}
\]
For each selected even position the least positive selected
even successor is within \(G'\) positions and has at most
\(N_1\) interior selected positions.  The number of selected
even predecessors in the preceding \(L_1\) positions is less
than \(r\), because \(L_1\le L_2\) and this is the actual
short-window condition in \(\mathrm{Good}\).
GJ3 identifies every one of these positions, successors and
predecessors with the stored orbit representative.
These are all four conjuncts of \(\mathrm{adm1}\).
They do not require a bound on the proportion of good orbits.
The separate predicate \(\mathrm{adm2}\) will still use the
actual selected-top input.

The middle line of (LGS.9) follows from the second inequality
in (LGS.10).  The definition of \(\rho_B\), nonnegativity of
the coefficients and the actual field
\(\mathcal F.\rho_{\mathrm{le}}\) consequently give
\[
              0\le\rho\le\rho_B\le x^{-2}.             \tag{LGS.11}
\]

For the last line, fix any \(\eta>0\) only after
\(K.\mathrm{La}\) and \(C_{\mathrm{orb}}\).  GJ2 proves all
four original logarithmic inequalities:
\[
 \begin{gathered}
 x^{\epsilon-\eta}\le\Lambda\le x^{\epsilon+\eta},\\
 p<p_0\ \Longrightarrow\ \log(p^{G'})
                       \le x^{\epsilon/8+\eta},\\
 p<X'\ \Longrightarrow\ \log(p^{G'})
                       \le x^{11\epsilon/24+\eta}.
 \end{gathered}                                                   \tag{LGS.12}
\]
For completeness, their single numerical threshold may be
chosen by requiring
\[
 (2K.\mathrm{La}\,C_{\mathrm{orb}}^2
       +38C_{\mathrm{orb}}+(K.\mathrm{La})^{-1})
                  (\log x)^3\le x^\eta,\qquad x\ge3.
 \]
The source bounds \(k\le x^6\), hence
\(\log p_0\le38\log x\); insert (LGS.10) and
\(\log a=\log x+x^\epsilon\).
The positive constant and logarithm powers are absorbed by
the displayed threshold.  It depends on \(\eta\), as the
literal \(\mathrm{Adm5F}\) permits.  \(C_{31}\) does not.
% END LGS3

% BEGIN LGS4
\paragraph{LGS4. The actual small-order prime contribution.}
\label{claim:LGS4}
Every supplied \(\mathrm{FamP}(\epsilon)\) satisfies the
literal \(\mathrm{Adm3F}\):
\[
 \sum_{\substack{p<X''=xnX'\\p\ \mathrm{prime},\
                      \mathrm{IsResidueClassZero}\\
                 \operatorname{ord}(h_P)<x^{\epsilon/4}}}
                    \frac{\log p}{p}
              \le x^{-h+\epsilon}                      \tag{LGS.13}
\]
eventually, uniformly in the family index.

We must first justify use of the dyadic estimate on this
particular Frobenius.  A residue-class-zero prime is odd and
does not divide the host conductor.  The original host
Frobenius fields therefore supply an arithmetic Frobenius
above it.  On the cyclotomic field of conductor \(2^{e+2}\),
arithmetic Frobenius sends \(\zeta\) to \(\zeta^p\).
The host's recorded Galois identification intertwines its
dyadic projection with the class of this \(p\) modulo
\(\{\pm1\}\).  This holds for every permissible prime and
Frobenius, and hence for the particular element used in the
definition of \(h_P\).  In the original notation it says
\[
 D.R.h(h_P)=
   \operatorname{ofMul}\!\bigl(\iota_H(h_{\mathrm{dyad}}(p))\bigr).
                                                               \tag{LGS.14}
\]
For a homomorphism, the order of the image divides the order
of the original element: apply the homomorphism to its
annihilating multiple.  Thus
\(\operatorname{ord}(h_{\mathrm{dyad}}(p))
 \le\operatorname{ord}(h_P)\).
This supplies the exact \(\mathrm{FrobeniusClassCompatibility}\)
used here; it is not a replacement of \(h_P\) by a chosen
conjugacy representative.

Put \(M=2^e\) and \(J=x^{\epsilon/4}\).
The rank bounds below give \(M\ge x^{h-\epsilon/50}>1\)
once \(x>1\), so the dyadic exponent satisfies \(e\ge1\).
The integer dyadic orders in this sum are at most
\(\lfloor J\rfloor\).  LFS2, with
\(\lfloor J\rfloor\le J\), gives
\[
 \sum\frac{\log p}{p}
       \le\frac{12J\log^2(MX'')}{M}.                  \tag{LGS.15}
\]
The source rank relations and its \(\mathrm{SizeN}\) and
\(\mathrm{o\_\allowbreak tie}\) fields are
\[
 M\le n\le oM,\quad o\le x^{\epsilon/100},\quad
 x^{h-\epsilon/100}\le n\le x^{h+\epsilon/100}.
 \]
Consequently \(M\ge x^{h-\epsilon/50}\) and
\(\log M\le\log n\le3\log x\).
Choose the threshold so that
\[
 8\log x\le x^{\epsilon/3},\qquad
                  48\le x^{19\epsilon/300}.
 \]
Then
\[
 0\le\log(MX'')\le
       7\log x+x^{\epsilon/3}\le2x^{\epsilon/3}.
 \]
The right side of (LGS.15) is at most
\[
 48x^{-h+(1/4+2/3+1/50)\epsilon}
 =48x^{-h+281\epsilon/300}
 \le x^{-h+\epsilon}.
 \]
Both thresholds precede every prime and every member.
The actual sums are finite; using \(p<X''\) instead of
\(p\le X''\) only restricts the summation set.
No new prime-distribution theorem is required.
% END LGS4

% BEGIN LGS5
\paragraph{LGS5. All five original power-band inequalities.}
\label{claim:LGS5}
Under \(K.\mathrm{CN}\ge C_{\mathrm{orb}}\), the actual family
satisfies \(\mathrm{Adm4F}\), including both prime sums with
their original bound \(x^{-1}\).

We prove the finite estimate used for both bands.  Let
\(X\ge1\), \(P_0\ge1\), \(0\le q\le1/3\), \(d\ge1\), and
let \(S\) be any finite set of positive integers such that
\[
 P_0\le p<X,\qquad X(1-q)\le p^d\le X(1+q)
                       \quad(p\in S).
 \]
If \(S\ne\varnothing\), let \(u=\min S\), \(v=\max S\).
Since \(v/u\ge1\),
\[
 \frac vu\le\left(\frac vu\right)^d
       \le\frac{1+q}{1-q}\le1+3q.
 \]
There are at most \(v-u+1\le3qu+1\) integers between the
endpoints, while \(\log p/p\le\log X/u\).
Thus, including the empty case,
\[
             \sum_{p\in S}\frac{\log p}{p}
                 \le\log X(3q+P_0^{-1}).             \tag{LGS.16}
\]
If each \(p\) is prime and some positive exponent \(d\) works,
choose one such exponent per \(p\).  As \(p\ge2\) and
\(p^d\le2X\),
\[
       d\le1+\frac{\log X}{\log2}\le1+2\log X.
 \]
Partition the finite set by the chosen exponent and apply
(LGS.16) on each part.  This yields
\[
 \sum_{p\in S}\frac{\log p}{p}
       \le(1+2\log X)\log X(3q+P_0^{-1}).             \tag{LGS.17}
\]
This proof counts an endpoint integer once even for a
zero-width band; it does not discard the \(P_0^{-1}\) term.

Apply (LGS.17) with \(X=X'\), \(P_0=p_0\), and first
\(q=\rho\), then \(q=\rho_B\).  LGS3 gives
\(0\le q\le x^{-2}\).  For \(x\ge2\) this is at most \(1/4\),
and \(p_0=k^2x^{26}\ge x^2\).  With
\(B=\log X'=x^{\epsilon/3}\ge1\), the right side is at most
\[
       (3B)B(4x^{-2})=12x^{2\epsilon/3-2}\le x^{-1}
                                                               \tag{LGS.18}
\]
as soon as \(12\le x^{1-2\epsilon/3}\).
The exponent is positive since \(\epsilon<1/52\).
Each original \(\mathrm{InBand0}\) or \(\mathrm{InBand}\)
prime satisfies precisely the premises of this finite
estimate.  Together with (LGS.11), the two bounds in
(LGS.18) are all five fields of the original
\(\mathrm{ArchimedeanSignatureMatching}\).
They use one common threshold, independent of the prime
and its witnessing exponent.
% END LGS5

% BEGIN LGS6
\paragraph{LGS6. The actual two projection certificates and canonical mean.}
\label{claim:LGS6}
Under (LGS.2), the same instance \(X=\mathcal F.Gx(x,i)\)
eventually supplies all three original predicates
\[
 \mathrm{adm6}(D,L,K),\qquad
 \mathrm{adm7}(D,L,X,K),\qquad
 \mathrm{adm8}(D,L,X,K,302114).                         \tag{LGS.19}
\]
In addition, the actual \(\mathrm{HypIAt}\) and
\(\mathrm{HypIIAt}\) hold at every permitted prime and good
orbit.  All coefficient settings and dominant weights in
these predicates retain their original universal quantifiers.

Here are the numerical and model identifications needed to use
the completed local theorems.  LGS3 supplies the same
\(C_{31}\), window count and \(\mathrm{adm5}\) with
\(\eta=\epsilon/48\).  LGS5 supplies exclusion of the original
power band.  In the large-prime range the fourteen inequalities
of the source's \(\mathrm{HypI}\), together with the two
spacing inequalities and actual degree windows, are exactly
GI1--GI5.  Their coefficients are
\[
 C_H=2800001,\quad C=302114,\quad C_0=K.\mathrm{La},\quad
 m_{\mathrm{loc}}=k-1,\quad c_{\mathrm{loc}}=\lfloor4\mathrm{Sc}\rfloor_+.
                                                               \tag{LGS.20}
\]
The lower bounds \(K.\mathrm{kI},K.\mathrm{DI}\ge C_H\),
\(K.\rho\ge2C_H\), \(K.\kappa_1\ge2C\),
\(K.\kappa_2\ge C\) and \(K.\mathrm{La}\ge880\)
are all consequences of (LGS.2).  The candidate excludes the
band of radius \(\rho_B\); LGS3 gives \(\rho\le\rho_B\),
so it also excludes the local forbidden band.  The source
windows are the literal restrictions of \(a_L,b_L\) to
the cyclic orbit, by GJ3.  No recentering is performed.

In the small-prime range GII4--GII7 supply every numerical
and geometric hypothesis of the Type-II theorem on the same
orbit.  GII8 applies with the single original constant \(4129\),
since it is no larger than
\[
 K.\mathrm{kII},\ K.\mathrm{tII},\ K.\mathrm{La},\
 K.\mathrm{DII},\ K.\mathrm{szII},\ K.\mathrm{kII}'.
\]
This is precisely the full \(\mathrm{adm6}\); it is not
restricted to coefficient rings occurring at a chosen
global place.  In particular it remains valid after the
standard-ring construction used below.

For \(\mathrm{adm7}\), the supplied family field
\(\mathrm{ex.kapI\_\allowbreak def}\) says that \(X.\kappa_I\) is the
infimum of the deficits for which the required certificate
holds universally over those coefficient settings.
GJ1 proves attainment.  To recall why this has no compactness
gap, in one fixed setting a certificate of real rank at least
\((1-q)d\) has an integral number of independent tests.
For \(q\) sufficiently close above its infimum \(q_0\),
that integer is at least \(\lceil(1-q_0)d\rceil\).
All other tests and divisibilities are unchanged.  The
setting was arbitrary, so \(q_0\) itself satisfies the
universal predicate.  This proves \(\mathrm{adm7}\) for the
original \(X.\kappa_I\), rather than a newly chosen function.

Finally GI6--GI8 compare the local theorem's time with the
original \(\lceil t_{If}\rceil_+\), retaining the common
signed leading sum.  Its exceptional function is chosen
before the coefficient setting and has mean at most
\[
                   302114\,N_vN_1/x^{13}.             \tag{LGS.21}
\]
Leastness of the attained canonical deficit bounds
\(X.\kappa_I\) by this same function.  Restriction to the
selected orbit positions is a bijection between the
original degree box and the local product box, with inverse
extension by zero.  Thus both sums and both cardinality
denominators agree, proving \(\mathrm{adm8}\) with
\(C_m=302114\).  Empty boxes obey the original zero-average
convention.  Taking the maximum of the thresholds just
described proves (LGS.19) uniformly in \(i\), before any
prime, degree assignment or coefficient extension is chosen.
% END LGS6

% BEGIN LGS7
\paragraph{LGS7. The literal mean bound on the original canonical deficit.}
\label{claim:LGS7}
Under (LGS.2), the supplied family eventually satisfies
\(\mathrm{MeanKappaBound}(\mathcal F.Gx,K)\).
Explicitly, for every candidate prime and every actual good
representative,
\[
 \begin{split}
 \operatorname{meanLoc}(\kappa_{\mathrm{loc}})
       &\le N_vx^{-12},\\
 p<p_0\quad&\Longrightarrow\quad
       \kappa_v(\mathbf a,\mathbf b)\le N_vx^{-12}
       \quad\hbox{for every }\mathbf a,\mathbf b.
 \end{split}                                                     \tag{LGS.22}
\]
This supplies both clauses, including the pointwise
small-prime clause, on the same \(X\) as LGS6.

Indeed the definition in the small-prime range is
\[
 \kappa_{\mathrm{loc}}=\kappa_v
     =\min\{1,K.\mathrm{kII}\,N_vr^2k x^{-26}\}.
 \]
It is independent of the degrees and nonnegative.  Since
\(k\le x^6\), the single requirement
\(K.\mathrm{kII}\,r^2\le x^8\) makes it at most
\(N_vx^{-12}\).  Its average has the same bound; if its
box is empty, the average is zero.
For \(p\ge p_0\), the candidate's order lower bound and
LGS3 supply the orbit, spacing and window hypotheses of
LGS6.  Consequently (LGS.21), together with
\[
                302114\,C_{\mathrm{orb}}\log x\le x,   \tag{LGS.23}
\]
gives the first line of (LGS.22).  Both requirements are
eventual in \(x\) after the fixed coefficients.  They are
independent of \(p,v,i\), of the coefficient DVR, and of
the degrees.  This is the numerical content of GW1--GW2
matched to the full original \(\mathrm{MeanKappaBound}\).

In particular no assumption that the actual selected
degrees already have a small deficit at every prime has
been inserted.  The later finite averaging theorem can
use (LGS.22) to choose a single degree assignment; it does
not change the function proved to satisfy this mean bound.
% END LGS7

% BEGIN LGS8
\paragraph{LGS8. What the common compatibility record actually supplies.}
\label{claim:LGS8}
Let the original \(\mathrm{GroupCompatibilityFamily}(\mathcal F)\)
be supplied.  Then, after one family threshold, the actual
local models have finite class sets and trivial stabilizers,
the standard coefficient-subring interface
\(\mathrm{StdSubAt}\), and \(\mathrm{NPDescF}\).
Both \(G.M_{\mathrm{na}}\) and \(G.M\) are finite free of the
same rank \(G.d\).  The two original operator-stability
conditions hold on these very lattices.

The matching uses every part of the record in its intended
role.  At fixed \(x,i\), its fields are
\[
 \mathrm{herm},\quad \mathrm{forms},\quad
 \mathrm{cmp},\quad \mathrm{ring},\quad \mathrm{gat}.    \tag{LGS.24}
\]
The last field identifies \(D.gat(u,w,h_g)\) with the
actual \(\mathrm{gatOfLevelP}\) and includes the original
\(\mathrm{Cmp.HypAtP}\); the level is the very
\(\mathrm{herm.lev0}(u)\).  It is not merely an isomorphism
of two abstract representation spaces.
The two distinct auxiliary primes and their
upper-unitriangular level conditions are part of this
level datum.  CR3 identifies its lattice orbits with its
finite class set, and CR4 makes the stabilizers trivial.
CR2 then identifies integral equivariant functions with
the finite product of their actual fibers.  CR1 gives
finite free Weyl fibers; CB2 gives finite free
dual-Weyl fibers.  LGS1 makes their rational spans equal.
Their common function-space rank is therefore exactly
\(G.d\), with the original character and determinant
twists.  This proves the literal
\(\mathrm{TrivialStabiliserFiniteClasses}\), not just a
finite-dimensionality assertion over the fraction field.

For the standard subring, \(\mathrm{ring}\) supplies four
compatible isomorphisms \(W,Q,\mathcal O,K\), including the
maps of the unramified subring, its fraction field, and
the chosen square root of \(p\).  Compose the finite map
constructed in BC1--BC2 with these isomorphisms.
BC8 shows that its three compatibility equations are
preserved.  Thus the target is the actual coefficient
DVR and the image of the chosen square root is the
recorded \(w\), even if it is not a uniformizer.
BC6--BC7 consequently transport the local rank
certificates without replacing the residue cardinality
of the unramified ring or changing any \(p\)-normalized
depth.

Finally the \(\mathrm{cmp}\) field is the same
\(\mathrm{CmpTied}\) used with the given
\(\mathrm{GatTies}\) and the family's actual comparison
map.  GLG14 applies its rational retraction and faithful
flatness argument to the ordered products and ramified
intersection, proving \(\mathrm{NPDescF}\).
GLG13 proves both operator-stability statements directly
for the actual polynomial models, including signed
determinant factors and nondominant zero fibers.
No extra integrality or comparison premise remains at
these steps.  Construction of (LGS.24) itself is the
explicit retained input of this chapter.
% END LGS8

% BEGIN LGS9
\paragraph{LGS9. The actual orbit certificates give joint exceptional generators.}
\label{claim:LGS9}
Assume (LGS.2) and the common compatibility record.
At every actual \((u,w,h_g)\), every served representative
\(t_0\) supplies the original \(\mathrm{LPprime}\) with
\[
 \begin{gathered}
 v=vOf(t_0),\qquad t_v=\lceil t_v^{\mathrm{source}}\rceil_+,
 \qquad D_0=D_p,\qquad \kappa_v=\kappa_v^{\mathrm{source}},\\
 c_*=8(D_p+1),\qquad
 r_{\mathrm{loc}}=
       \begin{cases}r,&p<p_0,\\1,&p\ge p_0.\end{cases}
 \end{gathered}                                                   \tag{LGS.25}
\]
It supplies \(\mathrm{OrbitLatticeBounds}\) for the original
data, not just for isomorphic standard models.
For every nonempty subset \(B\) of the served local
coordinates, the actual dual-slot section lattice \(L_B\)
has a finite set \(E_B\subset L_B\) such that
\[
 \begin{split}
 |E_B|&\le d\prod_{v\in B}\kappa_v,\qquad d=G.d,\\
 p^{D_p+1}L_B
   &\subseteq
     \sum_{v\in B}y_v L_{B;\,v=\mathrm{ball}}
                         +\langle E_B\rangle_{\mathcal O}.
 \end{split}                                                     \tag{LGS.26}
\]

For (LGS.25), a served representative belongs to the actual
good representative set of a candidate prime.  Its cyclic
coordinate bijection is GJ3.  The \(\mathrm{LatTied}\)
equalities give the same period, selected positions,
top and bottom degrees, filling weights and local rank
parameter.  The remaining coordinate bound is at most
\(\lfloor4\mathrm{Sc}\rfloor_+\).
LGS6 provides the applicable Type-II or Type-I certificate
over the standard subring; LGS8 and BC6--BC7 transport it
to the actual \(\mathcal O\).  GJ4 checks the four
quantities in (LGS.25) exactly.  Thus no height, depth or
exceptional-rank allowance is enlarged in this transport.

We now check all geometric hypotheses for the product
estimate in (LGS.26).  In each coefficient factor, EL5--EL7
construct the source, midpoint and target lattices from an
actual edge in Smith coordinates.  With
\[
 A_v=\sum_s\sum_{j=1}^{r_{\mathrm{loc}}}\lambda_{s,j},
 \]
the normalized edge pairing is
\(\beta_e(F,\ell)=\ell(p^{A_v}F)\).
Its integral perfection follows from
\(w^{-A_v}p^{A_v}w^{-A_v}=1\); \(A_v\) may be negative.
The ball and dual slots are the actual intersections and
polars in EL7, so for fixed other slots
\[
 P_t\subseteq H\subseteq B_t,\qquad
 w^tB_t\subseteq H,\quad w^tH\subseteq P_t,\quad
 y_v(B_t)\subseteq P_t.                                \tag{LGS.27}
\]
EL8--EL10 prove their equivariance and transport the
unwanted-edge errors.  These statements are valid after
tensoring all unchanged coordinates and after taking
the actual equivariant sections.  All slot lattices are
finite free and full by LGS8 and their two-sided
comparison with the base fibers.
The map \(vOf\) is injective on the representative set
by \(\mathrm{LatTied}\), so distinct representatives
produce distinct coordinates.  Operators at distinct
coordinates commute: both compositions are the same
finite double sum over independent incoming edges.
Their scalar factors also commute.

For clarity, the integral rank argument is as follows.
The \(\mathrm{LPprime}\) tests give an integral matrix
identity equal to \(p^{D_p}I_n\), with
\(n\ge(1-\kappa_v)d_v\), on the local dual-slot lattice.
JR4 applies Smith form and finds at least \(n\) basis
directions divisible by these tests at depth \(D_p\).
At most \(\kappa_vd_v\) directions are bad.
At an actual orbit representative tensor the bases and
retain a tensor only if every selected coordinate is bad.
Its count is at most
\((\prod_{v\in B}\kappa_v)\prod_vd_v\).
The finite orbit and trivial-stabilizer facts of LGS8
extend each retained value uniquely to a delta section.
The total base rank is \(d\), so the count is the first
line of (LGS.26).

Every discarded tensor has a good selected coordinate.
Applying its neighbor relation gives the required value
once at that target, since its stabilizer is trivial.
All unintended targets lie in the stated
\(w^{c_*}\)-error lattice; EL10 supplies the integral
identification with the target dual slot.
Here \(c_*=8(D_p+1)\ge4(D_p+1)\), so the error has
at least the \(p^{D_p+1}\) divisibility used in JR5.
The resulting finite quotient satisfies \(Q=pQ\);
the elementary determinant proof of Nakayama, or JR7,
makes it zero.  This proves the second line of
(LGS.26) with the original exponent \(D_p+1\).
The argument has now supplied the quantitative,
polynomial, orbit, and stabilizer inputs that were
separate in JR1--JR12 and LW4.
% END LGS9

% BEGIN LGS10
\paragraph{LGS10. The literal first-growth inclusion and length.}
\label{claim:LGS10}
Under the same hypotheses, fix any actual place and word.
Put \(M=G.M\), \(M_{\mathrm{na}}=G.M_{\mathrm{na}}\),
\(U=Y.Vram\), and \(M_{\mathrm{ram}}=M\cap U\).
Then eventually, with the exact original \(\tau_P\),
\(\mathrm{io}_P\), and \(\mathrm{NP}\),
\[
 \begin{split}
 \mathrm{NP}&\subseteq p^{\lceil\tau_P\rceil_+}M_{\mathrm{ram}},\\
 \ell:=\operatorname{length}_{\mathcal O}
                 (M_{\mathrm{ram}}/\mathrm{NP})&\in\mathbb N,
 \qquad
 \ell\le vUnit(G)\,G.d\,(\tau_P+\mathrm{io}_P).
 \end{split}                                                     \tag{LGS.28}
\]
These are the first two conjuncts of the original
\(\mathrm{admG1}\), including the finite natural-valued
length, for every \(\mathrm{GatTies}\) and
\(\mathrm{LatTied}\) witness on the fixed family data.

Let \(\mathcal V\) denote the actual served set.
If it is empty, the empty product term in the definition
of \(\mathrm{NP}\) is \(M\), while
\(\tau_P=\mathrm{io}_P=0\).  Hence
\(\mathrm{NP}=M_{\mathrm{ram}}\) and (LGS.28) follows.
Suppose it is nonempty.  Set
\[
 \begin{gathered}
 S_v=\mathrm{Sv}(v),\quad S=\max_{v\in\mathcal V}S_v,\quad
 \Sigma=m\sum_{v\in\mathcal V}\kappa_v,\\
 \tau_0=\sum_{v\in\mathcal V}m(t_v+D_p+1),\qquad
 \tau=e_p+\tau_0 .
 \end{gathered}                                                   \tag{LGS.29}
\]
Each \(S_v\) is already the integer ceiling of a
nonnegative sum, so is a nonnegative integer.
The source's outer natural-valued ceiling therefore
does not add one.  The source definition similarly
gives \(\tau_P=\tau\) exactly.
Because a nonempty served set is precisely the whole
good representative set at a candidate with
\(\mathrm{Sig}\le x^{-\epsilon}\), one has
\[
       \Sigma=\mathrm{Sig}\le x^{-\epsilon}\le\tfrac12,
       \qquad \mathrm{io}_P=2S\Sigma                  \tag{LGS.30}
\]
as soon as \(x\ge2^{1/\epsilon}\).
The original rank record already gives \(m\ge2\).

Use the actual slot lattices and joint relations of
LGS9 in the finite walk LW1--LW3.  For
\[
 J_{\mathrm{na}}=
 \sum_{B\subseteq\mathcal V}
     p^{\sum_{v\notin B}S_v}
       \Bigl(\prod_{v\in B}y_v^m\Bigr)M_{\mathrm{na}},
 \]
it gives, for every integer \(\beta\ge0\), at most
\[
               d\,m^{\beta+1}e_{\beta+1}(\kappa)       \tag{LGS.31}
\]
exceptional directions at depth \(\tau_0+\beta S\).
One sees the multiplicity directly in that walk:
choose the \(\beta+1\) active coordinates and one of
the \(m\) counter positions at each active coordinate.
Its depth cost at a coordinate is exactly
\(m(t_v+D_p+1)\).  Thus neither \(m\) nor the \(+1\)
is omitted.

The actual sandwich LGS2 and operator stability LGS8
give \(p^{e_p}J_{\mathrm{na}}\subseteq J\), where \(J\)
is defined by the same formula with \(M\).
Multiply the walk inclusion by \(p^{e_p}\).
The exceptional vectors then lie in \(M\), their number
is still (LGS.31), and their depth in \(M\) is
\(\tau+\beta S\), exactly as in SR3.
In particular the empty product term shows that \(J\)
contains a nonzero scalar multiple of \(M\), so all
lengths in this calculation are finite.

Write \(e_{\mathcal O}=v_{\mathcal O}(p)=vUnit(G)>0\).
By LC1--LC2, summing the Smith tails gives
\[
 \begin{split}
 \operatorname{length}_{\mathcal O}
       \bigl(M/(J\cap p^\tau M)\bigr)
 &\le e_{\mathcal O}d\left[
       \tau+S\left(\prod_v(1+m\kappa_v)-1\right)\right]\\
 &\le e_{\mathcal O}d(\tau+2S\Sigma).
 \end{split}                                                     \tag{LGS.32}
\]
The last inequality follows from
\(\prod_v(1+m\kappa_v)\le e^\Sigma\) and
\(e^\Sigma-1\le2\Sigma\) for \(0\le\Sigma\le1/2\).
For \(S=0\), the vanishing of (LGS.31) past
\(|\mathcal V|\) already puts every Smith exponent
at depth at most \(\tau\); the same formula holds.

The actual field \(\mathrm{ex.lat.vram}\) identifies
\(U\) with the coefficient-field span of the original
ramified subspace.  If \(a\ne0\) and \(az\in U\), its
inverse in that field gives \(z\in U\).  Therefore
\(M_{\mathrm{ram}}\) is saturated in \(M\), and
\(p^\tau M\cap M_{\mathrm{ram}}=p^\tau M_{\mathrm{ram}}\).
Now
\[
 \mathrm{NP}=J\cap p^\tau M\cap U.
 \]
This proves the inclusion in (LGS.28); the quotient
\(M_{\mathrm{ram}}/\mathrm{NP}\) injects into the
finite quotient in (LGS.32).  Its natural-valued length
therefore obeys the same bound.  Substituting (LGS.29)
and (LGS.30) proves (LGS.28) with the original full
rank \(d=G.d\) and normalization \(vUnit(G)\).
% END LGS10

% BEGIN LGS11
\paragraph{LGS11. An auxiliary zero chain preserves the complete interface.}
\label{claim:LGS11}
Define \(Y^0\) from the original \(\mathrm{LatInst}\)
\(Y\) by keeping \(vOf,Fl,Lv,Vram\) and replacing every
linear map of its chain \(Ch\) by zero.
Then the same family eventually satisfies
\[
 \begin{gathered}
 \mathrm{unitBoundAtTwistedPlaces}(D,L,X,Y^0,K),\\
 \mathrm{UnitBoundDifferenceTerm}(D,L,X,Y^0,K),\qquad
 \mathrm{admG1}(D,L,X,Y^0,K).
 \end{gathered}                                                   \tag{LGS.33}
\]
This construction requires no lower-valuation estimate
for the original theta chain.

We check the record rather than only its final numerical
expression.  A \(\mathrm{Chain}\) consists solely of
coefficient-field linear maps \(K_j\) between the
successive spaces \(Lv.Fm(j)\).  Zero maps satisfy that
definition.  Since \(m\ge2\), their composite
\(K_{m-1}\cdots K_0\) is zero on every input.
The predicate \(Y.vKge(F,g)\) asks for a uniformizer
\(\varpi\) such that this composite lies in
\[
 \varpi^{\lceil g\,vUnit(G)\rceil}\,Lv.Int(m).
 \]
A DVR has a uniformizer, and zero lies in this scaled
module for every real \(g\), including negative \(g\).
Thus \(Y^0.vKge(F,g)\) holds for every \(F,g\).

Every clause of \(\mathrm{LatTied}\) uses only \(vOf\),
\(Fl\), and the fixed \(D,X\).  Its period, rank,
selected-coordinate, filling, and injection conditions
are consequently unchanged.  The ordered product
\(\mathrm{yProd}\) uses only \(vOf\) and the actual
\(G.y_v^m\), so it too is unchanged.  In the complete
definition of \(\mathrm{NP}\), the three factors are
the span of these products on \(M\), the scalar-depth
copy of \(M\), and \(Vram\).  Hence
\[
 \mathrm{NP}(X,Y^0,K)=\mathrm{NP}(X,Y,K),\qquad
 M_{\mathrm{ram}}(Y^0)=M_{\mathrm{ram}}(Y).             \tag{LGS.34}
\]
The rational descent and scalar-extension identities
transfer exactly.  So does every original
\(\mathrm{LatticeLengthsAtPlace}\) identity.  The numerical
functions \(\kappa_v,S_v,\mathrm{Sig}\), \(VV\),
\(\tau_P,\mathrm{io}_P\), and \(\mathrm{gainP}\)
do not mention \(Ch\).
Neither do \(\mathrm{ParametersFill}\),
\(\mathrm{NuDef}\), the support condition, or any of
the five fields of \(\mathrm{LinkG\_\allowbreak at}\).
The actual selected global lattice \(L.NuLat\) and
its comparison with \(M_{\mathrm{ram}}\) are unchanged.
More explicitly, \(\mathrm{LatDefs}\) has just the
\(\mathrm{vram}\), \(\mathrm{served}\), and
\(\mathrm{atT}\) fields: the first two transfer by
(LGS.34), and the third does not mention \(Y\).
The larger arithmetic \(\mathrm{Extra}\) record
also has a \(\mathrm{ker}\) field about its geometric
chain.  That record remains the original record;
we do not transport its \(\mathrm{ker}\) field to
\(Y^0\).  Neither auxiliary admission nor the
\(\mathrm{LinkG}\) installer requires that transport.

The first two lines of (LGS.33) are now immediate even
with their full quantifiers: for every pair of subsets
\(VV'\subseteq VVs\subseteq\mathrm{reps}\), every
original section, and the printed lower term, the
left composite is zero.  The same observation gives
the final valuation conjunct of \(\mathrm{admG1}\)
for every \(F\in\mathrm{NP}\) and the exact lower term
\(\sum_j\mathrm{gainP}_j+\mathrm{lowFE}\).
The first two conjuncts are LGS10, which apply to the
unchanged \(\mathrm{NP}\) and \(M_{\mathrm{ram}}\).
This proves the full original \(\mathrm{admG1}\) on
the auxiliary \(Y^0\).

Only this auxiliary argument has a zero chain.
The given arithmetic datum, its geometric chain,
its valuations, and all divisibility hypotheses used
elsewhere have not been altered.  Nothing in
(LGS.33) asserts a new estimate for those original
theta maps.
% END LGS11

% BEGIN LGS12
\paragraph{LGS12. The complete original admitted family is now supplied.}
\label{claim:LGS12}
Assume (LGS.2), the original
\(\mathrm{GroupCompatibilityFamily}(\mathcal F)\), and
the actual \(\mathrm{TopsF}(\mathcal F)\).
Then
\[
 \mathrm{AuxBlock8.AdmF}
   \bigl(\mathcal F,K,\mathcal F.Gx,
                     (x,i)\mapsto(\mathcal F.Gy(x,i))^0\bigr)      \tag{LGS.35}
\]
holds.  The constants may be fixed as
\[
       C_{31}=C_{\mathrm{orb}}(\epsilon),\qquad
                         C_e=1,\qquad C_m=302114.     \tag{LGS.36}
\]
Equivalently, these same constants supply the full
\(\mathrm{AdmFam}\).  The assumption \(\mathrm{TopsF}\)
may instead be replaced by the literal \(\mathrm{Adm2F}\).

We give every field and its threshold, since a partial
admission record would not suffice here.
LGS3 gives \(\mathrm{adm1}\) with this \(C_{31}>0\).
GLG16 applies to the original selected tops and gives
the full \(\mathrm{adm2}\): both its cardinality
inequalities and its weighted bad-top bound.
It takes a threshold after the one actual constant
in \(\mathrm{ChosenBad}\), without resampling the tops.
LGS4 and LGS5 give \(\mathrm{adm3}\) and all five
fields of \(\mathrm{adm4}\).
For any \(\delta>0\), LGS3 gives \(\mathrm{adm5}\)
with that \(\delta\).  Its threshold may depend on
\(\delta\), and (LGS.36) does not.
LGS6 gives \(\mathrm{adm6},\mathrm{adm7},\mathrm{adm8}\)
with the canonical deficit and the constant \(C_m\)
in (LGS.36).  LGS2 gives \(\mathrm{adm9}\) with
the actual fixed \(C_e=1\).
LGS11 supplies the unit bound and every conjunct of
\(\mathrm{admG1}\) for \(Y^0\).

It remains to check the original
\(\mathrm{TypePrimeSizeConds}\), which has three fields.
The family's \(\mathrm{sizes}\) field gives all of
\(\mathrm{SizesAt}(L)\).
The actual parameter record and \(\mathrm{fills}\)
give, respectively, its \(S\)-bound and the equality
of the two \(\Lambda\)'s, hence
\[
                  4\mathrm{Sc}\le\Lambda b'/k.
 \]
GLG1 gives \(N_m\le\mathrm{Sc}\),
\(\mathrm{Sc}\ge100\), \(k\ge5\) and
\(k\mathrm{Sc}\le w''\).  Therefore
\(N_m+1\le\mathrm{Sc}+1\le k\mathrm{Sc}\le w''\).
This is the third field, with the original \(N_m\).

Fix (LGS.36) first.  Given \(\delta>0\), choose the
maximum of the finitely many preceding thresholds,
the compatibility threshold, the actual family start
plus one, \(2^{1/\epsilon}\), and \(2\).
For every later \(x\) and every \(i\in\mathcal F.Ix(x)\)
all twelve fields hold simultaneously.  This is
exactly
\[
 \exists C_{31}>0,\ C_e,C_m\ge0\
 \ \forall\delta>0\
 \ \exists x_0>1\
 \ \forall x\ge x_0\
 \ \forall i,\quad
             \mathrm{admitted}(D,L,X,Y^0,K;
                                      C_{31},C_e,C_m,\delta).
                                                               \tag{LGS.37}
\]
There is no index-dependent coefficient choice and
no change of family.  If \(\mathrm{Adm2F}\) is supplied
directly, use its threshold in place of GLG16; the
remaining proof is identical.
% END LGS12

% BEGIN LGS13
\paragraph{LGS13. Full \(\mathrm{LinkG}\) from the remaining common record.}
\label{claim:LGS13}
Under (LGS.2), the same family satisfies
\[
 \bigl(\mathrm{GroupCompatibilityFamily}(\mathcal F)
       \ \wedge\ \mathrm{TopsF}(\mathcal F)\bigr)
               \ \Longrightarrow\ \mathrm{LinkG}(\mathcal F).
                                                               \tag{LGS.38}
\]
One may use \(\mathrm{Adm2F}\) in place of
\(\mathrm{TopsF}\).  There are a constant \(C_G\) and
one threshold \(x_0\), before all later members and
places, such that the error is
\[
           b_o=b^{\delta_G(x)},\qquad
           \delta_G(x)=x^{-\epsilon/2}\longrightarrow0,
           \qquad C_G=4C_{\mathrm{count}}+1.          \tag{LGS.39}
\]
Here \(C_{\mathrm{count}}\ge0\) is exactly the
constant of GLG9 with (LGS.36) and the given \(K\).
No constant or loss term in the literal \(\mathrm{LinkG}\)
is changed.

Apply the complete remaining counting theorem GLG9
to the actual \(X\) and the auxiliary \(Y^0\), whose
full admission is LGS12.  This gives its
\(\mathrm{GainBoundAll}\) and ask comparison for
every pair in the original degree window, in
particular for \(a_L,b_L\).
LGS7 supplies the actual \(\mathrm{MeanKappaBound}\).
LGS8 supplies rational descent, and GLG15 then
identifies all local lengths with the actual
\(L.\mathrm{lenMN}\), including the supported place
with empty served set.  These identifications
continue to hold for \(Y^0\) by LGS11.
The family's original \(\mathrm{GatTies}\),
\(\mathrm{LatTied}\), filling identities, support,
\(\mathrm{NuDef}\), and unserved identity are the
unchanged records on the same member.
GLG10 and GLG12 supply the served-place and depth
conditions.  Finally LGS11 supplies every conjunct
of the local \(\mathrm{admG1}\).
Thus all premises of GLG17 have now been produced.

For precision, its five conclusions here are:
\begin{enumerate}
\item For every \(u,w\) with \(p\ne T_u\), the original
depth inclusion
\[
 \mathrm{LocLe}\bigl(\lfloor\mathrm{ur}_w(\tau(w))\rfloor_+,
                          L.NuLat(u),\mathrm{MramL}(u)\bigr)
 \]
holds, and
\[
 L.\mathrm{lenMN}(u,w)
       \le D.\mathrm{duS}(u)\{\tau(w)+\mathrm{io}(w)\}.
 \]
An empty served set forces all three quantities
\(\tau(w),\mathrm{io}(w),L.\mathrm{lenMN}(u,w)\) to zero.
\item For every prime and served representative,
\(S_v\le2mN_va/p\).
\item At every original signal or inert place with
\(p\le X''\), each gain lies in
\([0,4|TT|a/p]\), and the exact class-sum inequality is
\[
 \begin{aligned}
 &\sum_{j\in\mathrm{lifts}}\ \sum_{t\in TT_w}
       \{\mathrm{vonMangoldtClassSum}(j,t,p)
                         -\mathrm{loG}(L,C_G,b_o,p)\}\\
 &\qquad-m\,\mathrm{flG}(L,C_G,p)
       \le\sum_{j\in\mathrm{lifts}}\mathrm{gain}(j,w)-\tau(w).
 \end{aligned}
 \]
\item The full original \(\mathrm{ServedOrbitsContained}\)
holds.
\item The full original \(\mathrm{MeanKappaBound}\)
holds for \(X\), with both clauses of (LGS.22).
\end{enumerate}
The scalar \(vUnit(G)\) in (LGS.28) cancels by the
exact local length identity in GLG15; the actual
\(\mathrm{GatTies}\) rank identity gives
\(G.d=D.\mathrm{duS}(u)\).
This explains in particular why no coefficient-field
ramification factor is missing in item 1.
The original losses \(\mathrm{loG},\mathrm{flG}\)
and the constant \(4|C_{\mathrm{count}}|+1\) are
the exact GLG11 comparison.  Its absolute value is
unnecessary only because \(C_{\mathrm{count}}\ge0\).

Take the maximum of the finitely many family thresholds
just used.  All five conclusions concern the same
\(x,i,u,w\), lattices, selected tops and degree functions.
Their independence of \(Ch\), checked in LGS11,
turns the auxiliary application into the literal
\(\mathrm{LinkG}(\mathcal F)\).  This proves (LGS.38).

The original class-level implication also follows
with its full choice order.  If a family class
\(\mathcal C\) has
\[
 \mathrm{Closed}(\mathcal C,\mathrm{GroupCompatibilityFamily})
 \quad\hbox{and}\quad
 \mathrm{Closed}(\mathcal C,\mathrm{Adm2F}),
 \]
then it has \(\mathrm{Closed}(\mathcal C,\mathrm{LinkG})\).
First decrease the two supplied positive exponent
cutoffs and \(1/104\) to a common positive cutoff.
For each later fixed \(\epsilon\), take the pointwise
maximum of their two \(\mathrm{TransportedAlt}\)
bounds and
\[
 K_{\mathrm{local}}.\mathrm{pre}=0,\quad
 K_{\mathrm{local}}.t=0,\quad
 K_{\mathrm{local}}.\mathrm{mid}=4129,\quad
 K_{\mathrm{local}}.\mathrm{last}
             =\max\{5600002,C_{\mathrm{orb}}(\epsilon)\}.
                                                               \tag{LGS.40}
\]
The actual \(\mathrm{cMid}\) contains \(K.\mathrm{La}\);
the actual \(\mathrm{cLast}\) contains every other
coefficient constrained in (LGS.2).
Thus \(\mathrm{FamP.Large}\) at this bound implies
(LGS.2) and both supplied class predicates.
Only then is \(\mathcal F\), its height parameter,
and its eventual \(x\)-threshold considered.
This proves the literal source implication
\(\mathrm{linkG\_\allowbreak concrete5}\), with its original
common-family requirement.

The construction of \(\mathrm{GroupCompatibilityFamily}\)
remains outside this theorem.  In particular separate
adelic forms, an abstract local representation, or
unrelated lattice comparisons do not supply its five
simultaneous fields (LGS.24).  The selected-top
assumption is also retained when not provided by
another completed construction.  Consequently this
chapter supplies the local-growth, mean and admission
producers and the exact original installer, while
making no whole-leaf closure claim for the global
local-gain target.
% END LGS13

\endgroup
\end{document}
